\newMonth{Julius}
\setrunningtitles{Julius}{July}

% ---- martyrology/mart07/mart0701.htm
\needspace{10\baselineskip}
\begin{paracol}{2}
\selectlanguage{latin}
\begin{center}{\color{gregoriocolor} Kaléndis Júlii. 
 Luna\dots\ }\end{center}
\switchcolumn
\selectlanguage{english}
\begin{center}{\color{gregoriocolor} The 
 1\textsuperscript{st} Day of 
 July. The\dots\ Day of the Moon.}\end{center}
\end{paracol}

\noindent\begin{tabularx}{\linewidth}{*{19}{>{\centering\arraybackslash}X}}
 \textcolor{gregoriocolor}{a} & \textcolor{gregoriocolor}{b} & \textcolor{gregoriocolor}{c} & \textcolor{gregoriocolor}{d} & \textcolor{gregoriocolor}{e} & \textcolor{gregoriocolor}{f} & \textcolor{gregoriocolor}{g} & \textcolor{gregoriocolor}{h} & \textcolor{gregoriocolor}{i} & \textcolor{gregoriocolor}{k} & \textcolor{gregoriocolor}{l} & \textcolor{gregoriocolor}{m} & \textcolor{gregoriocolor}{n} & \textcolor{gregoriocolor}{p} & \textcolor{gregoriocolor}{q} & \textcolor{gregoriocolor}{r} & \textcolor{gregoriocolor}{s} & \textcolor{gregoriocolor}{t} & \textcolor{gregoriocolor}{u} \\
 6 & 7 & 8 & 9 & 10 & 11 & 12 & 13 & 14 & 15 & 16 & 17 & 18 & 19 & 20 & 21 & 22 & 23 & 24 \\
\end{tabularx}
\vspace{0.5\baselineskip}
\noindent\begin{tabularx}{\linewidth}{*{12}{>{\centering\arraybackslash}X}}
 \textcolor{gregoriocolor}{A} & \textcolor{gregoriocolor}{B} & \textcolor{gregoriocolor}{C} & \textcolor{gregoriocolor}{D} & \textcolor{gregoriocolor}{E} & F & \textcolor{gregoriocolor}{F} & \textcolor{gregoriocolor}{G} & \textcolor{gregoriocolor}{H} & \textcolor{gregoriocolor}{M} & \textcolor{gregoriocolor}{N} & \textcolor{gregoriocolor}{P} \\
 25 & 26 & 27 & 28 & 29 & 30 & 29 & 1 & 2 & 3 & 4 & 5 \\
\end{tabularx}

\begin{paracol}{2}
\selectlanguage{latin}
\lettrine[lines=1]{O}{ctáva} Nativitátis sancti Joánnis Baptístæ.
\switchcolumn
\selectlanguage{english}
\lettrine[lines=1]{T}{he} Octave of the Nativity of St. John the Baptist.
\switchcolumn*
\selectlanguage{latin}
Festum pretiosíssimi 
 Sánguinis Dómini nostri Jesu Christi.
\switchcolumn
\selectlanguage{english}
The feast of the most Precious Blood 
 of our Lord Jesus Christ.
\switchcolumn*
\selectlanguage{latin}
In monte Hor deposítio 
 sancti Aaron, primi ex órdine Levítico Sacerdótis.
\switchcolumn
\selectlanguage{english}
On Mt. Hor, the death of St. Aaron, 
 the first priest of the Levitical order.
\switchcolumn*
\selectlanguage{latin}
Viénnæ, in Gállia, 
 sancti Martíni Epíscopi, Apostolórum discípuli.
\switchcolumn
\selectlanguage{english}
At Vienne in France, St. Martin, a 
 bishop who was a disciple of the apostles.
\switchcolumn*
\selectlanguage{latin}
Sinuéssæ, in Campánia, 
 sanctórum Mártyrum Casti et Secundíni Episcopórum.
\switchcolumn
\selectlanguage{english}
At Sinuessa in Campánia the holy 
 martyrs Castus and Secundinus, bishops.
\switchcolumn*
\selectlanguage{latin}
In Británnia sanctórum 
 Mártyrum Júlii et Aaron, qui post sanctum Albánum, in persecutióne 
 Diocletiáni Imperatóris, passi sunt; quo témpore ibídem quamplúrimi, 
 divérsis cruciátibus torti et sævíssime laceráti, ad supérnæ civitátis 
 gáudia, consummáto agóne, pervenérunt.
\switchcolumn
\selectlanguage{english}
In England, the holy martyrs Julius 
 and Aaron, who suffered after St. Alban in the persecution of Diocletian. 
 In the same country a great number were tortured at that time in different 
 ways and barbarously lacerated, ended their combat, and attained to the joys 
 of the heavenly city.
\switchcolumn*
\selectlanguage{latin}
Arvérnis, in Gállia, 
 sancti Galli Epíscopi.
\switchcolumn
\selectlanguage{english}
In Auvergne in France, St. Gall, 
 bishop.
\switchcolumn*
\selectlanguage{latin}
In território 
 Lugdunénsi deposítio sancti Domitiáni Abbátis, qui primus illic eremitícam 
 vitam exércuit, et, cum plúrimos ibi in Dei servítio congregásset, tandem, 
 magnis virtútibus et gloriósis miráculis valde clarus, ad patres, in 
 senectúte bona, colléctus est.
\switchcolumn
\selectlanguage{english}
In the diocese of Lyons, the death 
 of St. Domitian, abbot, who was first to lead the life of a monk in that 
 district. After having called together many servants of God to that 
 place, and having gained great renown for virtues and miracles, he was 
 summoned to his fathers at an advanced age.
\switchcolumn*
\selectlanguage{latin}
Engolísmæ, in Gállia, 
 sancti Epárchii Abbátis.
\switchcolumn
\selectlanguage{english}
At Angoulême, St. Eparchius, abbot.
\switchcolumn*
\selectlanguage{latin}
In território Rheménsi 
 sancti Theodoríci Presbyteri, qui fuit discípulus beáti Remígii Epíscopi.
\switchcolumn
\selectlanguage{english}
In the diocese of Rheims, St. 
 Theodoric, priest and disciple of the blessed Bishop Remigius.
\switchcolumn*
\selectlanguage{latin}
Apud Eméssam, in 
 Phœnícia, sancti Simeónis Confessóris, cognoménto Sali, qui stultus propter 
 Christum factus est; sed altam ejus sapiéntiam Deus magnis miráculis 
 declarávit.
\switchcolumn
\selectlanguage{english}
At Emesa, St. Simeon, surnamed Salus, 
 confessor. He had feigned to be an idiot for the sake of Christ, but 
 God manifested his high wisdom by great miracles.
\switchcolumn*
\selectlanguage{latin}
\end{paracol}

% \continueMonth{Julius}

% ---- martyrology/mart07/mart0702.htm
\needspace{10\baselineskip}
\begin{paracol}{2}
\selectlanguage{latin}
\begin{center}{\color{gregoriocolor} Sexto Nonas Júlii. 
 Luna\dots\ }\end{center}
\switchcolumn
\selectlanguage{english}
\begin{center}{\color{gregoriocolor} The 2\textsuperscript{nd} Day of July. The\dots\ Day of the 
 Moon.}\end{center}
\end{paracol}

\noindent\begin{tabularx}{\linewidth}{*{19}{>{\centering\arraybackslash}X}}
 \textcolor{gregoriocolor}{a} & \textcolor{gregoriocolor}{b} & \textcolor{gregoriocolor}{c} & \textcolor{gregoriocolor}{d} & \textcolor{gregoriocolor}{e} & \textcolor{gregoriocolor}{f} & \textcolor{gregoriocolor}{g} & \textcolor{gregoriocolor}{h} & \textcolor{gregoriocolor}{i} & \textcolor{gregoriocolor}{k} & \textcolor{gregoriocolor}{l} & \textcolor{gregoriocolor}{m} & \textcolor{gregoriocolor}{n} & \textcolor{gregoriocolor}{p} & \textcolor{gregoriocolor}{q} & \textcolor{gregoriocolor}{r} & \textcolor{gregoriocolor}{s} & \textcolor{gregoriocolor}{t} & \textcolor{gregoriocolor}{u} \\
 7 & 8 & 9 & 10 & 11 & 12 & 13 & 14 & 15 & 16 & 17 & 18 & 19 & 20 & 21 & 22 & 23 & 24 & 25 \\
\end{tabularx}
\vspace{0.5\baselineskip}
\noindent\begin{tabularx}{\linewidth}{*{12}{>{\centering\arraybackslash}X}}
 \textcolor{gregoriocolor}{A} & \textcolor{gregoriocolor}{B} & \textcolor{gregoriocolor}{C} & \textcolor{gregoriocolor}{D} & \textcolor{gregoriocolor}{E} & F & \textcolor{gregoriocolor}{F} & \textcolor{gregoriocolor}{G} & \textcolor{gregoriocolor}{H} & \textcolor{gregoriocolor}{M} & \textcolor{gregoriocolor}{N} & \textcolor{gregoriocolor}{P} \\
 26 & 27 & 28 & 29 & 30 & 1 & 1 & 2 & 3 & 4 & 5 & 6 \\
\end{tabularx}

\begin{paracol}{2}
\selectlanguage{latin}
\lettrine[lines=2]{V}{isitátio} beátæ Maríæ 
 Vírginis ad Elísabeth.
\switchcolumn
\selectlanguage{english}
\lettrine[lines=2]{T}{he} Visitation of the Blessed Virgin 
 Mary to Elizabeth.
\switchcolumn*
\selectlanguage{latin}
Romæ, via Aurélia, natális sanctórum Mártyrum Procéssi et Martiniáni, qui, a beáto Petro 
 Apóstolo in Mamertíni custódia baptizáti, et, sub Neróne, oris contusiónem, 
 equúleum, nervos, fustes, flammas scorpionésque perpéssi, novíssime, gládio 
 percússi, martyrio coronáti sunt.
\switchcolumn
\selectlanguage{english}
At Rome, on the Aurelian Way, the 
 birthday of the holy martyrs Processus and Martinian, who were baptized by 
 the blessed apostle Peter in the Mamertine Prison. After being struck 
 on the mouth, racked, scourged with knotted ropes and whips strung with 
 pieces of metal; after being beaten with rods and exposed to the flames, 
 they were beheaded in the days of Nero, thus obtaining the crown of 
 martyrdom.
\switchcolumn*
\selectlanguage{latin}
Item Romæ pássio 
 sanctórum trium mílitum, qui, ad Christum convérsi in martyrio beáti Pauli 
 Apóstoli, cum eo partícipes fíeri cæléstis glóriæ meruérunt.
\switchcolumn
\selectlanguage{english}
Also at Rome, three holy soldiers, 
 who were converted to Christ by the martyrdom of the blessed apostle Paul, 
 and with him merited to be made partakers of the heavenly glory.
\switchcolumn*
\selectlanguage{latin}
Eódem die sanctórum 
 Mártyrum Aristónis, Crescentiáni, Eutychiáni, Urbáni, Vitális, Justi, 
 Felicíssimi, Felícis, Márciæ et Symphorósæ; qui omnes apud Campániam, cum 
 Diocletiáni Imperatóris persecútio desævíret, martyrio coronáti sunt.
\switchcolumn
\selectlanguage{english}
The same day, the holy martyrs 
 Ariston, Crescentian, Eutychian, Urbanus, Vitalis, Justus, Felicissimus, 
 Felix, Marcia, and Symphorosa, all of whom were crowned with martyrdom when 
 the persecution of Emperor Diocletian was raging.
\switchcolumn*
\selectlanguage{latin}
Wintóniæ, in Anglia, 
 sancti Swithúni Epíscopi, cujus sánctitas miráculis effúlsit.
\switchcolumn
\selectlanguage{english}
At Winchester in England, St. 
 Swithin, bishop, whose sanctity was verified by the gift of miracles.
\switchcolumn*
\selectlanguage{latin}
Bambérgæ sancti Othónis 
 Epíscopi, qui Pomeránis Evangélium prædicávit, et eos ad fidem convértit.
\switchcolumn
\selectlanguage{english}
At Bamberg, the holy bishop Otho, 
 who preached the Gospel to the people of Pomerania, and converted them to 
 the faith.
\switchcolumn*
\selectlanguage{latin}
Lyciis, in Apúlia, sancti Bernardíni Realíno, Confessóris, qui magistrátus múnere egrégie 
 functus, Societátem Jesu ingréssus et sacerdótio auctus, caritáte ac 
 miráculis incláruit.
\switchcolumn
\selectlanguage{english}
At Lecce in Apulia, St. Bernardino 
 Realino, confessor, who after practising the legal profession as a judge, 
 entered the Society of Jesus, was ordained to the priesthood, and was 
 renowned for his charity and miracles.
\switchcolumn*
\selectlanguage{latin}
Turónis, in Gállia, 
 deposítio sanctæ Monegúndis, religiósæ féminæ.
\switchcolumn
\selectlanguage{english}
At Tours, the death of St. 
 Monegundes, a pious woman.
\switchcolumn*
\selectlanguage{latin}
\end{paracol}


% ---- martyrology/mart07/mart0703.htm
\needspace{10\baselineskip}
\begin{paracol}{2}
\selectlanguage{latin}
\begin{center}{\color{gregoriocolor} Quinto Nonas Júlii. 
 Luna\dots\ }\end{center}
\switchcolumn
\selectlanguage{english}
\begin{center}{\color{gregoriocolor} The 
 3\textsuperscript{rd} Day of 
 July. The\dots\ Day of the Moon.}\end{center}
\end{paracol}

\noindent\begin{tabularx}{\linewidth}{*{19}{>{\centering\arraybackslash}X}}
 \textcolor{gregoriocolor}{a} & \textcolor{gregoriocolor}{b} & \textcolor{gregoriocolor}{c} & \textcolor{gregoriocolor}{d} & \textcolor{gregoriocolor}{e} & \textcolor{gregoriocolor}{f} & \textcolor{gregoriocolor}{g} & \textcolor{gregoriocolor}{h} & \textcolor{gregoriocolor}{i} & \textcolor{gregoriocolor}{k} & \textcolor{gregoriocolor}{l} & \textcolor{gregoriocolor}{m} & \textcolor{gregoriocolor}{n} & \textcolor{gregoriocolor}{p} & \textcolor{gregoriocolor}{q} & \textcolor{gregoriocolor}{r} & \textcolor{gregoriocolor}{s} & \textcolor{gregoriocolor}{t} & \textcolor{gregoriocolor}{u} \\
 8 & 9 & 10 & 11 & 12 & 13 & 14 & 15 & 16 & 17 & 18 & 19 & 20 & 21 & 22 & 23 & 24 & 25 & 26 \\
\end{tabularx}
\vspace{0.5\baselineskip}
\noindent\begin{tabularx}{\linewidth}{*{12}{>{\centering\arraybackslash}X}}
 \textcolor{gregoriocolor}{A} & \textcolor{gregoriocolor}{B} & \textcolor{gregoriocolor}{C} & \textcolor{gregoriocolor}{D} & \textcolor{gregoriocolor}{E} & F & \textcolor{gregoriocolor}{F} & \textcolor{gregoriocolor}{G} & \textcolor{gregoriocolor}{H} & \textcolor{gregoriocolor}{M} & \textcolor{gregoriocolor}{N} & \textcolor{gregoriocolor}{P} \\
 27 & 28 & 29 & 30 & 1 & 2 & 2 & 3 & 4 & 5 & 6 & 7 \\
\end{tabularx}

\begin{paracol}{2}
\selectlanguage{latin}
\lettrine[lines=2]{R}{omæ} natális sancti 
 Leónis Secúndi, Papæ et Confessóris, qui primo sui Pontificátus anno, plenus 
 méritis, migrávit in cælum.
\switchcolumn
\selectlanguage{english}
\lettrine[lines=2]{A}{t} Rome, the birthday of Pope St. 
 Leo II, confessor, who passed to heaven filled with merits during the first 
 year of his pontificate.
\switchcolumn*
\selectlanguage{latin}
Clúsii, in Etrúria, 
 sanctórum Mártyrum Irenæi Diáconi, et Mustíolæ matrónæ; qui sub Aureliáno 
 Imperatóre, divérsis atrocibúsque supplíciis cruciáti, corónam martyrii 
 meruérunt.
\switchcolumn
\selectlanguage{english}
At Chiusi in Tuscany, in the reign 
 of Emperor Aurelian, the holy martyrs Irenaeus, a deacon, and Mustiola, a 
 matron, who were subjected to various atrocious tortures and merited the 
 crown of martyrdom.
\switchcolumn*
\selectlanguage{latin}
Alexandríæ sanctórum 
 Mártyrum Tryphónis et aliórum duódecim.
\switchcolumn
\selectlanguage{english}
At Alexandria, the holy martyrs 
 Trypho and twelve others.
\switchcolumn*
\selectlanguage{latin}
Constantinópoli 
 sanctórum Eulógii et Sociórum Mártyrum.
\switchcolumn
\selectlanguage{english}
At Constantinople, the holy martyrs 
 Eulogius and his companions.
\switchcolumn*
\selectlanguage{latin}
Cæsaréæ, in Cappadócia, 
 sancti Hyacínthi, qui Trajáni Imperatóris cubiculárius fuit; et, accusátus 
 quod Christiánus esset, plagis várie afflíctus est, atque, in cárcerem 
 trusus, ibi fame consúmptus exspirávit.
\switchcolumn
\selectlanguage{english}
At Caesarea in Cappadocia, St. 
 Hyacinth, chamberlain of the emperor Trajan. Accused of being a 
 Christian, he was scourged and thrown into prison where he died of hunger.
\switchcolumn*
\selectlanguage{latin}
Eódem die sanctórum 
 Mártyrum Marci et Muciáni, qui pro Christo sunt gládio cæsi. Hos cum 
 puer párvulus, ne immolárent idólis, alta voce monéret, verbéribus jussus 
 est cædi; cumque veheméntius Christum confiterétur, tandem, una cum quodam 
 Paulo, Mártyres exhortánte, necátus est.
\switchcolumn
\selectlanguage{english}
The same day, the holy martyrs Mark 
 and Mucian, who were put to the sword for Christ. A small boy who 
 cried out to them not to sacrifice to idols was then whipped, but confessing 
 Christ still more vehemently, he was put to death with a man named Paul, who 
 had also exhorted the martyrs.
\switchcolumn*
\selectlanguage{latin}
Laodicéæ, in Syria, 
 sancti Anatólii Epíscopi, qui non solum religiósis viris, sed et philósophis 
 admiránda scripta relíquit.
\switchcolumn
\selectlanguage{english}
At Laodicea in Syria, St. Anatolius, 
 a bishop whose writings were admired not only by religious men, but by 
 philosophers.
\switchcolumn*
\selectlanguage{latin}
Altíni, in Venetórum 
 fínibus, sancti Heliodóri Epíscopi, doctrína et sanctitáte insígnis.
\switchcolumn
\selectlanguage{english}
At Altino, St. Heliodorus, a bishop 
 distinguished for holiness and learning.
\switchcolumn*
\selectlanguage{latin}
Ravénnæ sancti Dathi, 
 Epíscopi et Confessóris.
\switchcolumn
\selectlanguage{english}
At Ravenna, St. Dathus, bishop and 
 confessor.
\switchcolumn*
\selectlanguage{latin}
Edéssæ, in Mesopotámia, Translátio sancti Thomæ Apóstoli ex India; cujus relíquiæ Ortónam, apud 
 Frentános, póstea translátæ sunt.
\switchcolumn
\selectlanguage{english}
At Edessa in Mesopotamia, the 
 translation of the apostle St. Thomas from India. His relics were 
 afterwards taken to Ortona.
\switchcolumn*
\selectlanguage{latin}
\end{paracol}


% ---- martyrology/mart07/mart0704.htm
\needspace{10\baselineskip}
\begin{paracol}{2}
\selectlanguage{latin}
\begin{center}{\color{gregoriocolor} Quarto Nonas Júlii. 
 Luna\dots\ }\end{center}
\switchcolumn
\selectlanguage{english}
\begin{center}{\color{gregoriocolor} The 
 4\textsuperscript{th} Day of 
 July. The\dots\ Day of the Moon.}\end{center}
\end{paracol}

\noindent\begin{tabularx}{\linewidth}{*{19}{>{\centering\arraybackslash}X}}
 \textcolor{gregoriocolor}{a} & \textcolor{gregoriocolor}{b} & \textcolor{gregoriocolor}{c} & \textcolor{gregoriocolor}{d} & \textcolor{gregoriocolor}{e} & \textcolor{gregoriocolor}{f} & \textcolor{gregoriocolor}{g} & \textcolor{gregoriocolor}{h} & \textcolor{gregoriocolor}{i} & \textcolor{gregoriocolor}{k} & \textcolor{gregoriocolor}{l} & \textcolor{gregoriocolor}{m} & \textcolor{gregoriocolor}{n} & \textcolor{gregoriocolor}{p} & \textcolor{gregoriocolor}{q} & \textcolor{gregoriocolor}{r} & \textcolor{gregoriocolor}{s} & \textcolor{gregoriocolor}{t} & \textcolor{gregoriocolor}{u} \\
 9 & 10 & 11 & 12 & 13 & 14 & 15 & 16 & 17 & 18 & 19 & 20 & 21 & 22 & 23 & 24 & 25 & 26 & 27 \\
\end{tabularx}
\vspace{0.5\baselineskip}
\noindent\begin{tabularx}{\linewidth}{*{12}{>{\centering\arraybackslash}X}}
 \textcolor{gregoriocolor}{A} & \textcolor{gregoriocolor}{B} & \textcolor{gregoriocolor}{C} & \textcolor{gregoriocolor}{D} & \textcolor{gregoriocolor}{E} & F & \textcolor{gregoriocolor}{F} & \textcolor{gregoriocolor}{G} & \textcolor{gregoriocolor}{H} & \textcolor{gregoriocolor}{M} & \textcolor{gregoriocolor}{N} & \textcolor{gregoriocolor}{P} \\
 28 & 29 & 30 & 1 & 2 & 3 & 3 & 4 & 5 & 6 & 7 & 8 \\
\end{tabularx}

\begin{paracol}{2}
\selectlanguage{latin}
\lettrine[lines=2]{S}{tremótii,} in Lusitánia, 
 natális sanctæ Elísabeth Víduæ, Lusitanórum Regínæ, quam, virtútibus et 
 miráculis claram, Urbánus Octávus, Póntifex Máximus, in Sanctórum númerum 
 rétulit. Ejus tamen celébritas octávo Idus mensis hujus recólitur, ex 
 dispositióne Innocéntii Papæ Duodécimi.
\switchcolumn
\selectlanguage{english}
\lettrine[lines=2]{A}{t} Estremos in Portugal, the 
 birthday of St. Elizabeth the Widow, queen of Portugal, whom Pope Urban 
 VIII, mindful of her virtues and miracles, placed among the number of the 
 saints. Pope Innocent XII ordered her feast to be kept on the 8th of 
 July.
\switchcolumn*
\selectlanguage{latin}
Sanctórum Osée et Aggǽi 
 Prophetárum.
\switchcolumn
\selectlanguage{english}
The holy prophets Osee and Aggaeus.
\switchcolumn*
\selectlanguage{latin}
In território 
 Bituricénsi sancti Lauriáni, Epíscopi Hispalénsis et Mártyris, cujus caput 
 Híspalim, in Hispánia, delátum est.
\switchcolumn
\selectlanguage{english}
In the diocese of Bourges, St. 
 Laurian, bishop of Seville and martyr, whose head was taken to Seville in 
 Spain.
\switchcolumn*
\selectlanguage{latin}
In Africa natális 
 sancti Jucundiáni Mártyris, pro Christo in mare demérsi.
\switchcolumn
\selectlanguage{english}
In Africa, the birthday of St. 
 Jucundian, a martyr who was drowned in the sea for Christ.
\switchcolumn*
\selectlanguage{latin}
Sírmii sanctórum 
 Mártyrum Innocéntii et Sebástiæ, cum áliis trigínta.
\switchcolumn
\selectlanguage{english}
At Sirmium, Saints Innocent and 
 Sebastia, with thirty other martyrs.
\switchcolumn*
\selectlanguage{latin}
Madáuri, in Africa, 
 sancti Namphaniónis Mártyris et Sociórum, quos ille roborávit ad pugnam, et 
 ad corónam provéxit.
\switchcolumn
\selectlanguage{english}
At Madaurus in Africa, the martyr 
 Namphanion and his companions, whom he strengthened for the combat and led 
 to the crown of martyrdom.
\switchcolumn*
\selectlanguage{latin}
Cyréne, in Líbya, sancti Theodóri Epíscopi, qui, in persecutióne Diocletiáni, sub Digniáno 
 Prǽside, plumbátis cæsus et lingua abscíssus est; atque in pace tandem 
 Conféssor occúbuit.
\switchcolumn
\selectlanguage{english}
At Cyrene in Libya, the holy bishop 
 Theodore. In the persecution of Diocletian, under the governor Dignian, 
 he was scourged with leaded whips and had his tongue cut out. Finally, 
 however, he died a confessor.
\switchcolumn*
\selectlanguage{latin}
Augústæ, in Rhǽtia, 
 sancti Uldaríci Epíscopi, miræ abstinéntiæ, largitátis et vigilántiæ virtúte, 
 ac miraculórum dono illústris.
\switchcolumn
\selectlanguage{english}
At Augsburg in Germany, St. Uldaric, 
 a bishop illustrious for extraordinary abstinence, liberality, vigilance, 
 and the gift of miracles.
\switchcolumn*
\selectlanguage{latin}
Turónis, in Gállia, 
 Translátio sancti Martíni, Epíscopi et Confessóris; et Dedicátio Basílicæ 
 suæ hoc ipso die, quo étiam is ante áliquot annos in Epíscopum fúerat 
 ordinátus.
\switchcolumn
\selectlanguage{english}
At Tours in France, the translation 
 of St. Martin, bishop and confessor, and the dedication of his basilica, 
 consecrated on the same day that he had been raised to the episcopate some 
 years previously.
\switchcolumn*
\selectlanguage{latin}
\end{paracol}


% ---- martyrology/mart07/mart0705.htm
\needspace{10\baselineskip}
\begin{paracol}{2}
\selectlanguage{latin}
\begin{center}{\color{gregoriocolor} Tértio Nonas Júlii. 
 Luna\dots\ }\end{center}
\switchcolumn
\selectlanguage{english}
\begin{center}{\color{gregoriocolor} The 
 5\textsuperscript{th} Day of 
 July. The\dots\ Day of the Moon.}\end{center}
\end{paracol}

\noindent\begin{tabularx}{\linewidth}{*{19}{>{\centering\arraybackslash}X}}
 \textcolor{gregoriocolor}{a} & \textcolor{gregoriocolor}{b} & \textcolor{gregoriocolor}{c} & \textcolor{gregoriocolor}{d} & \textcolor{gregoriocolor}{e} & \textcolor{gregoriocolor}{f} & \textcolor{gregoriocolor}{g} & \textcolor{gregoriocolor}{h} & \textcolor{gregoriocolor}{i} & \textcolor{gregoriocolor}{k} & \textcolor{gregoriocolor}{l} & \textcolor{gregoriocolor}{m} & \textcolor{gregoriocolor}{n} & \textcolor{gregoriocolor}{p} & \textcolor{gregoriocolor}{q} & \textcolor{gregoriocolor}{r} & \textcolor{gregoriocolor}{s} & \textcolor{gregoriocolor}{t} & \textcolor{gregoriocolor}{u} \\
 10 & 11 & 12 & 13 & 14 & 15 & 16 & 17 & 18 & 19 & 20 & 21 & 22 & 23 & 24 & 25 & 26 & 27 & 28 \\
\end{tabularx}
\vspace{0.5\baselineskip}
\noindent\begin{tabularx}{\linewidth}{*{12}{>{\centering\arraybackslash}X}}
 \textcolor{gregoriocolor}{A} & \textcolor{gregoriocolor}{B} & \textcolor{gregoriocolor}{C} & \textcolor{gregoriocolor}{D} & \textcolor{gregoriocolor}{E} & F & \textcolor{gregoriocolor}{F} & \textcolor{gregoriocolor}{G} & \textcolor{gregoriocolor}{H} & \textcolor{gregoriocolor}{M} & \textcolor{gregoriocolor}{N} & \textcolor{gregoriocolor}{P} \\
 29 & 30 & 1 & 2 & 3 & 4 & 4 & 5 & 6 & 7 & 8 & 9 \\
\end{tabularx}

\begin{paracol}{2}
\selectlanguage{latin}
\lettrine[lines=2]{C}{remónæ,} in Insúbria, 
 sancti Antónii-Maríæ Zaccaría, Confessóris, qui Clericórum Regulárium sancti 
 Pauli et Angelicárum Vírginum fuit Institútor; atque, virtútibus ómnibus et 
 miráculis insígnis, a Leóne Papa Décimo tértio inter Sanctos adscríptus est. 
 Ejus corpus Medioláni, in Ecclésia sancti Bárnabæ cólitur.
\switchcolumn
\selectlanguage{english}
\lettrine[lines=2]{A}{t} Cremona in Italy, St. 
 Anthony-Mary Zacharias, confessor, founder of the Clerks Regular of St. Paul 
 and also of the Angelic Virgins. Distinguished for all the virtues and 
 for miracles, he was placed among the saints by Pope Leo XIII. His 
 body is venerated in the Church of St. Barnabas at Milan.
\switchcolumn*
\selectlanguage{latin}
Romæ sanctæ Zoæ 
 Mártyris, uxóris beáti Nicostráti Mártyris, quæ, sub Diocletiáno Imperatóre, 
 dum ad confessiónem beáti Petri Apóstoli oráret, a persecutóribus est 
 arctáta et in custódiam obscuríssimam trusa; deínde collo et capíllis in 
 árbore, adhíbito subter horríbili fumo, suspénsa est, et ita, in confessióne 
 Dómini, emísit spíritum.
\switchcolumn
\selectlanguage{english}
At Rome, St. Zoe, martyr, wife of 
 the blessed martyr Nicostratus. While praying at the tomb of the 
 apostle St. Peter, during the time of Diocletian, she was seized by the 
 persecutors, cast into a dark dungeon, then hanged on a tree by her neck and 
 hair, and suffocated by a loathsome smoke, finally yielding up her soul in 
 the confession of the Lord.
\switchcolumn*
\selectlanguage{latin}
Hierosólymis sancti 
 Athanásii Diáconi, qui, pro sancta Synodo Chalcedonénsi, ab hæréticis 
 comprehénsus et ómnium cruciátum expértus génera, ferro tandem necátus est.
\switchcolumn
\selectlanguage{english}
At Jerusalem, St. Athanasius, a 
 deacon, who was apprehended by the heretics for defending the Council of 
 Chalcedon, and after experiencing all kinds of torments, was finally put to 
 the sword.
\switchcolumn*
\selectlanguage{latin}
In Syria natális 
 sancti Domítii Mártyris, qui virtútibus suis multa íncolis præstat benefícia.
\switchcolumn
\selectlanguage{english}
In Syria, the birthday of St. 
 Domitius, martyr, who confers many favours on the people of that country by 
 his miracles.
\switchcolumn*
\selectlanguage{latin}
In Sicília sanctórum 
 Mártyrum Agathónis et Triphínæ.
\switchcolumn
\selectlanguage{english}
In Sicily, the holy martyrs Agatho 
 and Triphina.
\switchcolumn*
\selectlanguage{latin}
Tomis, in Scythia, 
 sanctórum Mártyrum Maríni, Theódoti et Sédophæ.
\switchcolumn
\selectlanguage{english}
At Tomis in Scythia, the holy 
 martyrs Marinus, Theodotus, and Sedopha.
\switchcolumn*
\selectlanguage{latin}
Cyréne, in Libya, 
 sanctæ Cyríllæ Mártyris, quæ, in persecutióne Diocletiáni, ardéntes carbónes 
 cum thure super manum pósitos diu ténuit, ne, prunas excutiéndo, thus 
 obtulísse viderétur; deínde, sævíssime laniáta, próprio decoráta sánguine 
 migrávit ad Sponsum.
\switchcolumn
\selectlanguage{english}
At Cyrene in Libya, St. Cyrilla, 
 martyr, in the persecution of Diocletian. She held burning coals with 
 incense on her hand for a long time, lest by shaking off the coals she 
 should seem to offer incense to the idols. She was afterwards cruelly 
 scourged, and went to her Spouse adorned with her own blood.
\switchcolumn*
\selectlanguage{latin}
Tréviris sancti 
 Numeriáni, Epíscopi et Confessóris.
\switchcolumn
\selectlanguage{english}
At Treves, St. Numerian, bishop and 
 confessor.
\switchcolumn*
\selectlanguage{latin}
Apud Septempedános, in 
 Picéno, sanctæ Philoménæ Vírginis.
\switchcolumn
\selectlanguage{english}
At San Severino in Piceno, St. 
 Philomena, virgin.
\switchcolumn*
\selectlanguage{latin}
\end{paracol}


% ---- martyrology/mart07/mart0706.htm
\needspace{10\baselineskip}
\begin{paracol}{2}
\selectlanguage{latin}
\begin{center}{\color{gregoriocolor} Prídie Nonas Júlii. 
 Luna\dots\ }\end{center}
\switchcolumn
\selectlanguage{english}
\begin{center}{\color{gregoriocolor} The 
 6\textsuperscript{th} Day of 
 July. The\dots\ Day of the Moon.}\end{center}
\end{paracol}

\noindent\begin{tabularx}{\linewidth}{*{19}{>{\centering\arraybackslash}X}}
 \textcolor{gregoriocolor}{a} & \textcolor{gregoriocolor}{b} & \textcolor{gregoriocolor}{c} & \textcolor{gregoriocolor}{d} & \textcolor{gregoriocolor}{e} & \textcolor{gregoriocolor}{f} & \textcolor{gregoriocolor}{g} & \textcolor{gregoriocolor}{h} & \textcolor{gregoriocolor}{i} & \textcolor{gregoriocolor}{k} & \textcolor{gregoriocolor}{l} & \textcolor{gregoriocolor}{m} & \textcolor{gregoriocolor}{n} & \textcolor{gregoriocolor}{p} & \textcolor{gregoriocolor}{q} & \textcolor{gregoriocolor}{r} & \textcolor{gregoriocolor}{s} & \textcolor{gregoriocolor}{t} & \textcolor{gregoriocolor}{u} \\
 11 & 12 & 13 & 14 & 15 & 16 & 17 & 18 & 19 & 20 & 21 & 22 & 23 & 24 & 25 & 26 & 27 & 28 & 29 \\
\end{tabularx}
\vspace{0.5\baselineskip}
\noindent\begin{tabularx}{\linewidth}{*{12}{>{\centering\arraybackslash}X}}
 \textcolor{gregoriocolor}{A} & \textcolor{gregoriocolor}{B} & \textcolor{gregoriocolor}{C} & \textcolor{gregoriocolor}{D} & \textcolor{gregoriocolor}{E} & F & \textcolor{gregoriocolor}{F} & \textcolor{gregoriocolor}{G} & \textcolor{gregoriocolor}{H} & \textcolor{gregoriocolor}{M} & \textcolor{gregoriocolor}{N} & \textcolor{gregoriocolor}{P} \\
 30 & 1 & 2 & 3 & 4 & 5 & 5 & 6 & 7 & 8 & 9 & 10 \\
\end{tabularx}

\begin{paracol}{2}
\selectlanguage{latin}
\lettrine[lines=1]{O}{ctava} sanctórum Apostolórum Petri et Pauli.
\switchcolumn
\selectlanguage{english}
\lettrine[lines=1]{T}{he} Octave of the holy apostles Peter and Paul.
\switchcolumn*
\selectlanguage{latin}
Hierosólymis sancti 
 Isaiæ Prophétæ, qui, sub Manásse Rege, in duas sectus partes occúbuit, 
 sepultúsque est sub quercu Rogel, juxta tránsitum aquárum.
\switchcolumn
\selectlanguage{english}
In Jerusalem, the holy prophet 
 Isaias. During the reign of King Manasses he was put to death by being 
 sawn in two and was buried beneath the oak of Rogel, near a running stream.
\switchcolumn*
\selectlanguage{latin}
Fæsulis, in Túscia, 
 sancti Romuli, Epíscopi et Mártyris, qui fuit beáti Petri Apóstoli 
 discípulus. Hic, ab eódem Apóstolo missus ad prædicándum Evangélium, 
 in multis Itáliæ locis Christum annuntiávit; ac tandem, Fæsulas regréssus, 
 ibi, sub Domitiáno Príncipe, martyrio coronátus est cum áliis Sóciis.
\switchcolumn
\selectlanguage{english}
At Fiesole in Tuscany, St. Romulus, 
 bishop and martyr, disciple of the blessed apostle Peter, who commissioned 
 him to preach the Gospel. After announcing Christ in many parts of 
 Italy, he returned to Fiesole, and was crowned with martyrdom along with 
 other Christians in the reign of Domitian.
\switchcolumn*
\selectlanguage{latin}
Romæ natális sancti 
 Tranquillíni Mártyris, patris sanctórum Marci et Marcelliáni, qui, ad 
 prædicatiónem sancti Sebastiáni Mártyris convérsus ad Christum, a beáto 
 Polycárpo Presbytero baptizátus est, et a sancto Cajo Papa Présbyter ordinátus. Ipse, die Octavárum Apostolórum, cum ad confessiónem beáti 
 Pauli oráret, ibídem, sub Diocletiáno Imperatóre, a Pagánis tentus est, et, 
 ab eis lapidátus, martyrium consummávit.
\switchcolumn
\selectlanguage{english}
At Rome, the birthday of St. 
 Tranquillinus, martyr, father of Saints Mark and Marcellianus, who had been 
 converted to Christ by the preaching of the martyr St. Sebastian. 
 Baptized by the blessed priest Polycarp, he was ordained priest by Pope St. 
 Caius. As he prayed at the tomb of blessed Paul on the octave of the 
 apostles, he was arrested and stoned to death by the pagans, and thus 
 completed his martyrdom.
\switchcolumn*
\selectlanguage{latin}
Londíni, in Anglia, 
 sancti Thomæ More, regni Cancellárii, qui, pro fide cathólica ac beáti Petri 
 primátu, jubénte Henríco Octávo Rege, decollátus est.
\switchcolumn
\selectlanguage{english}
At London in England, on Tower Hill, 
 St. Thomas More, chancellor of the entire realm, who was beheaded by order 
 of King Henry VIII for the defence of the Catholic faith and the primacy of 
 blessed Peter.
\switchcolumn*
\selectlanguage{latin}
In Campánia sanctæ 
 Domínicæ, Vírginis et Mártyris, quæ, sub Diocletiáno Imperatóre, cum 
 fregísset idóla, hinc ad béstias damnáta, sed ab illis nil læsa, demum, 
 cápite obtruncáta, migrávit ad Dóminum. Ipsíus vero corpus Tropéæ, in 
 Calábria, summa veneratióne asservátur.
\switchcolumn
\selectlanguage{english}
In Campania, St. Dominica, virgin 
 and martyr, in the time of Emperor Diocletian. For having destroyed 
 idols, she was condemned to the beasts, but being left uninjured by them, 
 she was beheaded and departed for heaven. Her body is kept with great 
 veneration at Tropea in Calabria.
\switchcolumn*
\selectlanguage{latin}
Eódem die sanctæ Lúciæ 
 Mártyris, quæ, natióne Campána, a Ríxio Varo Vicário tenta et ácriter 
 cruciáta, eúndem convértit ad Christum; quibus adjúncti sunt Antonínus, 
 Severínus, Diodórus, Dion, et álii decem et septem, qui in passióne collégæ 
 et in coróna fuére consórtes.
\switchcolumn
\selectlanguage{english}
The same day, St. Lucia, martyr, a 
 native of Campania. Being arrested and severely tortured by the 
 lieutenant-governor Rictiovarus, she converted him to Christ. To them 
 were added Antoninus, Severinus, Diodorus, Dion, and seventeen others who 
 shared their sufferings and their crowns.
\switchcolumn*
\selectlanguage{latin}
Neptúni, in Látio, 
 sanctæ Maríæ Gorétti, piíssimæ adolescéntis, in defendénda virginitáte 
 crudelíssime necátæ, quam Pius Papa Duodécimus sanctárum Mártyrum catálogo 
 solémniter accénsuit.
\switchcolumn
\selectlanguage{english}
At Nettuno in Lazio, St. Maria 
 Goretti, a most devout young girl, who was savagely murdered for the defence 
 of her virginity, and whom Pope Pius XII solemnly added to the 
 catalogue of holy martyrs.
\switchcolumn*
\selectlanguage{latin}
In pago Trevirénsi 
 sancti Góaris, Presbyteri et Confessóris.
\switchcolumn
\selectlanguage{english}
In the vicinity of Treves, St. Goar, 
 priest and confessor.
\switchcolumn*
\selectlanguage{latin}
\end{paracol}


% ---- martyrology/mart07/mart0707.htm
\needspace{10\baselineskip}
\begin{paracol}{2}
\selectlanguage{latin}
\begin{center}{\color{gregoriocolor} Nonis Júlii. 
 Luna\dots\ }\end{center}
\switchcolumn
\selectlanguage{english}
\begin{center}{\color{gregoriocolor} The 
 7\textsuperscript{th} Day of 
 July. The\dots\ Day of the Moon.}\end{center}
\end{paracol}

\noindent\begin{tabularx}{\linewidth}{*{19}{>{\centering\arraybackslash}X}}
 \textcolor{gregoriocolor}{a} & \textcolor{gregoriocolor}{b} & \textcolor{gregoriocolor}{c} & \textcolor{gregoriocolor}{d} & \textcolor{gregoriocolor}{e} & \textcolor{gregoriocolor}{f} & \textcolor{gregoriocolor}{g} & \textcolor{gregoriocolor}{h} & \textcolor{gregoriocolor}{i} & \textcolor{gregoriocolor}{k} & \textcolor{gregoriocolor}{l} & \textcolor{gregoriocolor}{m} & \textcolor{gregoriocolor}{n} & \textcolor{gregoriocolor}{p} & \textcolor{gregoriocolor}{q} & \textcolor{gregoriocolor}{r} & \textcolor{gregoriocolor}{s} & \textcolor{gregoriocolor}{t} & \textcolor{gregoriocolor}{u} \\
 12 & 13 & 14 & 15 & 16 & 17 & 18 & 19 & 20 & 21 & 22 & 23 & 24 & 25 & 26 & 27 & 28 & 29 & 30 \\
\end{tabularx}
\vspace{0.5\baselineskip}
\noindent\begin{tabularx}{\linewidth}{*{12}{>{\centering\arraybackslash}X}}
 \textcolor{gregoriocolor}{A} & \textcolor{gregoriocolor}{B} & \textcolor{gregoriocolor}{C} & \textcolor{gregoriocolor}{D} & \textcolor{gregoriocolor}{E} & F & \textcolor{gregoriocolor}{F} & \textcolor{gregoriocolor}{G} & \textcolor{gregoriocolor}{H} & \textcolor{gregoriocolor}{M} & \textcolor{gregoriocolor}{N} & \textcolor{gregoriocolor}{P} \\
 1 & 2 & 3 & 4 & 5 & 6 & 6 & 7 & 8 & 9 & 10 & 11 \\
\end{tabularx}

\begin{paracol}{2}
\selectlanguage{latin}
\lettrine[lines=2]{S}{anctórum} Episcopórum 
 et Confessórum Cyrílli et Methódii fratrum, quorum natális respectíve ágitur 
 sextodécimo Kaléndas Mártii et octávo Idus Aprílis.
\switchcolumn
\selectlanguage{english}
\lettrine[lines=2]{T}{he} holy bishops Cyril and Methodius, 
 whose respective birthdays are on the 14th of February and the 6th of 
 April.
\switchcolumn*
\selectlanguage{latin}
Romæ sanctórum Mártyrum 
 Cláudii Commentariénsis, Nicóstrati primiscrínii, qui fuit vir beátæ 
 Mártyris Zoæ, Castórii, Victoríni et Symphoriáni; quos omnes sanctus 
 Sebastiánus ad fidem Christi perdúxit, et beátus Polycárpus Presbyter 
 baptizávit. Eósdem, in perquiréndis sanctórum Mártyrum corpóribus 
 occupátos, Fabiánus Judex comprehéndi jussit, et, cum eos, per decem dies 
 minis et blandítiis exágitans, in nullo pénitus posset commovére, jussit 
 tértio torquéri, ac póstea in mare præcípites dari.
\switchcolumn
\selectlanguage{english}
At Rome, the holy martyrs Claudius, 
 a notary, Nicostratus, an assistant prefect who had been the husband of the 
 blessed Martyr Zoa, Castorius, Victorinus, and 
 Symphorian. All these had been brought to the faith of Christ by St. Sebastian, 
 and baptized by the blessed priest Polycarp. While they were engaged 
 in searching for the bodies of the holy martyrs, the judge Fabian had them 
 arrested, and for ten days he tried to shake their constancy by threats and 
 flatteries, but being utterly unable to succeed, he ordered them to be 
 thrice tortured, then thrown into the sea.
\switchcolumn*
\selectlanguage{latin}
Dyrráchii, in 
 Macedónia, sanctórum Mártyrum Peregríni, Luciáni, Pompéji, Hesychii, Pápii, 
 Saturníni et Germáni; qui, natióne Itali, in persecutióne Trajáni, cum ad 
 eam urbem confugíssent, et ibi sanctum Astium Epíscopum, pro fide Christi, 
 in cruce pendéntem vidérent, ac se páriter Christiános esse palam 
 confiteréntur, ideo, Præsidis jussu, tenti sunt atque in mare demérsi.
\switchcolumn
\selectlanguage{english}
At Durazzo in Macedonia, the holy 
 martyrs Peregrínus, Lucian, Pompeius, Hesychius, Papius, Saturninus, and 
 Germanus, all natives of Italy. In the persecution of Trajan they took 
 refuge in the town of Durazzo where they saw the saintly bishop Astius 
 hanging on a cross for the faith of Christ. They then publicly 
 declared themselves to be Christians, when, by order of the governor, they 
 were arrested and cast into the sea.
\switchcolumn*
\selectlanguage{latin}
Bríxiæ sancti Apollónii, 
 Epíscopi et Confessóris.
\switchcolumn
\selectlanguage{english}
At Brescia, St. Apollonius, bishop 
 and confessor.
\switchcolumn*
\selectlanguage{latin}
Eystádii, in Germánia, sancti Willebáldi, ejúsdem urbis primi Epíscopi, qui fílius 
 éxstitit sancti 
 Richárdi, Anglórum Regis, et frater sanctæ Walbúrgæ Vírginis; atque, una cum 
 sancto Bonifátio, labórans in Evangélio, multas gentes convértit ad Christum.
\switchcolumn
\selectlanguage{english}
At Eichstadt in Germany, St. 
 Willibald, the first bishop of that city. He was the son of St. 
 Richard, king of England, and brother of St. Walburga, virgin. He 
 laboured with St. Boniface in preaching the Gospel and converted many 
 nations to Christ.
\switchcolumn*
\selectlanguage{latin}
Arvérnis, in Gállia, 
 sancti Illídii Epíscopi.
\switchcolumn
\selectlanguage{english}
In Auvergne, St. Illidius, bishop.
\switchcolumn*
\selectlanguage{latin}
Urgéllæ, in Hispánia 
 Tarraconénsi, sancti Odónis Epíscopi.
\switchcolumn
\selectlanguage{english}
At Urgal in Spain, St. Odo, bishop.
\switchcolumn*
\selectlanguage{latin}
In Anglia sancti Heddæ, 
 Epíscopi Sáxonum occidentálium.
\switchcolumn
\selectlanguage{english}
In England, St. Hedda, bishop of the 
 West Saxons.
\switchcolumn*
\selectlanguage{latin}
Alexandríæ natális 
 sancti Pantæni, viri Apostólici et omni sapiéntia adornáti, cui tantum 
 stúdii et amóris erga verbum Dei quæ in Oriéntis últimis secéssibus 
 recondúntur, fídei et devotiónis calóre succénsus, proféctus sit; ac demum, 
 Alexandríam revérsus, ibi, sub Antoníno Caracálla, in pace quiévit.
\switchcolumn
\selectlanguage{english}
At Alexandria, the birthday of St. 
 Pantaenus, a man of apostolic manner, filled with wisdom. He had such 
 an affection and love for the word of God, and was so inflamed with the 
 ardour of faith and devotion, that he set out to preach the Gospel of Christ 
 to the nations living in the farthest districts of the East. Returning 
 at last to Alexandria, he rested in peace, in the time of Antoninus 
 Caracalla.
\switchcolumn*
\selectlanguage{latin}
Eboríaci, in território Meldénsi, sanctæ Edilbúrgæ, Abbatíssæ et Vírginis, Anglórum Regis fíliæ.
\switchcolumn
\selectlanguage{english}
At Faremoutier, in the neighbourhood 
 of Meaux, St. Ethelburga, virgin, daughter of the English king.
\switchcolumn*
\selectlanguage{latin}
Perúsiæ Beáti Benedícti 
 Papæ Undécimi, Tarvisíni, ex Ordine Prædicatórum, Confessóris; qui, brevi 
 Pontificátus sui spátio, Ecclésiæ pacem, disciplínæ restauratiónem, 
 religiónis increméntum mirífice promóvit.
\switchcolumn
\selectlanguage{english}
At Perugia, blessed Pope Benedict 
 XI, a native of Treviso, of the Order of Preachers, who in the brief space 
 of his pontificate greatly promoted the peace of the Church, the restoration 
 of discipline, and the spread of religion.
\switchcolumn*
\selectlanguage{latin}
\end{paracol}


% ---- martyrology/mart07/mart0708.htm
\needspace{10\baselineskip}
\begin{paracol}{2}
\selectlanguage{latin}
\begin{center}{\color{gregoriocolor} Octávo Idus Júlii. 
 Luna\dots\ }\end{center}
\switchcolumn
\selectlanguage{english}
\begin{center}{\color{gregoriocolor} The 
 8\textsuperscript{th} Day of 
 July. The\dots\ Day of the Moon.}\end{center}
\end{paracol}

\noindent\begin{tabularx}{\linewidth}{*{19}{>{\centering\arraybackslash}X}}
 \textcolor{gregoriocolor}{a} & \textcolor{gregoriocolor}{b} & \textcolor{gregoriocolor}{c} & \textcolor{gregoriocolor}{d} & \textcolor{gregoriocolor}{e} & \textcolor{gregoriocolor}{f} & \textcolor{gregoriocolor}{g} & \textcolor{gregoriocolor}{h} & \textcolor{gregoriocolor}{i} & \textcolor{gregoriocolor}{k} & \textcolor{gregoriocolor}{l} & \textcolor{gregoriocolor}{m} & \textcolor{gregoriocolor}{n} & \textcolor{gregoriocolor}{p} & \textcolor{gregoriocolor}{q} & \textcolor{gregoriocolor}{r} & \textcolor{gregoriocolor}{s} & \textcolor{gregoriocolor}{t} & \textcolor{gregoriocolor}{u} \\
 13 & 14 & 15 & 16 & 17 & 18 & 19 & 20 & 21 & 22 & 23 & 24 & 25 & 26 & 27 & 28 & 29 & 30 & 1 \\
\end{tabularx}
\vspace{0.5\baselineskip}
\noindent\begin{tabularx}{\linewidth}{*{12}{>{\centering\arraybackslash}X}}
 \textcolor{gregoriocolor}{A} & \textcolor{gregoriocolor}{B} & \textcolor{gregoriocolor}{C} & \textcolor{gregoriocolor}{D} & \textcolor{gregoriocolor}{E} & F & \textcolor{gregoriocolor}{F} & \textcolor{gregoriocolor}{G} & \textcolor{gregoriocolor}{H} & \textcolor{gregoriocolor}{M} & \textcolor{gregoriocolor}{N} & \textcolor{gregoriocolor}{P} \\
 2 & 3 & 4 & 5 & 6 & 7 & 7 & 8 & 9 & 10 & 11 & 12 \\
\end{tabularx}

\begin{paracol}{2}
\selectlanguage{latin}
\lettrine[lines=2]{S}{anctæ} Elísabeth Víduæ, 
 Lusitanórum Regínæ, quæ ad regnum cæléste quarto Nonas hujus mensis 
 transívit.
\switchcolumn
\selectlanguage{english}
\lettrine[lines=2]{S}{t.} Elisabeth, widow, queen of 
 Portugal, whose birthday is observed on the 4th of July.
\switchcolumn*
\selectlanguage{latin}
In Asia minóre 
 sanctórum Aquilæ et Priscíllæ uxóris, de quibus in Actibus Apostolórum 
 scríbitur.
\switchcolumn
\selectlanguage{english}
In Asia Minor, the Saints Aquilla 
 and his wife Priscilla, of whom mention is made in the Acts of the Apostles.
\switchcolumn*
\selectlanguage{latin}
Herbípoli, in Germánia, 
 sancti Chiliáni Epíscopi, qui, a Románo Pontífice ad prædicándum Evangélium 
 missus, ibi, cum multos ad Christum perduxísset, una cum Sóciis Colománno 
 Presbytero et Totnáno Diácono, trucidátus est.
\switchcolumn
\selectlanguage{english}
At Wurtzburg in Germany, St. Kilian, 
 bishop, who was commissioned by the Roman Pontiff to preach the Gospel. 
 After having converted many to Christ, he was put to death along with his 
 companions Colman, a priest, and Totnan, a deacon.
\switchcolumn*
\selectlanguage{latin}
In Portu Románo 
 sanctórum quinquagínta mílitum Mártyrum, qui, in sanctæ Bonósæ confessióne 
 ad fidem addúcti, et a beáto Felíce Papa Primo baptizáti, in Aureliáno 
 Imperatóris persecutióne occísi sunt.
\switchcolumn
\selectlanguage{english}
At Porto, fifty holy martyrs, all 
 soldiers, who were led to the faith by the martyrdom of St. Bonosa, and 
 baptized by the blessed Pope Felix. They were put to death in the 
 persecution of Aurelian.
\switchcolumn*
\selectlanguage{latin}
Cæsaréæ, in Palæstína, 
 sancti Procópii Mártyris, qui, sub Diocletiáno Imperatóre, a Scythópoli 
 ductus Cæsaréam, illic, ad primam responsiónum ejus confidéntiam, a Júdice 
 Fabiáno est cápite cæsus.
\switchcolumn
\selectlanguage{english}
In Palestine, in the reign of 
 Diocletian, St. Procopius, martyr, who was brought from Scythopolis to 
 Caesarea, and upon his first resolute answer was beheaded by the judge 
 Fabian.
\switchcolumn*
\selectlanguage{latin}
Constantinópoli pássio 
 sanctórum Monachórum Abrahamitárum, qui ob cultum sanctárum Imáginum, 
 resisténtes Theóphilo Imperatóri, martyrium consummárunt.
\switchcolumn
\selectlanguage{english}
At Constantinople, the holy 
 Abrahamite monks, who resisted Emperor Theophilus by defending the 
 veneration of sacred images, and suffered martyrdom.
\switchcolumn*
\selectlanguage{latin}
Spinæ Lambérti, in 
 Æmília, sancti Hadriáni Papæ Tértii, stúdio conciliándi Románæ Ecclésias 
 Orientáles insígnis, ac miráculis clari; cujus corpus in monastérium 
 Nonantulénse fuit delátum, et in æde sancti Silvéstri honorífice cónditum.
\switchcolumn
\selectlanguage{english}
At Spina Lamberti in Emília, Pope 
 St. Adrian III, famous for his zeal in reconciling the Eastern to the Roman 
 Church, and renowned for his miracles. His body was taken to the 
 monastery of Nonantola and buried with honours in the Church of St. 
 Sylvester.
\switchcolumn*
\selectlanguage{latin}
Tréviris sancti 
 Auspícii, Epíscopi et Confessóris.
\switchcolumn
\selectlanguage{english}
At Treves, St. Auspicius, bishop and 
 confessor.
\switchcolumn*
\selectlanguage{latin}
Romæ beáti Eugénii Papæ 
 Tértii, qui, postquam cœnóbium sanctórum Vincéntii et Anastásii ad Aquas 
 Sálvias magna sanctimóniæ ac prudéntiæ laude rexísset, Ecclésiam univérsam, 
 Póntifex Máximus renuntiátus, sanctíssime gubernávit. Cultum autem, ab 
 immemorábili témpore ipsi exhíbitum, Pius Papa Nonus ratum hábuit et 
 confirmávit.
\switchcolumn
\selectlanguage{english}
At Rome, blessed Eugene III, pope. 
 Having gained a great reputation for sanctity and prudence in his government 
 of the monastery of Saints Vincent and Anastasius, he was raised to the 
 Sovereign Pontificate and ruled the universal Church in much holiness. 
 Pope Pius IX approved and confirmed the veneration paid to him.
\switchcolumn*
\selectlanguage{latin}
\end{paracol}


% ---- martyrology/mart07/mart0709.htm
\needspace{10\baselineskip}
\begin{paracol}{2}
\selectlanguage{latin}
\begin{center}{\color{gregoriocolor} Séptimo Idus Júlii. 
 Luna\dots\ }\end{center}
\switchcolumn
\selectlanguage{english}
\begin{center}{\color{gregoriocolor} The 
 9\textsuperscript{th} Day of 
 July. The\dots\ Day of the Moon.}\end{center}
\end{paracol}

\noindent\begin{tabularx}{\linewidth}{*{19}{>{\centering\arraybackslash}X}}
 \textcolor{gregoriocolor}{a} & \textcolor{gregoriocolor}{b} & \textcolor{gregoriocolor}{c} & \textcolor{gregoriocolor}{d} & \textcolor{gregoriocolor}{e} & \textcolor{gregoriocolor}{f} & \textcolor{gregoriocolor}{g} & \textcolor{gregoriocolor}{h} & \textcolor{gregoriocolor}{i} & \textcolor{gregoriocolor}{k} & \textcolor{gregoriocolor}{l} & \textcolor{gregoriocolor}{m} & \textcolor{gregoriocolor}{n} & \textcolor{gregoriocolor}{p} & \textcolor{gregoriocolor}{q} & \textcolor{gregoriocolor}{r} & \textcolor{gregoriocolor}{s} & \textcolor{gregoriocolor}{t} & \textcolor{gregoriocolor}{u} \\
 14 & 15 & 16 & 17 & 18 & 19 & 20 & 21 & 22 & 23 & 24 & 25 & 26 & 27 & 28 & 29 & 30 & 1 & 2 \\
\end{tabularx}
\vspace{0.5\baselineskip}
\noindent\begin{tabularx}{\linewidth}{*{12}{>{\centering\arraybackslash}X}}
 \textcolor{gregoriocolor}{A} & \textcolor{gregoriocolor}{B} & \textcolor{gregoriocolor}{C} & \textcolor{gregoriocolor}{D} & \textcolor{gregoriocolor}{E} & F & \textcolor{gregoriocolor}{F} & \textcolor{gregoriocolor}{G} & \textcolor{gregoriocolor}{H} & \textcolor{gregoriocolor}{M} & \textcolor{gregoriocolor}{N} & \textcolor{gregoriocolor}{P} \\
 3 & 4 & 5 & 6 & 7 & 8 & 8 & 9 & 10 & 11 & 12 & 13 \\
\end{tabularx}

\begin{paracol}{2}
\selectlanguage{latin}
\lettrine[lines=2]{R}{omæ} ad Guttam júgiter 
 manéntem, natális sanctórum Mártyrum Zenónis, et aliórum decem míllium ac 
 ducentórum trium.
\switchcolumn
\selectlanguage{english}
\lettrine[lines=2]{A}{t} Rome, at the Ever-flowing Spring, 
 the birthday of St. Zeno and ten thousand two hundred and three other 
 martyrs.
\switchcolumn*
\selectlanguage{latin}
Gortynæ, in Creta, 
 sancti Cyrílli Epíscopi, qui, in persecutióne Décii, sub Lúcio Præside, 
 flammis est injéctus, et, cum ab igne, incénsis vínculis, illæsus evasísset, 
 ac stupóre tanti miráculi a Júdice dimíssus esset, rursus ab eódem, pro 
 instánti et álacri fídei prædicatióne facta de Christo, comprehénsus et 
 cápite plexus est.
\switchcolumn
\selectlanguage{english}
At Gortyna in Crete, in the 
 persecution of Decius, under the governor Lucius, Bishop St. Cyril. 
 When he was thrown into the flames his bonds were burned, but he was not 
 injured. The judge, struck with awe at so great a miracle, set him at 
 liberty, but as the saint began again immediately to preach with zeal the 
 faith of Christ, he was beheaded.
\switchcolumn*
\selectlanguage{latin}
Brilæ, in Hollándia, 
 pássio novémdecim Mártyrum, Gorcomiénsium nuncupatórum; quorum ex número 
 novem Sacerdótes ac duo Láici erant Fratres Minóres, quátuor Presbyteri 
 sæculáres, duo Præmonstraténses, unus Reguláris Canónicus sancti Augústíni, 
 et unus Dominicánus. Hi omnes, ob tuéndam Ecclésiæ Románæ auctoritátem 
 et reálem Christi in Eucharístia præséntiam, a Calviniánis hæréticis varia 
 ludíbria et torménta perpéssi, tandem, in trabem acti, agónem suum 
 adstríctis láqueo fáucibus, consummárunt; et a Pio Nono, Pontífice Máximo, 
 inter sanctos Mártyres reláti sunt.
\switchcolumn
\selectlanguage{english}
At Briel in Holland, the passion of 
 the nineteen martyrs of Gorcum. Of these, nine priests and two lay 
 brothers were of the Order of Friars Minor, four were secular priests, two 
 Premonstratensians, one Canon Regular of St. Augustine, and one Dominican. 
 For vindicating the authority of the Roman Church and the real presence of 
 Christ in the Eucharist, they endured various insults and torments from the 
 Calvinist heretics, and their great suffering was ended by all of them being 
 hanged. Pope Pius IX included them in the number of holy martyrs.
\switchcolumn*
\selectlanguage{latin}
In civitáte Thora, apud 
 lacum Velínum, item pássio sanctórum Anatóliæ et Audácis, sub Décio 
 Imperatóre. Ex his Anatóliæ, Christi Virgo, postquam plúrimos per 
 totam Picéni provínciam váriis languóribus afféctos curásset et in Christum 
 credéntes fecísset, Júdicis Faustiniáni jussu divérsis pœnárum genéribus est 
 vexáta, et, cum ab immísso serpénte líbera evádens Audácem convertísset ad 
 fidem, novíssime, exténsis mánibus orans, gládio transverberáta est; Audax 
 quoque, in custódiam tráditus, sine mora senténtia capitáli coronátur.
\switchcolumn
\selectlanguage{english}
In the town of Thora, on Lake Velino 
 in Italy, the martyrdom of the Saints Anatolia and Audax, under Emperor 
 Decius. Anatolia, a virgin consecrated to Christ, cured many persons 
 afflicted with various infirmities throughout the province of Piceno, and 
 made them believe in Christ. By order of the judge Faustinian she was 
 condemned to different kinds of punishment. She was cured of the sting 
 of a serpent to which she had been exposed, a miracle that converted Audax 
 to the faith. At last, praying with outstretched hands, she was 
 pierced with a sword. Audax was sent to prison, and without delay 
 sentenced to capital punishment, thus obtaining the crown of martyrdom.
\switchcolumn*
\selectlanguage{latin}
Alexandríæ sanctórum 
 Mártyrum Patermúthii, Coprétis et Alexándri; qui sub Juliáno Apóstata cæsi 
 sunt.
\switchcolumn
\selectlanguage{english}
At Alexandria, the holy martyrs 
 Patermuthius, Copres, and Alexander, who were put to death under Julian the 
 Apostate.
\switchcolumn*
\selectlanguage{latin}
Mártulæ, in Umbria, 
 sancti Bríctii Epíscopi, qui, sub Marciáno Júdice, ob confessiónem Dómini, 
 multa passus est; ac tandem, cum magnam pópuli multitúdinem ad Christum 
 convertísset, Conféssor in pace quiévit.
\switchcolumn
\selectlanguage{english}
At Martula in Umbria, St. Brictius, 
 bishop. Under the judge Marcian, after having suffered much for the 
 confession of our Lord, and having converted to Christ a great multitude of 
 people, he rested in peace, a confessor.
\switchcolumn*
\selectlanguage{latin}
Tiférni, in Umbria, 
 sanctæ Verónicæ de Juliánis, Vírginis, in Urbaniénsis diœcésis óppido 
 Mercatéllo natæ, Moniális e secúndo sancti Francísci Ordine ac Tifernátis 
 ascetérii Abbatíssæ; quam, insígni patiéndi stúdio, ceterísque virtútibus et 
 cæléstibus charismátibus illústrem, Gregórius Papa Décimus sextus in 
 sanctárum Vírginum collégium adscrípsit.
\switchcolumn
\selectlanguage{english}
At Tiferno in Umbria, St. Veronica 
 Giuliani, a nun of the second Order of St. Francis and abbess of the 
 monastery in that town. Born at Mercatello in the diocese of Urbania, 
 she became illustrious by her great love for suffering and other virtues, 
 and by her heavenly gifts. She was inscribed among the holy virgins by 
 Pope Gregory XVI.
\switchcolumn*
\selectlanguage{latin}
\end{paracol}


% ---- martyrology/mart07/mart0710.htm
\needspace{10\baselineskip}
\begin{paracol}{2}
\selectlanguage{latin}
\begin{center}{\color{gregoriocolor} Sexto Idus Júlii. 
 Luna\dots\ }\end{center}
\switchcolumn
\selectlanguage{english}
\begin{center}{\color{gregoriocolor} The 
 10\textsuperscript{th} Day of 
 July. The\dots\ Day of the Moon.}\end{center}
\end{paracol}

\noindent\begin{tabularx}{\linewidth}{*{19}{>{\centering\arraybackslash}X}}
 \textcolor{gregoriocolor}{a} & \textcolor{gregoriocolor}{b} & \textcolor{gregoriocolor}{c} & \textcolor{gregoriocolor}{d} & \textcolor{gregoriocolor}{e} & \textcolor{gregoriocolor}{f} & \textcolor{gregoriocolor}{g} & \textcolor{gregoriocolor}{h} & \textcolor{gregoriocolor}{i} & \textcolor{gregoriocolor}{k} & \textcolor{gregoriocolor}{l} & \textcolor{gregoriocolor}{m} & \textcolor{gregoriocolor}{n} & \textcolor{gregoriocolor}{p} & \textcolor{gregoriocolor}{q} & \textcolor{gregoriocolor}{r} & \textcolor{gregoriocolor}{s} & \textcolor{gregoriocolor}{t} & \textcolor{gregoriocolor}{u} \\
 15 & 16 & 17 & 18 & 19 & 20 & 21 & 22 & 23 & 24 & 25 & 26 & 27 & 28 & 29 & 30 & 1 & 2 & 3 \\
\end{tabularx}
\vspace{0.5\baselineskip}
\noindent\begin{tabularx}{\linewidth}{*{12}{>{\centering\arraybackslash}X}}
 \textcolor{gregoriocolor}{A} & \textcolor{gregoriocolor}{B} & \textcolor{gregoriocolor}{C} & \textcolor{gregoriocolor}{D} & \textcolor{gregoriocolor}{E} & F & \textcolor{gregoriocolor}{F} & \textcolor{gregoriocolor}{G} & \textcolor{gregoriocolor}{H} & \textcolor{gregoriocolor}{M} & \textcolor{gregoriocolor}{N} & \textcolor{gregoriocolor}{P} \\
 4 & 5 & 6 & 7 & 8 & 9 & 9 & 10 & 11 & 12 & 13 & 14 \\
\end{tabularx}

\begin{paracol}{2}
\selectlanguage{latin}
\lettrine[lines=2]{R}{omæ} pássio sanctórum 
 septem Mártyrum fratrum, filiórum sanctæ Felicitátis Mártyris, id est 
 Januárii, Felícis, Philíppi, Silváni, Alexándri, Vitális et Martiális, 
 témpore Antoníni Imperatóris, sub Præfécto Urbis Públio. Ex ipsis vero 
 Januárius, post virgárum vérbera et cárceris maceratiónem, plumbátis occísus; 
 Felix et Philíppus fústibus mactáti; Silvánus præcipítio interémptus; 
 Alexánder, Vitális et Martiális capitáli senténtia puníti sunt.
\switchcolumn
\selectlanguage{english}
\lettrine[lines=2]{A}{t} Rome, the martyrdom of the seven 
 holy brothers, sons of the saintly martyr Felicitas. They are 
 Januarius, Felix, Philip, Sylvanus, Alexander, Vitalis, and Martial. 
 They died in the time of Emperor Antoninus, under Publius, prefect of the 
 city. Januarius, after being scourged with rods and detained in 
 prison, died from the blows inflicted with leaded whips. Felix and 
 Philip were scourged to death. Sylvanus was thrown headlong from a 
 great height. Alexander, Vitalis, and Martial were beheaded.
\switchcolumn*
\selectlanguage{latin}
Item Romæ sanctárum 
 Vírginum et Mártyrum Rufínæ et Secúndæ sorórum, quæ, in persecutióne 
 Valeriáni et Galliéni, torméntis subáctæ sunt; atque ad últimum, cum álteri 
 caput illísum esset gládio, et álteri cervix cæsa, migrárunt in cælum. 
 Ipsárum vero córpora in Basílica Lateranénsi, prope Baptistérium, débito 
 honóre servántur.
\switchcolumn
\selectlanguage{english}
Also at Rome, in the persecution of 
 Valerian and Gallienus, the holy virgins and martyrs Rufina and Secunda, 
 sisters. After being subjected to torments, and one having her head 
 crushed with a sword, the other beheaded, they departed for heaven. 
 Their bodies are kept with due honour in the Lateran basilica, near the baptistry.
\switchcolumn*
\selectlanguage{latin}
In Africa sanctórum 
 Mártyrum Januárii, Maríni, Nabóris et Felícis, qui gládio decolláti sunt.
\switchcolumn
\selectlanguage{english}
In Africa, the holy martyrs 
 Januarius, Marinus, Nabor and Felix, all of whom were beheaded.
\switchcolumn*
\selectlanguage{latin}
Nicópoli, in Arménia, sanctórum Mártyrum Leóntii, Maurítii, Daniélis, et Sociórum; qui, sub 
 Licínio Imperatóre et Lysia Præside, várie excruciáti, tandem, in ignem 
 conjécti, martyrii cursum confecérunt.
\switchcolumn
\selectlanguage{english}
At Nicopolis in Armenia, the holy 
 martyrs Leontius, Mauritius, Daniel, and their companions, who were tortured 
 in different ways, and being lastly cast into the fire, ended their long 
 martyrdom in the time of Emperor Licinius and the governor Lysias.
\switchcolumn*
\selectlanguage{latin}
In Pisídia sanctórum 
 Mártyrum Biánoris et Silváni, qui, sævíssima pro Christi nómine passi, demum, 
 cervícibus abscíssis, coronántur.
\switchcolumn
\selectlanguage{english}
In Pisidia, the holy martyrs Bianor 
 and Silvanus, who were merited an immortal crown by being beheaded, after 
 enduring most bitter torments for the name of Christ.
\switchcolumn*
\selectlanguage{latin}
Icónii, in Lycaónia, 
 sancti Apollónii Mártyris, qui per crucem insígne martyrium consummávit.
\switchcolumn
\selectlanguage{english}
At Iconium, St. Apollonius, martyr, 
 whose glorious martyrdom was fulfilled by death on the cross.
\switchcolumn*
\selectlanguage{latin}
Apud Gandávum, in 
 Flándria, sanctæ Amelbérgæ Vírginis.
\switchcolumn
\selectlanguage{english}
At Ghent in Flanders, St. Amelberga, 
 virgin.
\switchcolumn*
\selectlanguage{latin}
\end{paracol}


% ---- martyrology/mart07/mart0711.htm
\needspace{10\baselineskip}
\begin{paracol}{2}
\selectlanguage{latin}
\begin{center}{\color{gregoriocolor} Quinto Idus Júlii. 
 Luna\dots\ }\end{center}
\switchcolumn
\selectlanguage{english}
\begin{center}{\color{gregoriocolor} The 
 11\textsuperscript{th} Day of 
 July. The\dots\ Day of the Moon.}\end{center}
\end{paracol}

\noindent\begin{tabularx}{\linewidth}{*{19}{>{\centering\arraybackslash}X}}
 \textcolor{gregoriocolor}{a} & \textcolor{gregoriocolor}{b} & \textcolor{gregoriocolor}{c} & \textcolor{gregoriocolor}{d} & \textcolor{gregoriocolor}{e} & \textcolor{gregoriocolor}{f} & \textcolor{gregoriocolor}{g} & \textcolor{gregoriocolor}{h} & \textcolor{gregoriocolor}{i} & \textcolor{gregoriocolor}{k} & \textcolor{gregoriocolor}{l} & \textcolor{gregoriocolor}{m} & \textcolor{gregoriocolor}{n} & \textcolor{gregoriocolor}{p} & \textcolor{gregoriocolor}{q} & \textcolor{gregoriocolor}{r} & \textcolor{gregoriocolor}{s} & \textcolor{gregoriocolor}{t} & \textcolor{gregoriocolor}{u} \\
 16 & 17 & 18 & 19 & 20 & 21 & 22 & 23 & 24 & 25 & 26 & 27 & 28 & 29 & 30 & 1 & 2 & 3 & 4 \\
\end{tabularx}
\vspace{0.5\baselineskip}
\noindent\begin{tabularx}{\linewidth}{*{12}{>{\centering\arraybackslash}X}}
 \textcolor{gregoriocolor}{A} & \textcolor{gregoriocolor}{B} & \textcolor{gregoriocolor}{C} & \textcolor{gregoriocolor}{D} & \textcolor{gregoriocolor}{E} & F & \textcolor{gregoriocolor}{F} & \textcolor{gregoriocolor}{G} & \textcolor{gregoriocolor}{H} & \textcolor{gregoriocolor}{M} & \textcolor{gregoriocolor}{N} & \textcolor{gregoriocolor}{P} \\
 5 & 6 & 7 & 8 & 9 & 10 & 10 & 11 & 12 & 13 & 14 & 15 \\
\end{tabularx}

\begin{paracol}{2}
\selectlanguage{latin}
\lettrine[lines=2]{R}{omæ} sancti Pii Primi, 
 Papæ et Mártyris; qui martyrio coronátus est in persecutióne Marci Aurélii 
 Antoníni.
\switchcolumn
\selectlanguage{english}
\lettrine[lines=2]{A}{t} Rome, Pope Pius I, who was 
 crowned with martyrdom in the persecution of Marcus Aurelius Antoninus.
\switchcolumn*
\selectlanguage{latin}
Bérgomi sancti Joánnis 
 Epíscopi, qui, ob tuéndam cathólicam fidem, ab Ariánis occísus est.
\switchcolumn
\selectlanguage{english}
At Bergamo, St. John, a bishop, who 
 was killed by the Arians for defending the Catholic faith.
\switchcolumn*
\selectlanguage{latin}
Sidæ, in Pamphylia, 
 sancti Cindéi Presbyteri, qui, sub Diocletiáno Imperatóre et Stratoníco 
 Præside, post multa torménta, injéctus in ignem et nil læsus, demum in 
 oratióne réddidit spíritum.
\switchcolumn
\selectlanguage{english}
At Sida in Pamphylia, St. Cindeus, 
 priest, in the time of Emperor Diocletian and the governor Stratonicus. 
 After suffering many torments, he was thrown into the fire, but was not 
 injured by it. He later yielded up his soul in prayer.
\switchcolumn*
\selectlanguage{latin}
Córdubæ, in Hispánia, 
 sancti Abúndii Presbyteri, qui, in persecutióne Arábica, cum in Mahumétis 
 sectam inveherétur, martyrio coronátus est.
\switchcolumn
\selectlanguage{english}
At Cordova in Spain, St. Abundius, a 
 priest, crowned with martyrdom while preaching against the sect of Mohammed.
\switchcolumn*
\selectlanguage{latin}
Nicópoli, in Arménia, natális sanctórum Mártyrum Januárii et Pelágiæ, qui, equúleo, 
 úngulis et 
 testárum fragméntis per dies quátuor cruciáti, martyrium implevérunt.
\switchcolumn
\selectlanguage{english}
At Nicopolis in Armenia, the 
 birthday of the holy martyrs Januarius and Pelagia, who for four days were 
 racked, torn with iron claws and pieces of earthenware, and thus achieved 
 their martyrdom.
\switchcolumn*
\selectlanguage{latin}
In território Senonénsi 
 sancti Sidrónii Mártyris.
\switchcolumn
\selectlanguage{english}
In the territory of Sens, St. 
 Sidronius, martyr.
\switchcolumn*
\selectlanguage{latin}
Icónii, in Lycaónia, 
 sancti Marciáni Mártyris, qui, sub Perénnio Præside, per multa torménta 
 pervénit ad palmam.
\switchcolumn
\selectlanguage{english}
At Iconium in Lycaonia, St. Marcian, 
 martyr. He obtained the palm of martyrdom by many torments, under the 
 governor Perennius.
\switchcolumn*
\selectlanguage{latin}
Bríxiæ sanctórum 
 Mártyrum Savíni et Cypriáni.
\switchcolumn
\selectlanguage{english}
At Brescia, the holy martyrs Savinus 
 and Cyprian.
\switchcolumn*
\selectlanguage{latin}
In território 
 Pictaviénsi sancti Sabíni Confessóris.
\switchcolumn
\selectlanguage{english}
In the territory of Poitiers, St. 
 Sabinus, confessor.
\switchcolumn*
\selectlanguage{latin}
\end{paracol}


% ---- martyrology/mart07/mart0712.htm
\needspace{10\baselineskip}
\begin{paracol}{2}
\selectlanguage{latin}
\begin{center}{\color{gregoriocolor} Quarto Idus Júlii. 
 Luna\dots\ }\end{center}
\switchcolumn
\selectlanguage{english}
\begin{center}{\color{gregoriocolor} The 
 12\textsuperscript{th} Day of 
 July. The\dots\ Day of the Moon.}\end{center}
\end{paracol}

\noindent\begin{tabularx}{\linewidth}{*{19}{>{\centering\arraybackslash}X}}
 \textcolor{gregoriocolor}{a} & \textcolor{gregoriocolor}{b} & \textcolor{gregoriocolor}{c} & \textcolor{gregoriocolor}{d} & \textcolor{gregoriocolor}{e} & \textcolor{gregoriocolor}{f} & \textcolor{gregoriocolor}{g} & \textcolor{gregoriocolor}{h} & \textcolor{gregoriocolor}{i} & \textcolor{gregoriocolor}{k} & \textcolor{gregoriocolor}{l} & \textcolor{gregoriocolor}{m} & \textcolor{gregoriocolor}{n} & \textcolor{gregoriocolor}{p} & \textcolor{gregoriocolor}{q} & \textcolor{gregoriocolor}{r} & \textcolor{gregoriocolor}{s} & \textcolor{gregoriocolor}{t} & \textcolor{gregoriocolor}{u} \\
 17 & 18 & 19 & 20 & 21 & 22 & 23 & 24 & 25 & 26 & 27 & 28 & 29 & 30 & 1 & 2 & 3 & 4 & 5 \\
\end{tabularx}
\vspace{0.5\baselineskip}
\noindent\begin{tabularx}{\linewidth}{*{12}{>{\centering\arraybackslash}X}}
 \textcolor{gregoriocolor}{A} & \textcolor{gregoriocolor}{B} & \textcolor{gregoriocolor}{C} & \textcolor{gregoriocolor}{D} & \textcolor{gregoriocolor}{E} & F & \textcolor{gregoriocolor}{F} & \textcolor{gregoriocolor}{G} & \textcolor{gregoriocolor}{H} & \textcolor{gregoriocolor}{M} & \textcolor{gregoriocolor}{N} & \textcolor{gregoriocolor}{P} \\
 6 & 7 & 8 & 9 & 10 & 11 & 11 & 12 & 13 & 14 & 15 & 16 \\
\end{tabularx}

\begin{paracol}{2}
\selectlanguage{latin}
\lettrine[lines=2]{I}{n} monastério 
 Passiniáno, prope Floréntiam, sancti Joánnis Gualbérti Abbátis, qui fuit 
 Institútor Ordinis Vallis Umbrósæ.
\switchcolumn
\selectlanguage{english}
\lettrine[lines=2]{I}{n} the monastery of Passignano, near 
 Florence, Abbot St. John Gualbert, founder of the Order of Vallombrosa.
\switchcolumn*
\selectlanguage{latin}
Laudæ, in Insúbria, 
 sanctórum Mártyrum Náboris et Felícis, qui, in persecutióne Maximiáni, post 
 vária torménta, cápitis decollatióne martyrium complevérunt; eorúmque 
 córpora a beáta Savína Mediolánum advécta sunt, ibíque honorífice sepúlta.
\switchcolumn
\selectlanguage{english}
At Milan, the holy martyrs Nabor and 
 Felix, who suffered in the persecution of Maximian. Their bodies were 
 brought into the city by blessed Savina, and were honourably buried there.
\switchcolumn*
\selectlanguage{latin}
In Cypro beáti Jasónis, 
 qui fuit unus ex antíquis Christi discípulis.
\switchcolumn
\selectlanguage{english}
In Cyprus, St. Jason, one of the 
 first disciples of Christ.
\switchcolumn*
\selectlanguage{latin}
Lucæ, in Túscia, beáti 
 Paulíni, qui, a sancto Petro Apóstolo primus ejúsdem civitátis Epíscopus est 
 ordinátus; et, sub Neróne, ad radíces montis Pisáni, post multos agónes, 
 martyrium suum cum áliis Sóciis consummávit.
\switchcolumn
\selectlanguage{english}
At Lucca in Tuscany, blessed 
 Paulinus, who was consecrated first bishop of that city by St. Peter. 
 Under Nero he completed his martyrdom along with many others at the foot of 
 Mt. Pisa, but only after many trials.
\switchcolumn*
\selectlanguage{latin}
Aquiléjæ natális sancti 
 Hermágoræ, qui éxstitit discípulus beáti Marci Evangelístæ, et primus 
 ejúsdem civitátis Epíscopus. Hic, inter mirácula sanitátum, et 
 prædicatiónis instántiam, ac populórum conversiónem, plúrima pœnárum génera 
 expértus, tandem, una cum Diácono suo Fortunáto, capitáli supplício 
 perpétuum méruit obtinére triúmphum.
\switchcolumn
\selectlanguage{english}
At Aquileia, the birthday of St. 
 Hermagoras, disciple of the blessed evangelist Mark, and first bishop of 
 that city. When performing miraculous cures, or while preaching, 
 frequently bringing souls to repentance, he suffered many torments. 
 Finally by capital punishment her merited an immortal triumph along with his 
 deacon Fortunatus.
\switchcolumn*
\selectlanguage{latin}
Eódem die pássio 
 sanctórum Proci et Hilariónis, qui, sub Trajáno Imperatóre et Máximo Præside, 
 per acerbíssima torménta ad palmam martyrii pervenérunt.
\switchcolumn
\selectlanguage{english}
The same day, the Saints Proclus and 
 Hilarion, who won the palm of martyrdom after most bitter torments, in the 
 time of Emperor Trajan and the governor Maximus.
\switchcolumn*
\selectlanguage{latin}
Toléti, in Hispánia, 
 sanctæ Marciánæ, Vírginis et Mártyris; quæ, pro fide Christi, objécta 
 béstiis atque a tauro discérpta, martyrio coronátur.
\switchcolumn
\selectlanguage{english}
At Toledo in Spain, St. Marciana, 
 virgin and martyr. For the faith of Christ, she was cast to the 
 beasts, torn to pieces by a bull, and was thus crowned with martyrdom.
\switchcolumn*
\selectlanguage{latin}
Apud Leontínos, in 
 Sicília, sanctæ Epíphanæ, quæ sub Diocletiáno Imperatóre et Tertyllo Præside, 
 ubéribus præcísis, réddidit spíritum.
\switchcolumn
\selectlanguage{english}
At Lentini, St. Epiphana, who, after 
 her breasts were cut away, died in the time of Emperor Diocletian and the 
 governor Tertillus.
\switchcolumn*
\selectlanguage{latin}
Lugdúni, in Gállia, 
 sancti Viventíoli Epíscopi.
\switchcolumn
\selectlanguage{english}
At Lyons in France, St. Viventiolus,bishop.
\switchcolumn*
\selectlanguage{latin}
Bonóniæ sancti 
 Paterniáni Epíscopi.
\switchcolumn
\selectlanguage{english}
At Bologna, St. Paternian, bishop.
\switchcolumn*
\selectlanguage{latin}
\end{paracol}


% ---- martyrology/mart07/mart0713.htm
\needspace{10\baselineskip}
\begin{paracol}{2}
\selectlanguage{latin}
\begin{center}{\color{gregoriocolor} Tértio Idus Júlii. 
 Luna\dots\ }\end{center}
\switchcolumn
\selectlanguage{english}
\begin{center}{\color{gregoriocolor} The 
 13\textsuperscript{th} Day of 
 July. The\dots\ Day of the Moon.}\end{center}
\end{paracol}

\noindent\begin{tabularx}{\linewidth}{*{19}{>{\centering\arraybackslash}X}}
 \textcolor{gregoriocolor}{a} & \textcolor{gregoriocolor}{b} & \textcolor{gregoriocolor}{c} & \textcolor{gregoriocolor}{d} & \textcolor{gregoriocolor}{e} & \textcolor{gregoriocolor}{f} & \textcolor{gregoriocolor}{g} & \textcolor{gregoriocolor}{h} & \textcolor{gregoriocolor}{i} & \textcolor{gregoriocolor}{k} & \textcolor{gregoriocolor}{l} & \textcolor{gregoriocolor}{m} & \textcolor{gregoriocolor}{n} & \textcolor{gregoriocolor}{p} & \textcolor{gregoriocolor}{q} & \textcolor{gregoriocolor}{r} & \textcolor{gregoriocolor}{s} & \textcolor{gregoriocolor}{t} & \textcolor{gregoriocolor}{u} \\
 18 & 19 & 20 & 21 & 22 & 23 & 24 & 25 & 26 & 27 & 28 & 29 & 30 & 1 & 2 & 3 & 4 & 5 & 6 \\
\end{tabularx}
\vspace{0.5\baselineskip}
\noindent\begin{tabularx}{\linewidth}{*{12}{>{\centering\arraybackslash}X}}
 \textcolor{gregoriocolor}{A} & \textcolor{gregoriocolor}{B} & \textcolor{gregoriocolor}{C} & \textcolor{gregoriocolor}{D} & \textcolor{gregoriocolor}{E} & F & \textcolor{gregoriocolor}{F} & \textcolor{gregoriocolor}{G} & \textcolor{gregoriocolor}{H} & \textcolor{gregoriocolor}{M} & \textcolor{gregoriocolor}{N} & \textcolor{gregoriocolor}{P} \\
 7 & 8 & 9 & 10 & 11 & 12 & 12 & 13 & 14 & 15 & 16 & 17 \\
\end{tabularx}

\begin{paracol}{2}
\selectlanguage{latin}
\lettrine[lines=2]{R}{omæ} sancti Anacléti, 
 Papæ et Mártyris, qui, post sanctum Cleméntem Ecclésiam Dei regens, eam 
 glorióso martyrio decorávit.
\switchcolumn
\selectlanguage{english}
\lettrine[lines=2]{A}{t} Rome, St. Anacletus, pope and 
 martyr, who governed the Church of God after St. Clement, and shed lustre 
 upon it by a glorious martyrdom.
\switchcolumn*
\selectlanguage{latin}
Bambérgæ natális sancti 
 Henríci Primi, Imperatóris Romanórum et Confessóris, qui cum sancta 
 Cunegúnde, uxóre sua, perpétuam virginitátem servávit, et sanctum Stéphanum, 
 Hungarórum Regem, cum univérso fere ipsíus regno, ad fidem Christi 
 suscipiéndam perdúxit. Ejus autem festívitas Idibus mensis hujus 
 celebrátur.
\switchcolumn
\selectlanguage{english}
At Bamberg, the birthday of the 
 Roman emperor St. Henry I, confessor. He led a life of perpetual 
 virginity with his wife St. Cunegunde, and converted St. Stephen, king of 
 Hungary,and almost all his people to the faith of Christ. His festival 
 is celebrated on the 15th of July.
\switchcolumn*
\selectlanguage{latin}
In Palæstína sanctórum 
 Joélis et Esdræ Prophetárum.
\switchcolumn
\selectlanguage{english}
In Palestine the holy prophets Joel 
 and Esdras.
\switchcolumn*
\selectlanguage{latin}
In Macedónia beáti Silæ, 
 qui, cum esset unus de primis frátribus et ab Apóstolis ad Ecclésias Géntium, 
 una cum Paulo et Bárnaba, destinátus, prædicatiónis offícium, grátia Dei 
 plenus, instánter consummávit, atque, in passiónibus suis Christum 
 claríficans, póstmodum requiévit.
\switchcolumn
\selectlanguage{english}
In Macedonia, blessed Silas, one of 
 the first Christians. The apostles sent him with Paul and Barnabas to 
 the churches of the gentiles. Filled with the grace of God, he 
 zealously discharged the office of preaching, and after glorifying Christ by 
 his sufferings, rested in peace.
\switchcolumn*
\selectlanguage{latin}
Item sancti Serapiónis 
 Mártyris, qui, sub Sevéro Imperatóre et Aquila Præside, per ignem pervénit 
 ad corónam martyrii.
\switchcolumn
\selectlanguage{english}
Also, St. Serapion, martyr, who 
 obtained the crown of martyrdom by fire, in the time of Emperor Severus and 
 the governor Aquila.
\switchcolumn*
\selectlanguage{latin}
In ínsula Chio sanctæ 
 Myrópis Mártyris, quæ sub Décio Imperatóre et Numeriáno Præside, véctibus 
 contúsa, migrávit ad Dóminum.
\switchcolumn
\selectlanguage{english}
In the island of Chio, in the time 
 of Emperor Decius and the governor Numerian, the martyr St. Myrope. 
 She went to the Lord after being beaten with clubs.
\switchcolumn*
\selectlanguage{latin}
In Africa sanctórum 
 Confessórum Eugénii, Carthaginénsis Epíscopi, fide ac virtútibus gloriósi, 
 et univérsi cleri ejúsdem Ecclésiæ, qui fere quingénti, vel eo ámplius, in 
 persecutióne Wandálica, sub Ariáno Rege Hunneríco, cæde inediáque maceráti 
 (inter quos plúrimi erant Lectóres infántuli), gaudéntes in Dómino, exsílio 
 crudéli procul extrúsi sunt. Erant étiam nobilíssimi in eis 
 Archidiáconus, nómine Salutáris, et Murítta, secúndus in offício ministrórum; 
 qui, tértio Confessóres effécti, ambo glorióse in Christo, perseverántiæ 
 título, illustráti sunt.
\switchcolumn
\selectlanguage{english}
In Africa, the holy confessors 
 Eugene, the faithful and virtuous bishop of Carthage, and all the clergy of 
 that Church, to the number of about five hundred or more, among whom were 
 many small children who performed the office of lector. In the 
 persecution of the Vandals, under the Arian king Hunneric, they were 
 subjected to scourging and starvation, and driven into a most painful 
 banishment which they bore with joy for God's sake. In their number 
 were also two distinguished persons, the archdeacon Salutaris, and Muritta, 
 occupying the second rank among the ministers of the Church. Both had 
 three times confessed the faith, and were illustrious by their sturdy 
 perseverance in Christianity.
\switchcolumn*
\selectlanguage{latin}
In Británnia minóre 
 sancti Turiávi, Epíscopi et Confessóris, miræ simplicitátis et innocéntiæ 
 viri.
\switchcolumn
\selectlanguage{english}
In Brittany, St. Turian, bishop and 
 confessor, a man of admirable simplicity and innocence.
\switchcolumn*
\selectlanguage{latin}
\end{paracol}


% ---- martyrology/mart07/mart0714.htm
\needspace{10\baselineskip}
\begin{paracol}{2}
\selectlanguage{latin}
\begin{center}{\color{gregoriocolor} Prídie Idus Júlii. 
 Luna\dots\ }\end{center}
\switchcolumn
\selectlanguage{english}
\begin{center}{\color{gregoriocolor} The 
 14\textsuperscript{th} Day of 
 July. The\dots\ Day of the Moon.}\end{center}
\end{paracol}

\noindent\begin{tabularx}{\linewidth}{*{19}{>{\centering\arraybackslash}X}}
 \textcolor{gregoriocolor}{a} & \textcolor{gregoriocolor}{b} & \textcolor{gregoriocolor}{c} & \textcolor{gregoriocolor}{d} & \textcolor{gregoriocolor}{e} & \textcolor{gregoriocolor}{f} & \textcolor{gregoriocolor}{g} & \textcolor{gregoriocolor}{h} & \textcolor{gregoriocolor}{i} & \textcolor{gregoriocolor}{k} & \textcolor{gregoriocolor}{l} & \textcolor{gregoriocolor}{m} & \textcolor{gregoriocolor}{n} & \textcolor{gregoriocolor}{p} & \textcolor{gregoriocolor}{q} & \textcolor{gregoriocolor}{r} & \textcolor{gregoriocolor}{s} & \textcolor{gregoriocolor}{t} & \textcolor{gregoriocolor}{u} \\
 19 & 20 & 21 & 22 & 23 & 24 & 25 & 26 & 27 & 28 & 29 & 30 & 1 & 2 & 3 & 4 & 5 & 6 & 7 \\
\end{tabularx}
\vspace{0.5\baselineskip}
\noindent\begin{tabularx}{\linewidth}{*{12}{>{\centering\arraybackslash}X}}
 \textcolor{gregoriocolor}{A} & \textcolor{gregoriocolor}{B} & \textcolor{gregoriocolor}{C} & \textcolor{gregoriocolor}{D} & \textcolor{gregoriocolor}{E} & F & \textcolor{gregoriocolor}{F} & \textcolor{gregoriocolor}{G} & \textcolor{gregoriocolor}{H} & \textcolor{gregoriocolor}{M} & \textcolor{gregoriocolor}{N} & \textcolor{gregoriocolor}{P} \\
 8 & 9 & 10 & 11 & 12 & 13 & 13 & 14 & 15 & 16 & 17 & 18 \\
\end{tabularx}

\begin{paracol}{2}
\selectlanguage{latin}
\lettrine[lines=2]{S}{ancti} Bonaventúræ, ex 
 Ordine Minórum, Cardinális et Epíscopi Albanénsis, Confessóris et Ecclésiæ 
 Doctóris; qui sequénti die migrávit ad Dóminum.
\switchcolumn
\selectlanguage{english}
\lettrine[lines=2]{S}{t.} Bonaventure of the Order of 
 Friars Minor, cardinal and bishop of Albano, confessor and doctor of the 
 Church, who passed to the Lord on the day following this.
\switchcolumn*
\selectlanguage{latin}
Romæ natális sancti 
 Camílli de Lellis, Presbyteri et Confessóris, Clericórum Regulárium infírmis 
 ministrántium Institutóris; quem virtútibus et miráculis clarum, Summus 
 Póntifex Benedíctus Décimus quartus in Sanctórum número recénsuit, et Leo 
 Décimus tértius cæléstem hospitálium et infirmórum Patrónum renuntiávit. 
 Ipsíus tamen festívitas quintodécimo Kaléndas Augústi recólitur.
\switchcolumn
\selectlanguage{english}
At Rome, the birthday of St. 
 Camillus de Lellis, priest and confessor, founder of the Clerks Regular for 
 Ministering to the Sick. Pope Benedict XIV numbered him among the 
 saints because of the fame of his miracles and virtues; Pope Leo XIII 
 appointed him heavenly protector of hospitals and of the sick. His 
 feast is observed on the 18th of July.
\switchcolumn*
\selectlanguage{latin}
Item Romæ sancti Justi 
 mílitis, qui, sub Cláudio Tribúno, apparénte sibi divínitus Cruce, crédidit 
 in Christum, et, mox baptizátus, ómnia sua paupéribus erogávit; tentus deínde 
 a Magnétio Præfécto, atque verberári nervis, gálea igníta cóntegi et in 
 rogum immítti jussus ac ne in capíllo quidem læsus, in Dómini confessióne 
 réddidit spíritum.
\switchcolumn
\selectlanguage{english}
Also at Rome, St. Justus, a soldier 
 under the tribune Claudius. When a miraculous cross appeared to him he 
 believed in Christ, was baptized, and gave away his goods to the poor. 
 Afterwards arrested by the prefect Magnetius, he was scourged with rods, had 
 a heated helmet put on his head, and was thrown into the fire, but received 
 no injury, not even to a hair of his head. In the end he yielded up 
 his soul confessing the Lord.
\switchcolumn*
\selectlanguage{latin}
Synópe, in Ponto, 
 sancti Phocæ Mártyris, ejúsdem civitátis Epíscopi, qui, sub Trajáno 
 Imperatóre, cárcerem, víncula, ferrum ignémque pro Christo súperans, 
 evolávit in cælum. Ejus relíquiæ póstea Viénnam, in Gállia, sunt 
 delátæ, et in Basílica sanctórum Apostolórum cónditæ.
\switchcolumn
\selectlanguage{english}
At Sinope in Pontus, the martyr St. 
 Phocas, bishop of the city. Under Emperor Trajan, after having been 
 imprisoned, bound, struck with the sword, and exposed to the fire for 
 Christ, he departed to heaven. His remains were brought to Vienne in 
 France, and deposited in the Church of the Holy Apostles.
\switchcolumn*
\selectlanguage{latin}
Alexandríæ sancti 
 Héraclæ Antístitis, ob cujus celebérrimam opiniónem Africánus 
 historiógraphus mémorat se Alexandriam, ad eum viséndum, properásse.
\switchcolumn
\selectlanguage{english}
At Alexandria, St. Heracles, bishop, 
 whose fame was so great that the historian Africanus testifies that he 
 journeyed to Alexandria to see him.
\switchcolumn*
\selectlanguage{latin}
Carthágine sancti Cyri 
 Epíscopi, in cujus festivitáte sanctus Augustínus de eo sermónem ad pópulum 
 hábuit.
\switchcolumn
\selectlanguage{english}
At Carthage, St. Cyrus, bishop, on 
 whose festival St. Augustine spoke of him to his people.
\switchcolumn*
\selectlanguage{latin}
Novocómi sancti Felícis, 
 qui fuit primus ejúsdem civitátis Epíscopus.
\switchcolumn
\selectlanguage{english}
At Como, St. Felix, first bishop of 
 that city.
\switchcolumn*
\selectlanguage{latin}
Bríxiæ sancti Optatiáni 
 Epíscopi.
\switchcolumn
\selectlanguage{english}
At Brescia, St. Optatian, bishop.
\switchcolumn*
\selectlanguage{latin}
Davéntriæ, in Belgis, 
 sancti Marcellíni, Presbyteri et Confessóris.
\switchcolumn
\selectlanguage{english}
At Deventer in Belgium, St. 
 Marcellinus, priest and confessor.
\switchcolumn*
\selectlanguage{latin}
Limæ, in Perúvia, 
 sancti Francísci Soláni, Sacerdótis ex Ordine Minórum et Confessóris; qui, 
 prædicatióne, signis et virtútibus apud Indos Occidentáles illústris, 
 migrávit ad Dóminum, atque a Benedícto Décimo tértio, Pontífice Máximo, 
 Sanctórum fastis adscríptus est.
\switchcolumn
\selectlanguage{english}
At Lima in Peru, St. Francis Solano, 
 a priest and confessor of the Order of Friars Minor. He passed to the 
 Lord in the West Indies, renowned for his preaching, miracles and virtues. 
 Pope Benedict XIII placed him on the canon of the saints.
\switchcolumn*
\selectlanguage{latin}
\end{paracol}


% ---- martyrology/mart07/mart0715.htm
\needspace{10\baselineskip}
\begin{paracol}{2}
\selectlanguage{latin}
\begin{center}{\color{gregoriocolor} Idibus Júlii. 
 Luna\dots\ }\end{center}
\switchcolumn
\selectlanguage{english}
\begin{center}{\color{gregoriocolor} The 
 15\textsuperscript{th} Day of 
 July. The\dots\ Day of the Moon.}\end{center}
\end{paracol}

\noindent\begin{tabularx}{\linewidth}{*{19}{>{\centering\arraybackslash}X}}
 \textcolor{gregoriocolor}{a} & \textcolor{gregoriocolor}{b} & \textcolor{gregoriocolor}{c} & \textcolor{gregoriocolor}{d} & \textcolor{gregoriocolor}{e} & \textcolor{gregoriocolor}{f} & \textcolor{gregoriocolor}{g} & \textcolor{gregoriocolor}{h} & \textcolor{gregoriocolor}{i} & \textcolor{gregoriocolor}{k} & \textcolor{gregoriocolor}{l} & \textcolor{gregoriocolor}{m} & \textcolor{gregoriocolor}{n} & \textcolor{gregoriocolor}{p} & \textcolor{gregoriocolor}{q} & \textcolor{gregoriocolor}{r} & \textcolor{gregoriocolor}{s} & \textcolor{gregoriocolor}{t} & \textcolor{gregoriocolor}{u} \\
 20 & 21 & 22 & 23 & 24 & 25 & 26 & 27 & 28 & 29 & 30 & 1 & 2 & 3 & 4 & 5 & 6 & 7 & 8 \\
\end{tabularx}
\vspace{0.5\baselineskip}
\noindent\begin{tabularx}{\linewidth}{*{12}{>{\centering\arraybackslash}X}}
 \textcolor{gregoriocolor}{A} & \textcolor{gregoriocolor}{B} & \textcolor{gregoriocolor}{C} & \textcolor{gregoriocolor}{D} & \textcolor{gregoriocolor}{E} & F & \textcolor{gregoriocolor}{F} & \textcolor{gregoriocolor}{G} & \textcolor{gregoriocolor}{H} & \textcolor{gregoriocolor}{M} & \textcolor{gregoriocolor}{N} & \textcolor{gregoriocolor}{P} \\
 9 & 10 & 11 & 12 & 13 & 14 & 14 & 15 & 16 & 17 & 18 & 19 \\
\end{tabularx}

\begin{paracol}{2}
\selectlanguage{latin}
\lettrine[lines=2]{S}{ancti} Henríci Primi, 
 Imperatóris Romanórum et Confessóris, cujus dies natális tértio Idus mensis 
 hujus recensétur.
\switchcolumn
\selectlanguage{english}
\lettrine[lines=2]{S}{t.} Henry I, Roman emperor and 
 confessor, whose birthday was noted on the 13th of this month.
\switchcolumn*
\selectlanguage{latin}
Lugdúni, in Gállia, 
 deposítio sancti Bonaventúræ, Cardinális et Epíscopi Albanénsis, Confessóris 
 et Ecclésiæ Doctóris, ex Ordine Minórum, doctrína et vitæ sanctitáte 
 celebérrimi. Ipsíus tamen festívitas prídie hujus diéi recólitur.
\switchcolumn
\selectlanguage{english}
At Lyons in France, the death of St. 
 Bonaventure, cardinal and bishop of Albano, confessor and doctor of the 
 Church, of the Order of Friars Minor, who is famed for his learning and the 
 sanctity of his life. His feast is celebrated on the previous day.
\switchcolumn*
\selectlanguage{latin}
Papíæ sancti Felícis, 
 Epíscopi et Mártyris.
\switchcolumn
\selectlanguage{english}
At Pavia, St. Felix, bishop and 
 martyr.
\switchcolumn*
\selectlanguage{latin}
In Portu Románo natális 
 sanctórum Mártyrum Eutrópii, atque Zósimæ et Bonósæ sorórum.
\switchcolumn
\selectlanguage{english}
At Porto, the birthday of the holy 
 martyrs Eutropius, and the sisters Zosima and Bonosa.
\switchcolumn*
\selectlanguage{latin}
Carthágine beáti 
 Catulíni Diáconi, de cujus láudibus sanctus Augustínus sermónem ad pópulum 
 hábuit, et sanctórum Januárii, Floréntii, Júliæ et Justæ Mártyrum; qui 
 pósiti sunt in Basílica Fausti.
\switchcolumn
\selectlanguage{english}
At Carthage, blessed Catulinus, 
 deacon, whose glories were proclaimed by St. Augustine in a sermon to his 
 people. Also the saints Januarius, Florentius, Julia, and Justa, 
 martyrs, who were entombed in the Church of St. Faustus.
\switchcolumn*
\selectlanguage{latin}
Alexandríæ sanctórum 
 Mártyrum Philíppi, Zenónis, Narséi et decem infántum.
\switchcolumn
\selectlanguage{english}
At Alexandria, the holy martyrs 
 Philip, Zeno, Narseus, and ten children.
\switchcolumn*
\selectlanguage{latin}
In ínsula Ténedo sancti 
 Abudémii Mártyris, qui sub Diocletiáno passus est.
\switchcolumn
\selectlanguage{english}
In the island of Tenedos, St. 
 Abudemius, martyr, who suffered under Diocletian.
\switchcolumn*
\selectlanguage{latin}
Sebáste, in Arménia, sancti Antíochi médici, qui, sub Hadriáno Præside, cápite obtruncátus est; cumque ex eo lac pro sánguine manáret, Cyríacus cárnifex, convérsus ad 
 Christum, et ipse martyrium súbiit.
\switchcolumn
\selectlanguage{english}
At Sebaste in Armenia, St. 
 Antiochus, a physician, who was beheaded under the governor Adrian. On 
 seeing milk flowing from his wounds instead of blood, Cyriacus, his 
 executioner, was converted to Christ and endured martyrdom.
\switchcolumn*
\selectlanguage{latin}
Nísibi, in Mesopotámia, natális sancti Jacóbi, ejúsdem urbis Epíscopi, magnæ sanctitátis viri. 
 Hic, miráculis et eruditióne clarus, unus fuit, sub persecutióne Galérii 
 Maximiáni, ex número Confessórum, qui in Nicæna deínde Synodo perversitátem 
 Aríi, Homoúsii oppositióne, damnárunt; cujus et sancti Alexándri Epíscopi 
 oratióne ipse Aríus condígnam suæ iniquitátis mercédem, effúsis viscéribus, 
 Constantinópoli recépit.
\switchcolumn
\selectlanguage{english}
At Nisibis in Mesopotamia, the 
 birthday of St. James, bishop of that city, a man celebrated for great 
 holiness, miracles and learning. He was one of those who confessed the 
 faith during the persecution of Galerius Maximian, and later condemned the 
 perverse heresy of Arius in the Council of Nicaea by opposing to the 
 doctrine of consubstantiality. It was also owing to his prayers, and 
 those of the bishop Alexander, that Arius received at Constantinople the 
 suitable punishment of his iniquity, his bowels gushing out.
\switchcolumn*
\selectlanguage{latin}
Neápoli, in Campánia, 
 sancti Athanásii, ejúsdem civitátis Epíscopi, qui, ab ímpio nepóte Sérgio 
 multa passus ac sede pulsus, tandem Vérulis, in Hérnicis, conféctus ærúmnis, 
 migrávit in cælum, témpore Cároli Calvi.
\switchcolumn
\selectlanguage{english}
At Naples in Campania, St. 
 Athanasius, bishop of that city, who suffered a great deal from his wicked 
 nephew Sergius, by whom he was driven from his diocese. Overcome with 
 afflictions, he departed for heaven at Veroli, in the time of Charles the 
 Bald.
\switchcolumn*
\selectlanguage{latin}
Campi Salentinórum, in 
 Apúlia, sancti Pompílii Maríæ Pirrótti, Confessóris, ex Ordine Clericórum 
 Páuperum Matris Dei Scholárum Piárum, vita apostólica præclári, a Pio 
 Undécimo, Pontífice Máximo, inter Sanctos reláti.
\switchcolumn
\selectlanguage{english}
At Campo in Italy, the birthday of 
 St. Pompilio Maria Pirrotti of St. Nicholas, confessor, a member of the 
 Congregation of Poor Clerks Regular of Pious Schools of the Mother of God, 
 who spent his entire life in safeguarding the salvation of souls. He 
 was registered among the saints by Pope Pius XI.
\switchcolumn*
\selectlanguage{latin}
Panórmi Invéntio 
 córporis sanctæ Rosáliæ, Vírginis Panormitánæ; quod, Urbáno Octávo Pontífice 
 Máximo, repértum divínitus, Jubilæi anno Sicíliam a peste liberávit.
\switchcolumn
\selectlanguage{english}
At Palermo, the finding of the body 
 of St. Rosalia, virgin of that city. Miraculously discovered in the 
 time of Pope Urban VIII, it delivered Sicily from the plague in the year of 
 the Jubilee.
\switchcolumn*
\selectlanguage{latin}
\end{paracol}


% ---- martyrology/mart07/mart0716.htm
\needspace{10\baselineskip}
\begin{paracol}{2}
\selectlanguage{latin}
\begin{center}{\color{gregoriocolor} Décimo séptimo Kaléndas Augústi. 
 Luna\dots\ }\end{center}
\switchcolumn
\selectlanguage{english}
\begin{center}{\color{gregoriocolor} The 
 16\textsuperscript{th} Day of 
 July. The\dots\ Day of the Moon.}\end{center}
\end{paracol}

\noindent\begin{tabularx}{\linewidth}{*{19}{>{\centering\arraybackslash}X}}
 \textcolor{gregoriocolor}{a} & \textcolor{gregoriocolor}{b} & \textcolor{gregoriocolor}{c} & \textcolor{gregoriocolor}{d} & \textcolor{gregoriocolor}{e} & \textcolor{gregoriocolor}{f} & \textcolor{gregoriocolor}{g} & \textcolor{gregoriocolor}{h} & \textcolor{gregoriocolor}{i} & \textcolor{gregoriocolor}{k} & \textcolor{gregoriocolor}{l} & \textcolor{gregoriocolor}{m} & \textcolor{gregoriocolor}{n} & \textcolor{gregoriocolor}{p} & \textcolor{gregoriocolor}{q} & \textcolor{gregoriocolor}{r} & \textcolor{gregoriocolor}{s} & \textcolor{gregoriocolor}{t} & \textcolor{gregoriocolor}{u} \\
 21 & 22 & 23 & 24 & 25 & 26 & 27 & 28 & 29 & 30 & 1 & 2 & 3 & 4 & 5 & 6 & 7 & 8 & 9 \\
\end{tabularx}
\vspace{0.5\baselineskip}
\noindent\begin{tabularx}{\linewidth}{*{12}{>{\centering\arraybackslash}X}}
 \textcolor{gregoriocolor}{A} & \textcolor{gregoriocolor}{B} & \textcolor{gregoriocolor}{C} & \textcolor{gregoriocolor}{D} & \textcolor{gregoriocolor}{E} & F & \textcolor{gregoriocolor}{F} & \textcolor{gregoriocolor}{G} & \textcolor{gregoriocolor}{H} & \textcolor{gregoriocolor}{M} & \textcolor{gregoriocolor}{N} & \textcolor{gregoriocolor}{P} \\
 10 & 11 & 12 & 13 & 14 & 15 & 15 & 16 & 17 & 18 & 19 & 20 \\
\end{tabularx}

\begin{paracol}{2}
\selectlanguage{latin}
\lettrine[lines=2]{F}{estum} beátæ Maríæ 
 Vírginis de monte Carmélo.
\switchcolumn
\selectlanguage{english}
\lettrine[lines=2]{T}{he} feast of the Blessed Virgin Mary 
 of Mount Carmel.
\switchcolumn*
\selectlanguage{latin}
Sebáste in Arménia, sanctórum Mártyrum Athenógenis Epíscopi, et decem discipulórum ejus, sub 
 Diocletiáno Imperatóre.
\switchcolumn
\selectlanguage{english}
At Sebaste in Armenia, the holy 
 martyrs Athenogenes, bishop, and ten of his disciples, in the time of 
 Emperor Diocletian.
\switchcolumn*
\selectlanguage{latin}
Tréviris sancti 
 Valentíni, Epíscopi et Mártyris.
\switchcolumn
\selectlanguage{english}
At Treves, St. Valentine, bishop and 
 martyr.
\switchcolumn*
\selectlanguage{latin}
Córdubæ, in Hispánia, 
 sancti Sisenándi, Levítæ et Mártyris; qui a Saracénis, pro Christi fide, 
 jugulátus est.
\switchcolumn
\selectlanguage{english}
At Cordova in Spain, St. Sisenand, 
 cleric and martyr, who was strangled by the Saracens for the faith of 
 Christ.
\switchcolumn*
\selectlanguage{latin}
Eódem die natális 
 sancti Fausti Mártyris, qui, sub Décio Imperatóre, cruci affíxus, quinque 
 diébus in ea vixit, ac demum, sagíttis confóssus, migrávit in cælum.
\switchcolumn
\selectlanguage{english}
The same day, the birthday of St. 
 Faustus, martyr, under Decius. He lived five days fastened to a cross, 
 and being then pierced with arrows, he went to heaven.
\switchcolumn*
\selectlanguage{latin}
Sántis, in Gállia, 
 sanctórum Mártyrum Rainéldis Vírginis, et Sociórum ejus, qui a bárbaris, pro 
 Christe fide, cæsi sunt.
\switchcolumn
\selectlanguage{english}
At Saintes in France, the holy 
 martyrs Raineld, virgin, and her companions who were slain by barbarians for 
 the Christian faith.
\switchcolumn*
\selectlanguage{latin}
Bérgomi sancti 
 Domniónis Mártyris.
\switchcolumn
\selectlanguage{english}
At Bergamo, St. Domnio, martyr.
\switchcolumn*
\selectlanguage{latin}
Antiochíæ natális beáti 
 Eustáchii, Epíscopi et Confessóris, doctrína et sanctitáte célebris; qui, 
 sub Ariáno Imperatóre Constántio, ob cathólicæ fídei defensiónem, exsílio 
 Trajanópolim Thráciæ pulsus, ibídem in Dómino requiévit.
\switchcolumn
\selectlanguage{english}
At Antioch, the birthday of blessed 
 Eustace, bishop and confessor, celebrated for learning and sanctity. 
 Under the Arian emperor Constantius, for the defence of the Catholic faith, 
 he was banished to Trajanopolis in Thrace, where he rested in the Lord.
\switchcolumn*
\selectlanguage{latin}
Cápuæ sancti Vitaliáni, 
 Epíscopi et Confessóris.
\switchcolumn
\selectlanguage{english}
At Capua, St. Vitalian, bishop and 
 confessor.
\switchcolumn*
\selectlanguage{latin}
Apud Abbatíam 
 Sanctíssimi Salvatóris, diœcésis Constantiénsis, in Gállia, sanctæ 
 Maríæ-Magdalénæ Postel, Fundatrícis Institúti Sorórum Scholárum 
 Christianárum a Misericórdia, a Pio Papa Undécimo in sanctárum Vírginum 
 album relátæ.
\switchcolumn
\selectlanguage{english}
At the abbey of our Most Holy 
 Redeemer, in the diocese of Coutances in France, St. Mary Magdalene Postel, 
 foundress of the Sisters of Mercy of the Christian Schools, who was added to 
 the list of the holy virgins by Pope Pius XI.
\switchcolumn*
\selectlanguage{latin}
Apud Ostia Tiberína 
 Translátio córporis sancti Hilaríni Mónachi, qui, una cum sancto Donáto, in 
 persecutióne Juliáni Apóstatæ, comprehénsus est, et, cum nollet sacrificáre, 
 demum, fústibus cæsus, Arétii, in Túscia, martyrium sumpsit séptimo Idus 
 Augústi.
\switchcolumn
\selectlanguage{english}
The translation of St. Hilarinus, a 
 monk, to Ostia. He was arrested with St. Donatus in the persecution of 
 Julian. Because he refused to sacrifice to idols, he was finally 
 scourged at Arezzo in Tuscany, and underwent martyrdom on the 7th of August.
\switchcolumn*
\selectlanguage{latin}
\end{paracol}


% ---- martyrology/mart07/mart0717.htm
\needspace{10\baselineskip}
\begin{paracol}{2}
\selectlanguage{latin}
\begin{center}{\color{gregoriocolor} Sextodécimo Kaléndas Augústi. 
 Luna\dots\ }\end{center}
\switchcolumn
\selectlanguage{english}
\begin{center}{\color{gregoriocolor} The 
 17\textsuperscript{th} Day of 
 July. The\dots\ Day of the Moon.}\end{center}
\end{paracol}

\noindent\begin{tabularx}{\linewidth}{*{19}{>{\centering\arraybackslash}X}}
 \textcolor{gregoriocolor}{a} & \textcolor{gregoriocolor}{b} & \textcolor{gregoriocolor}{c} & \textcolor{gregoriocolor}{d} & \textcolor{gregoriocolor}{e} & \textcolor{gregoriocolor}{f} & \textcolor{gregoriocolor}{g} & \textcolor{gregoriocolor}{h} & \textcolor{gregoriocolor}{i} & \textcolor{gregoriocolor}{k} & \textcolor{gregoriocolor}{l} & \textcolor{gregoriocolor}{m} & \textcolor{gregoriocolor}{n} & \textcolor{gregoriocolor}{p} & \textcolor{gregoriocolor}{q} & \textcolor{gregoriocolor}{r} & \textcolor{gregoriocolor}{s} & \textcolor{gregoriocolor}{t} & \textcolor{gregoriocolor}{u} \\
 22 & 23 & 24 & 25 & 26 & 27 & 28 & 29 & 30 & 1 & 2 & 3 & 4 & 5 & 6 & 7 & 8 & 9 & 10 \\
\end{tabularx}
\vspace{0.5\baselineskip}
\noindent\begin{tabularx}{\linewidth}{*{12}{>{\centering\arraybackslash}X}}
 \textcolor{gregoriocolor}{A} & \textcolor{gregoriocolor}{B} & \textcolor{gregoriocolor}{C} & \textcolor{gregoriocolor}{D} & \textcolor{gregoriocolor}{E} & F & \textcolor{gregoriocolor}{F} & \textcolor{gregoriocolor}{G} & \textcolor{gregoriocolor}{H} & \textcolor{gregoriocolor}{M} & \textcolor{gregoriocolor}{N} & \textcolor{gregoriocolor}{P} \\
 11 & 12 & 13 & 14 & 15 & 16 & 16 & 17 & 18 & 19 & 20 & 21 \\
\end{tabularx}

\begin{paracol}{2}
\selectlanguage{latin}
\lettrine[lines=2]{R}{omæ} sancti Aléxii 
 Confessóris, ex Euphemiáno Senatóre progéniti. Hic, prima nocte 
 nuptiárum, sponsa intácta, e domo sua abscédens, ac, post longam 
 peregrinatiónem, ad Urbem rédiens, decem et septem annos tamquam egénus in 
 domo patérna recéptus hospítio, nova mundum arte delúdens, incógnitus mansit; 
 sed post óbitum, et voce per Urbis Ecclésias audíta et scripto suo ágnitus, 
 Innocéntio Primo Pontífice Máximo, ad sancti Bonifátii Ecclésiam summo 
 honóre delátus est, ibíque multis miráculis cláruit.
\switchcolumn
\selectlanguage{english}
\lettrine[lines=2]{A}{t} Rome, St. Alexius, confessor, son 
 of the senator Euphemian. Leaving his spouse before the night of 
 marriage, he withdrew from his house, and after a long pilgrimage, returned 
 to Rome where he was for seventeen years harboured in his father's house as 
 an unknown pauper, thus deluding the world in this strange way. After 
 his death, however, becoming known through a voice heard in the churches of 
 the city, and by his own writings, he was, under the sovereign Pontiff 
 Innocent I, translated to the Church of St. Boniface, where he wrought many 
 miracles.
\switchcolumn*
\selectlanguage{latin}
Carthágine natális 
 sanctórum Mártyrum Scillitanórum, id est Speráti, Narzális, Cythíni, Vetúrii, 
 Felícis, Acyllíni, Lætántii, Januáriæ, Generósæ, Vestínæ, Donátæ et Secúndæ; 
 qui, jussu Saturníni Præfécti, post primam Christi confessiónem, in cárcerem 
 trusi et in ligno confíxi, deínde gládio decolláti sunt. Speráti autem 
 relíquiæ, cum óssibus beáti Cypriáni et cápite sancti Pantaleónis Mártyris, 
 ex Africa in Gállias translátæ, Lugdúni, in Basílica sancti Joánnis Baptístæ, 
 religióse cónditæ fuérunt.
\switchcolumn
\selectlanguage{english}
At Carthage, the birthday of the 
 holy Scillitan martyrs Speratus, Narzales, Cythinus, Venturius, Felix, 
 Acyllinus, Laetantius, Januaria, Generosa, Vestina, Donata, and Secunda. 
 By order of the prefect Saturninus, after their first confession of the 
 faith, they were sent to prison, nailed to a cross, and finally beheaded. 
 The relics of Speratus, with the bones of blessed Cyprian and the head of 
 the martyr, St. Pantaleon, were carried from Africa into France and 
 honourably buried in the basilica of St. John the Baptist at Lyons.
\switchcolumn*
\selectlanguage{latin}
Amástride, in 
 Paphlagónia, sancti Hyacínthi Mártyris, qui, sub Castrítio Præside, multa 
 passus, quiévit in cárcere.
\switchcolumn
\selectlanguage{english}
At Amastris in Paphlagonia, St. 
 Hyacinth, martyr, who died in prison after much suffering, under the prefect 
 Castritus.
\switchcolumn*
\selectlanguage{latin}
Tíbure sancti Generósi 
 Mártyris.
\switchcolumn
\selectlanguage{english}
At Tivoli, St. Generosus, martyr.
\switchcolumn*
\selectlanguage{latin}
Constantinópoli sanctæ 
 Theódotæ Mártyris, sub Leóne Iconoclásta.
\switchcolumn
\selectlanguage{english}
At Constantinople, St. Theodota, 
 martyr, under Leo the Iconoclast.
\switchcolumn*
\selectlanguage{latin}
Romæ deposítio sancti 
 Leónis Papæ Quarti.
\switchcolumn
\selectlanguage{english}
At Rome, the death of Pope St. Leo 
 IV.
\switchcolumn*
\selectlanguage{latin}
Papíæ sancti Ennódii, 
 Epíscopi et Confessóris.
\switchcolumn
\selectlanguage{english}
At Pavia, St. Ennodius, bishop and 
 confessor.
\switchcolumn*
\selectlanguage{latin}
Antisiodóri sancti 
 Theodósii Epíscopi.
\switchcolumn
\selectlanguage{english}
At Auxerre, St. Theodosius, bishop.
\switchcolumn*
\selectlanguage{latin}
Medioláni sanctæ 
 Marcellínæ Vírginis, soróris beáti Ambrósii Epíscopi, quæ Romæ, in Basílica 
 sancti Petri, a Libério Papa velum consecratiónis accépit; cujus étiam 
 sanctitátem idem beátus Ambrósius scriptis suis testátam facit.
\switchcolumn
\selectlanguage{english}
At Milan, the virgin saint 
 Marcellina, sister of the blessed bishop Ambrose, who received the religious 
 veil from Pope Liberius, in the basilica of St. Peter at Rome. Her 
 sanctity is attested to by St. Ambrose in his writings.
\switchcolumn*
\selectlanguage{latin}
Venétiis Translátio 
 sanctæ Marínæ Vírginis.
\switchcolumn
\selectlanguage{english}
At Venice, the translation of St. 
 Marina, virgin.
\switchcolumn*
\selectlanguage{latin}
\end{paracol}


% ---- martyrology/mart07/mart0718.htm
\needspace{10\baselineskip}
\begin{paracol}{2}
\selectlanguage{latin}
\begin{center}{\color{gregoriocolor} Quintodécimo Kaléndas Augústi. 
 Luna\dots\ }\end{center}
\switchcolumn
\selectlanguage{english}
\begin{center}{\color{gregoriocolor} The 
 18\textsuperscript{th} Day of 
 July. The\dots\ Day of the Moon.}\end{center}
\end{paracol}

\noindent\begin{tabularx}{\linewidth}{*{19}{>{\centering\arraybackslash}X}}
 \textcolor{gregoriocolor}{a} & \textcolor{gregoriocolor}{b} & \textcolor{gregoriocolor}{c} & \textcolor{gregoriocolor}{d} & \textcolor{gregoriocolor}{e} & \textcolor{gregoriocolor}{f} & \textcolor{gregoriocolor}{g} & \textcolor{gregoriocolor}{h} & \textcolor{gregoriocolor}{i} & \textcolor{gregoriocolor}{k} & \textcolor{gregoriocolor}{l} & \textcolor{gregoriocolor}{m} & \textcolor{gregoriocolor}{n} & \textcolor{gregoriocolor}{p} & \textcolor{gregoriocolor}{q} & \textcolor{gregoriocolor}{r} & \textcolor{gregoriocolor}{s} & \textcolor{gregoriocolor}{t} & \textcolor{gregoriocolor}{u} \\
 23 & 24 & 25 & 26 & 27 & 28 & 29 & 30 & 1 & 2 & 3 & 4 & 5 & 6 & 7 & 8 & 9 & 10 & 11 \\
\end{tabularx}
\vspace{0.5\baselineskip}
\noindent\begin{tabularx}{\linewidth}{*{12}{>{\centering\arraybackslash}X}}
 \textcolor{gregoriocolor}{A} & \textcolor{gregoriocolor}{B} & \textcolor{gregoriocolor}{C} & \textcolor{gregoriocolor}{D} & \textcolor{gregoriocolor}{E} & F & \textcolor{gregoriocolor}{F} & \textcolor{gregoriocolor}{G} & \textcolor{gregoriocolor}{H} & \textcolor{gregoriocolor}{M} & \textcolor{gregoriocolor}{N} & \textcolor{gregoriocolor}{P} \\
 12 & 13 & 14 & 15 & 16 & 17 & 17 & 18 & 19 & 20 & 21 & 22 \\
\end{tabularx}

\begin{paracol}{2}
\selectlanguage{latin}
\lettrine[lines=2]{S}{ancti} Camílli de 
 Lellis, Presbyteri et Confessóris, Clericórum Regulárium infírmis 
 ministrántium Institutóris, cæléstis hospitálium et infirmórum Patróni; 
 cujus dies natális prídie Idus Júlii recólitur.
\switchcolumn
\selectlanguage{english}
\lettrine[lines=2]{S}{t.} Camillus de Lellis, priest and 
 confessor, founder of the Clerks Regular Ministering to the Sick, the 
 heavenly patron of hospitals and of the sick, whose birthday is the 14th day 
 of July.
\switchcolumn*
\selectlanguage{latin}
Tíbure sanctæ 
 Symphorósæ, uxóris sancti Getúlii Mártyris, cum septem fíliis suis, scílicet 
 Crescénte, Juliáno, Nemésio, Primitívo, Justíno, Stácteo et Eugénio. 
 Horum mater, sub Hadriáno Príncipe, ob insuperábilem constántiam, primo cæsa 
 diu palmis, deínde crínibus suspénsa, novíssime saxo alligáta, in flumen 
 præcipitáta est; fílii autem, stipítibus ad tróchleas exténsi, divérso 
 mortis éxitu martyrium complevérunt. Eorúndem córpora póstea 
 Romam transláta, et, Pio Quarto Summo Pontífice, in Diacónia sancti Angeli 
 in Piscína fuérunt invénta.
\switchcolumn
\selectlanguage{english}
At Tivoli, in the time of Emperor 
 Hadrian, St. Symphorosa, wife of the martyr St. Getulius, with her seven 
 sons, Crescens, Julian, Nemesius, Primitivus, Justin, Stacteus, and Eugene. 
 The mother, because of her invincible constancy, was first beaten a long 
 time, then suspended by her hair, and lastly thrown into the river with a 
 stone tied to her body. Her sons were stretched by pulleys attached to 
 stakes, and completed their martyrdom in divers ways. Afterwards, 
 their bodies were taken to Rome, and in the pontificate of Pius IV, were 
 found in the sacristy of St. Angelo in Piscina.
\switchcolumn*
\selectlanguage{latin}
Trajécti sancti 
 Frideríci, Epíscopi et Mártyris.
\switchcolumn
\selectlanguage{english}
At Utrecht, St. Frederick, bishop 
 and martyr.
\switchcolumn*
\selectlanguage{latin}
Doróstori, in Mysia 
 inferióre, sancti Æmiliáni Mártyris, qui, témpore Juliáni Apóstatæ, sub 
 Capitolíno Præside, in fornácem injéctus, martyrii palmam accépit.
\switchcolumn
\selectlanguage{english}
At Silistria in Bulgaria, St. 
 Emilian, martyr, who was cast into a furnace, in the time of Julian the 
 Apostate, under the governor Capitolinus, and received the palm of 
 martyrdom.
\switchcolumn*
\selectlanguage{latin}
Carthágine sanctæ 
 Gundénis Vírginis, quæ, jussu Rufíni Procónsulis, ob Christi confessiónem, 
 quater divérsis tempóribus per extensiónem in equúleo torta, et úngulis 
 horrénde lacerántibus cruciáta, cárceris squalóre diu afflícta, novíssime 
 gládio cæsa est.
\switchcolumn
\selectlanguage{english}
At Carthage, St. Gundenes, virgin. 
 By order of the proconsul Ruffinus, she was at four different times 
 stretched on the rack for the faith of Christ, horribly lacerated with iron 
 hooks, confined for a long time in a filthy prison, and finally put to the 
 sword.
\switchcolumn*
\selectlanguage{latin}
Gallǽciæ, in Hispánia, 
 sanctæ Marínæ, Vírginis et Mártyris.
\switchcolumn
\selectlanguage{english}
In Spanish Galicia, St. Marina, 
 virgin and martyr.
\switchcolumn*
\selectlanguage{latin}
Medioláni sancti 
 Matérni Epíscopi, qui, sub Maximiáni Imperatóre, pro fide Christi et pro 
 Ecclésia sibi commíssa, in cárcerem detrúsus et sæpe verbéribus cæsus est; 
 ac tandem, multis confessiónibus clarus, obdormívit in Dómino.
\switchcolumn
\selectlanguage{english}
At Milan, in the reign of Maximian, 
 the holy bishop Maternus. For the faith of Christ and the Church 
 entrusted to him, he went to his rest in the Lord with a great renown for 
 his repeated confession of the faith.
\switchcolumn*
\selectlanguage{latin}
Bríxiæ natális sancti 
 Philástrii, qui fuit ejúsdem civitátis Epíscopus. Hic advérsus 
 hæréticos, præsértim Ariános, a quibus multa passus est, plúrimum verbis 
 scriptísque pugnávit; demum, clarus miráculis, Conféssor in pace quiévit.
\switchcolumn
\selectlanguage{english}
At Brescia, the birthday of St. 
 Philastrius, bishop of that city, who both by word and writing opposed the 
 heretics, especially the Arians, from whom he suffered greatly. 
 Finally he died in peace, a confessor renowned for miracles.
\switchcolumn*
\selectlanguage{latin}
Metis, in Gállia, 
 sancti Arnúlfi Epíscopi, qui, sanctitáte et miráculis illústris, eremíticam 
 delégit vitam, et beáto fine quiévit.
\switchcolumn
\selectlanguage{english}
At Metz in France, St. Arnulf, a 
 bishop illustrious for holiness and miracles. He chose the life of a 
 hermit and ended his blessed career in peace.
\switchcolumn*
\selectlanguage{latin}
Sígniæ sancti Brunónis, 
 Epíscopi et Confessóris.
\switchcolumn
\selectlanguage{english}
At Segni, St. Bruno, bishop and 
 confessor.
\switchcolumn*
\selectlanguage{latin}
Foro Pompílii, in 
 Æmília, sancti Ruffilli, ejúsdem civitátis Epíscopi.
\switchcolumn
\selectlanguage{english}
At Forlimpopoli in Emilia, St. 
 Ruffillus, bishop of that city.
\switchcolumn*
\selectlanguage{latin}
\end{paracol}


% ---- martyrology/mart07/mart0719.htm
\needspace{10\baselineskip}
\begin{paracol}{2}
\selectlanguage{latin}
\begin{center}{\color{gregoriocolor} Quartodécimo Kaléndas Augústi. Luna\dots\ }\end{center}
\switchcolumn
\selectlanguage{english}
\begin{center}{\color{gregoriocolor} The 
 19\textsuperscript{th} Day of 
 July. The\dots\ Day of the Moon.}\end{center}
\end{paracol}

\noindent\begin{tabularx}{\linewidth}{*{19}{>{\centering\arraybackslash}X}}
 \textcolor{gregoriocolor}{a} & \textcolor{gregoriocolor}{b} & \textcolor{gregoriocolor}{c} & \textcolor{gregoriocolor}{d} & \textcolor{gregoriocolor}{e} & \textcolor{gregoriocolor}{f} & \textcolor{gregoriocolor}{g} & \textcolor{gregoriocolor}{h} & \textcolor{gregoriocolor}{i} & \textcolor{gregoriocolor}{k} & \textcolor{gregoriocolor}{l} & \textcolor{gregoriocolor}{m} & \textcolor{gregoriocolor}{n} & \textcolor{gregoriocolor}{p} & \textcolor{gregoriocolor}{q} & \textcolor{gregoriocolor}{r} & \textcolor{gregoriocolor}{s} & \textcolor{gregoriocolor}{t} & \textcolor{gregoriocolor}{u} \\
 24 & 25 & 26 & 27 & 28 & 29 & 30 & 1 & 2 & 3 & 4 & 5 & 6 & 7 & 8 & 9 & 10 & 11 & 12 \\
\end{tabularx}
\vspace{0.5\baselineskip}
\noindent\begin{tabularx}{\linewidth}{*{12}{>{\centering\arraybackslash}X}}
 \textcolor{gregoriocolor}{A} & \textcolor{gregoriocolor}{B} & \textcolor{gregoriocolor}{C} & \textcolor{gregoriocolor}{D} & \textcolor{gregoriocolor}{E} & F & \textcolor{gregoriocolor}{F} & \textcolor{gregoriocolor}{G} & \textcolor{gregoriocolor}{H} & \textcolor{gregoriocolor}{M} & \textcolor{gregoriocolor}{N} & \textcolor{gregoriocolor}{P} \\
 13 & 14 & 15 & 16 & 17 & 18 & 18 & 19 & 20 & 21 & 22 & 23 \\
\end{tabularx}

\begin{paracol}{2}
\selectlanguage{latin}
\lettrine[lines=2]{S}{ancti} Vincénti a 
 Paulo, Presbyteri et Confessóris, Congregatiónis Presbyterórum Missiónis et 
 Puellárum Caritátis Fundatóris, cæléstis ómnium caritátis Societátum Patróni; 
 qui in Dómino obdormívit quinto Kaléndas Octóbris.
\switchcolumn
\selectlanguage{english}
\lettrine[lines=2]{S}{t.} Vincent de Paul, priest and 
 confessor, founder of the priests of the Congregation of the Mission and the 
 Sisters of Charity, the heavenly patron of all charitable organizations. 
 He fell asleep in the Lord on the 27th of September.
\switchcolumn*
\selectlanguage{latin}
Colóssis, in Phrygia, 
 natális sancti Epaphræ, quem sanctus Paulus Apóstolus concaptívum appéllat. 
 Hic, ab eódem Apóstolo Colóssis Epíscopus ordinátus, ibídem, clarus 
 virtútibus, martyrii palmam, pro óvibus sibi commendátis, viríli agóne 
 percépit; cujus corpus Romæ, in Basílica sanctæ Maríæ Majóris, cónditum est.
\switchcolumn
\selectlanguage{english}
At Colossae in Phrygia, the birthday 
 of St. Epaphras, whom the apostle St. Paul calls his fellow-prisoner. 
 By the same apostle he was consecrated bishop of Colossae, where, becoming 
 renowned for his virtues, he received the palm of martyrdom for defending 
 courageously the flock committed to his charge. His body lies at Rome 
 in the basilica of St. Mary Major.
\switchcolumn*
\selectlanguage{latin}
Tréviris sancti 
 Martíni, Epíscopi et Mártyris.
\switchcolumn
\selectlanguage{english}
At Treves, St. Martin, bishop and 
 martyr.
\switchcolumn*
\selectlanguage{latin}
Híspali, in Hispánia, 
 pássio sanctárum Vírginum Justæ et Rufínæ, quæ, a Præside Diogeniáno 
 comprehénsæ, primo equúlei extensióne et ungulárum laniatióne vexátæ, póstea cárcere, inédia et váriis torsiónibus sunt afflíctæ; tandem Justa in cárcere 
 spíritum exhalávit, Rufínæ vero cervix, in confessióne Dómini, confracta 
 est.
\switchcolumn
\selectlanguage{english}
At Seville in Spain, the martyrdom 
 of the holy virgins Justa and Rufina. Arrested by the governor 
 Diogenian, they were stretched on the rack and lacerated with iron claws, 
 then imprisoned and subjected to starvation and various tortures. 
 Justa died in prison, but Rufina's neck was broken for the confession of the 
 Lord.
\switchcolumn*
\selectlanguage{latin}
Córdubæ, in Hispánia, 
 sanctæ Aureæ Vírginis, beatórum Adúlfi et Joánnis Mártyrum soróris; quæ 
 aliquándo in apostasíæ crimen a Mahumetáno Júdice indúcta est, sed mox, 
 facti pænitens, iteráto certámine, hostem effúso sánguine superávit.
\switchcolumn
\selectlanguage{english}
At Cordova in Spain, St. Aura, 
 virgin, the sister of the holy martyrs Adulphus and John. A Mohammedan 
 judge had persuaded her to apostatize for a while, but quickly repenting of 
 what she had done, in the second trial overcame the enemy by the shedding of 
 her blood.
\switchcolumn*
\selectlanguage{latin}
Romæ sancti Symmachi 
 Papæ, qui, a schismaticórum factióne diútius fatigátus, demum, sanctitáte 
 conspícuus, migrávit ad Dóminum.
\switchcolumn
\selectlanguage{english}
At Rome, Pope St. Symmachus, who for 
 a long time had much to bear, from a faction of schismatics. At last, 
 distinguished by holiness, he went to God.
\switchcolumn*
\selectlanguage{latin}
Verónæ sancti Felícis 
 Epíscopi.
\switchcolumn
\selectlanguage{english}
At Verona, St. Felix, bishop.
\switchcolumn*
\selectlanguage{latin}
Apud Scetim, Ægypti 
 montem, sancti Arsénii, Románæ Ecclésiæ Diáconi; qui, Theodósii témpore, in 
 solitúdinem secéssit, ibíque, virtútibus ómnibus consummátus et jugi 
 lacrimárum imbre perfúsus, spíritum Deo réddidit.
\switchcolumn
\selectlanguage{english}
At Scete, a mountain in Egypt, St. 
 Arsenius, a deacon of the Roman Church. In the time of Theodosius he 
 retired into a desert where, endowed with every virtue and shedding 
 continual tears, he yielded his soul unto God.
\switchcolumn*
\selectlanguage{latin}
In Cappadócia sanctæ 
 Macrínæ Vírginis, fíliæ sanctórum Basilíi et Emméliæ, atque soróris item 
 sanctórum Episcopórum Basilíi Magni, Gregórii Nysséni et Petri Sebastiénsis.
\switchcolumn
\selectlanguage{english}
In Cappadocia, St. Macrina, virgin. 
 She was the daughter of Saints Basil and Emmelia, and the sister of the holy 
 bishops, St. Basil the Great, St. Gregory of Nyssa, and St. Peter of Sebaste.
\switchcolumn*
\selectlanguage{latin}
\end{paracol}


% ---- martyrology/mart07/mart0720.htm
\needspace{10\baselineskip}
\begin{paracol}{2}
\selectlanguage{latin}
\begin{center}{\color{gregoriocolor} Tertiodécimo Kaléndas Augústi. 
 Luna\dots\ }\end{center}
\switchcolumn
\selectlanguage{english}
\begin{center}{\color{gregoriocolor} The 
 20\textsuperscript{th} Day of 
 July. The\dots\ Day of the Moon.}\end{center}
\end{paracol}

\noindent\begin{tabularx}{\linewidth}{*{19}{>{\centering\arraybackslash}X}}
 \textcolor{gregoriocolor}{a} & \textcolor{gregoriocolor}{b} & \textcolor{gregoriocolor}{c} & \textcolor{gregoriocolor}{d} & \textcolor{gregoriocolor}{e} & \textcolor{gregoriocolor}{f} & \textcolor{gregoriocolor}{g} & \textcolor{gregoriocolor}{h} & \textcolor{gregoriocolor}{i} & \textcolor{gregoriocolor}{k} & \textcolor{gregoriocolor}{l} & \textcolor{gregoriocolor}{m} & \textcolor{gregoriocolor}{n} & \textcolor{gregoriocolor}{p} & \textcolor{gregoriocolor}{q} & \textcolor{gregoriocolor}{r} & \textcolor{gregoriocolor}{s} & \textcolor{gregoriocolor}{t} & \textcolor{gregoriocolor}{u} \\
 25 & 26 & 27 & 28 & 29 & 30 & 1 & 2 & 3 & 4 & 5 & 6 & 7 & 8 & 9 & 10 & 11 & 12 & 13 \\
\end{tabularx}
\vspace{0.5\baselineskip}
\noindent\begin{tabularx}{\linewidth}{*{12}{>{\centering\arraybackslash}X}}
 \textcolor{gregoriocolor}{A} & \textcolor{gregoriocolor}{B} & \textcolor{gregoriocolor}{C} & \textcolor{gregoriocolor}{D} & \textcolor{gregoriocolor}{E} & F & \textcolor{gregoriocolor}{F} & \textcolor{gregoriocolor}{G} & \textcolor{gregoriocolor}{H} & \textcolor{gregoriocolor}{M} & \textcolor{gregoriocolor}{N} & \textcolor{gregoriocolor}{P} \\
 14 & 15 & 16 & 17 & 18 & 19 & 19 & 20 & 21 & 22 & 23 & 24 \\
\end{tabularx}

\begin{paracol}{2}
\selectlanguage{latin}
\lettrine[lines=2]{S}{ancti} Hierónymi 
 Æmiliáni Confessóris, Congregatiónis Somáschæ Institutóris, cæléstis ómnium 
 orphanórum ac derelíctæ juventútis Patróni; qui sexto Idus Februárii 
 obdormívit in Dómino.
\switchcolumn
\selectlanguage{english}
\lettrine[lines=2]{S}{t.} Jerome Emiliani, confessor, 
 founder of the Congregation of Somascha, the heavenly patron of all orphans 
 and destitute children. He fell asleep in the Lord on the 8th of 
 February.
\switchcolumn*
\selectlanguage{latin}
Antiochíæ pássio sanctæ 
 Margarítæ, Vírginis et Mártyris.
\switchcolumn
\selectlanguage{english}
At Antioch, the passion of St. 
 Margaret, virgin and martyr.
\switchcolumn*
\selectlanguage{latin}
In monte Carmélo sancti 
 Elíæ Prophétæ.
\switchcolumn
\selectlanguage{english}
On Mount Carmel, the holy prophet 
 Elijah.
\switchcolumn*
\selectlanguage{latin}
In Judæa natális beáti 
 Joseph, qui cognominátus est Justus, quem Apóstoli cum beáto Matthía 
 statuérunt ut locum apostolátus Judæ proditóris impléret; sed, cum sors 
 cecidísset, super Matthíam, ipse, nihilóminus prædicatiónis et sanctitátis 
 offício insérviens, multámque pro fide Christi a Judæis persecutiónem 
 sústinens, victorióso fine quiévit. De quo étiam refértur quod venénum 
 bíberit, et nihil ex hoc triste propter Dómini fidem pertúlerit.
\switchcolumn
\selectlanguage{english}
In Judea, the birthday of blessed 
 Joseph, surnamed the Just, whom the apostles selected with blessed Matthias 
 for the apostleship to replace the traitor Judas. The lot having 
 fallen upon Matthias, Joseph, notwithstanding, continued to preach and to 
 advance in virtue, and after having sustained from the Jews many 
 persecutions for the faith of Christ, he happily completed his life. 
 It is related of him that having drunk poison, he received no injury from 
 it, because of his confidence in the Lord.
\switchcolumn*
\selectlanguage{latin}
Córdubæ, in Hispánia, 
 sancti Pauli, Diáconi et Mártyris; qui, cum infidéles príncipes argúeret 
 Mahuméticæ impietátis ac sævítiæ, et Christum constantíssime prædicáret, 
 idcírco, eórum jussu necátus, ad præmia evolávit in cælum.
\switchcolumn
\selectlanguage{english}
At Cordova in Spain, St. Paul, 
 deacon and martyr. For rebuking Mohammedan princes for their impiety 
 and cruelty, and preaching Christ with constancy, he was put to death and 
 went to his reward in heaven.
\switchcolumn*
\selectlanguage{latin}
Damásci sanctórum 
 Mártyrum Sabíni, Juliáni, Máximi, Macróbii, Cássiæ et Paulæ, cum áliis decem.
\switchcolumn
\selectlanguage{english}
At Damascus, the holy martyrs 
 Sabinus, Julian, Maximus, Macrobius, Cassia, and Paul, with ten others.
\switchcolumn*
\selectlanguage{latin}
In Lusitánia sanctæ 
 Wilgefórtis, Vírginis et Mártyris; quæ, pro Christiána fide ac pudicítia 
 decértans, in cruce méruit gloriósum obtinére triúmphum.
\switchcolumn
\selectlanguage{english}
In Portugal, St. Wilgefortis, virgin 
 and martyr, who merited the crown of martyrdom on a cross in defence of the 
 faith and her chastity.
\switchcolumn*
\selectlanguage{latin}
Eódem die natális 
 sanctórum Flaviáni Secúndi, Epíscopi Antiochéni, et Elíæ, Epíscopi 
 Hierosolymitáni; qui, pro Synodo Chalcedonénsi ab Anastásio Imperatóre ambo 
 in exsílium acti, victóres migrárunt ad Dóminum.
\switchcolumn
\selectlanguage{english}
The same day, the birthday of St. 
 Flavian II, bishop of Antioch, and St. Elias, bishop of Jerusalem. 
 They were driven into exile by Emperor Anastasius for their defence of the 
 Council of Chalcedon, and there they went victoriously to the Lord.
\switchcolumn*
\selectlanguage{latin}
In pago Bononiénsi, in 
 Gállia, sancti Vulmári Abbátis, admirándæ sanctitátis viri.
\switchcolumn
\selectlanguage{english}
At Boulogne in France, the abbot St. 
 Wulmar, a man of admirable sanctity.
\switchcolumn*
\selectlanguage{latin}
Tréviris sanctæ Sevéræ 
 Vírginis.
\switchcolumn
\selectlanguage{english}
At Treves, St. Severa, virgin.
\switchcolumn*
\selectlanguage{latin}
\end{paracol}


% ---- martyrology/mart07/mart0721.htm
\needspace{10\baselineskip}
\begin{paracol}{2}
\selectlanguage{latin}
\begin{center}{\color{gregoriocolor} Duodécimo Kaléndas Augústi. 
 Luna\dots\ }\end{center}
\switchcolumn
\selectlanguage{english}
\begin{center}{\color{gregoriocolor} The 
 21\textsuperscript{st} Day of 
 July. The\dots\ Day of the Moon.}\end{center}
\end{paracol}

\noindent\begin{tabularx}{\linewidth}{*{19}{>{\centering\arraybackslash}X}}
 \textcolor{gregoriocolor}{a} & \textcolor{gregoriocolor}{b} & \textcolor{gregoriocolor}{c} & \textcolor{gregoriocolor}{d} & \textcolor{gregoriocolor}{e} & \textcolor{gregoriocolor}{f} & \textcolor{gregoriocolor}{g} & \textcolor{gregoriocolor}{h} & \textcolor{gregoriocolor}{i} & \textcolor{gregoriocolor}{k} & \textcolor{gregoriocolor}{l} & \textcolor{gregoriocolor}{m} & \textcolor{gregoriocolor}{n} & \textcolor{gregoriocolor}{p} & \textcolor{gregoriocolor}{q} & \textcolor{gregoriocolor}{r} & \textcolor{gregoriocolor}{s} & \textcolor{gregoriocolor}{t} & \textcolor{gregoriocolor}{u} \\
 26 & 27 & 28 & 29 & 30 & 1 & 2 & 3 & 4 & 5 & 6 & 7 & 8 & 9 & 10 & 11 & 12 & 13 & 14 \\
\end{tabularx}
\vspace{0.5\baselineskip}
\noindent\begin{tabularx}{\linewidth}{*{12}{>{\centering\arraybackslash}X}}
 \textcolor{gregoriocolor}{A} & \textcolor{gregoriocolor}{B} & \textcolor{gregoriocolor}{C} & \textcolor{gregoriocolor}{D} & \textcolor{gregoriocolor}{E} & F & \textcolor{gregoriocolor}{F} & \textcolor{gregoriocolor}{G} & \textcolor{gregoriocolor}{H} & \textcolor{gregoriocolor}{M} & \textcolor{gregoriocolor}{N} & \textcolor{gregoriocolor}{P} \\
 15 & 16 & 17 & 18 & 19 & 20 & 20 & 21 & 22 & 23 & 24 & 25 \\
\end{tabularx}

\begin{paracol}{2}
\selectlanguage{latin}
\lettrine[lines=2]{R}{omæ} sanctæ Praxédis 
 Vírginis, quæ, in omni castitáte et lege divína cum esset erudíta, vigíliis 
 et oratiónibus atque jejúniis assídue vacans, quiévit in Christo, sepúltaque est, juxta sorórem suam Pudentiánam, via Salária.
\switchcolumn
\selectlanguage{english}
\lettrine[lines=2]{A}{t} Rome, the holy virgin Praxedes, 
 who was brought up in all chastity and in the knowledge of the divine law. 
 Diligently attending to watching, prayer, and fasting, she rested in Christ, 
 and was buried near her sister Pudentiana on the Salarian Way.
\switchcolumn*
\selectlanguage{latin}
Babylóne sancti 
 Daniélis Prophétæ.
\switchcolumn
\selectlanguage{english}
At Babylon, the holy prophet Daniel.
\switchcolumn*
\selectlanguage{latin}
Cománæ, in Arménia, sancti Zótici, Epíscopi et Mártyris; qui sub Sevéro coronátus est.
\switchcolumn
\selectlanguage{english}
At Comana in Armenia, the holy 
 bishop and martyr Zoticus, who was crowned under Severus.
\switchcolumn*
\selectlanguage{latin}
Massíliæ, in Gállia, 
 natális sancti Victóris, qui, cum esset miles, et nec militáre neque idólis 
 sacrificáre vellet, hinc, primo in cárcerem trusus ibíque ab Angelo 
 visitátus, deínde váriis cruciátibus punítus, novíssime, contrítus in mola 
 pistória, martyrium consummávit. Passi sunt cum ipso et tres mílites, 
 id est Alexánder, Feliciánus et Longínus.
\switchcolumn
\selectlanguage{english}
At Marseilles in France, the 
 birthday of St. Victor, a soldier. Because he refused to serve in the 
 army and sacrifice to idols, he was thrust into prison, where he was visited 
 by an angel. He was subjected to various torments, and finally being 
 crushed under a millstone, he ended his martyrdom. With him also 
 suffered three soldiers, Alexander, Felician, and Longinus.
\switchcolumn*
\selectlanguage{latin}
Trecis, in Gállia, 
 pássio sanctórum Cláudii, Justi, Jucundíni et Sociórum quinque, sub 
 Aureliáno Imperatóre.
\switchcolumn
\selectlanguage{english}
At Troyes in France, the martyrdom 
 of the saints Claudius, Justus, Jucundinus, and five companions, in the time 
 of Emperor Aurelian.
\switchcolumn*
\selectlanguage{latin}
Ibídem sanctæ Júliæ, 
 Vírginis et Mártyris.
\switchcolumn
\selectlanguage{english}
In the same place, St. Julia, virgin 
 and martyr.
\switchcolumn*
\selectlanguage{latin}
Argentoráti sancti 
 Arbogásti Epíscopi, miráculis clari.
\switchcolumn
\selectlanguage{english}
At Strasbourg, St. Arbogast, a 
 bishop, renowned for miracles.
\switchcolumn*
\selectlanguage{latin}
In Syria sancti Joánnis 
 Mónachi, qui éxstitit colléga sancti Simeónis.
\switchcolumn
\selectlanguage{english}
In Syria, the holy monk John, a 
 companion of St. Simeon.
\switchcolumn*
\selectlanguage{latin}
\end{paracol}


% ---- martyrology/mart07/mart0722.htm
\needspace{10\baselineskip}
\begin{paracol}{2}
\selectlanguage{latin}
\begin{center}{\color{gregoriocolor} Undécimo Kaléndas Augústi. 
 Luna\dots\ }\end{center}
\switchcolumn
\selectlanguage{english}
\begin{center}{\color{gregoriocolor} The 
 22\textsuperscript{nd} Day of 
 July. The\dots\ Day of the Moon.}\end{center}
\end{paracol}

\noindent\begin{tabularx}{\linewidth}{*{19}{>{\centering\arraybackslash}X}}
 \textcolor{gregoriocolor}{a} & \textcolor{gregoriocolor}{b} & \textcolor{gregoriocolor}{c} & \textcolor{gregoriocolor}{d} & \textcolor{gregoriocolor}{e} & \textcolor{gregoriocolor}{f} & \textcolor{gregoriocolor}{g} & \textcolor{gregoriocolor}{h} & \textcolor{gregoriocolor}{i} & \textcolor{gregoriocolor}{k} & \textcolor{gregoriocolor}{l} & \textcolor{gregoriocolor}{m} & \textcolor{gregoriocolor}{n} & \textcolor{gregoriocolor}{p} & \textcolor{gregoriocolor}{q} & \textcolor{gregoriocolor}{r} & \textcolor{gregoriocolor}{s} & \textcolor{gregoriocolor}{t} & \textcolor{gregoriocolor}{u} \\
 27 & 28 & 29 & 30 & 1 & 2 & 3 & 4 & 5 & 6 & 7 & 8 & 9 & 10 & 11 & 12 & 13 & 14 & 15 \\
\end{tabularx}
\vspace{0.5\baselineskip}
\noindent\begin{tabularx}{\linewidth}{*{12}{>{\centering\arraybackslash}X}}
 \textcolor{gregoriocolor}{A} & \textcolor{gregoriocolor}{B} & \textcolor{gregoriocolor}{C} & \textcolor{gregoriocolor}{D} & \textcolor{gregoriocolor}{E} & F & \textcolor{gregoriocolor}{F} & \textcolor{gregoriocolor}{G} & \textcolor{gregoriocolor}{H} & \textcolor{gregoriocolor}{M} & \textcolor{gregoriocolor}{N} & \textcolor{gregoriocolor}{P} \\
 16 & 17 & 18 & 19 & 20 & 21 & 21 & 22 & 23 & 24 & 25 & 26 \\
\end{tabularx}

\begin{paracol}{2}
\selectlanguage{latin}
\lettrine[lines=2]{A}{pud} Massíliam, in 
 Gállia, natális sanctæ Maríæ Magdalénæ, de qua Dóminus ejécit septem dæmónia, 
 et quæ ipsum Salvatórem a mórtuis resurgéntem prima vidére méruit.
\switchcolumn
\selectlanguage{english}
\lettrine[lines=2]{A}{t} Marseilles in France, the 
 birthday of St. Mary Magdalene, out of whom our Lord expelled seven demons, 
 and who deserved to be the first to see the Saviour after he had risen from 
 the dead.
\switchcolumn*
\selectlanguage{latin}
Philíppis, in 
 Macedónia, sanctæ Syntyches, cujus méminit beátus Paulus Apóstolus.
\switchcolumn
\selectlanguage{english}
At Philippi in Macedonia, St. 
 Syntyche, mentioned by the blessed apostle Paul.
\switchcolumn*
\selectlanguage{latin}
Ancyræ, in Galátia, natális sancti Platónis Mártyris, qui, sub Agrippíno Vicário, verbéribus 
 cæsus, uncis laniátus férreis, aliísque immaníssimis tormentórum genéribus 
 cruciátus, demum, abscísso cápite, invíctam ánimam Deo réddidit. 
 Ipsíus vero mirácula in subveniéndis captívis, Acta secúndæ Synodi Nicǽnæ 
 testántur.
\switchcolumn
\selectlanguage{english}
At Ancyra in Galatia, the birthday 
 of the martyr St. Plato. Under the lieutenant-governor Agrippinus, he 
 was scourged, lacerated with iron hooks, and subjected to the most atrocious 
 torments, and finally being beheaded, he rendered his invincible soul to 
 God. The Acts of the Second Council of Nicaea bear witness to his 
 miracles in helping captives.
\switchcolumn*
\selectlanguage{latin}
In Cypro sancti 
 Theóphili Prætóris, qui ab Arábibus tentus, et, cum nec donis nec minis 
 flecti posset ut Christum negáret, gládio cæsus est.
\switchcolumn
\selectlanguage{english}
In Cyprus, St. Theophilus, a 
 praetor, who was apprehended by the Arabs, and as he could not be induced 
 either by gifts or by threats to deny Christ, was put to the sword.
\switchcolumn*
\selectlanguage{latin}
Antiochíæ sancti 
 Cyrílli Epíscopi, doctrína et sanctitáte conspícui.
\switchcolumn
\selectlanguage{english}
At Antioch, the holy bishop Cyril, 
 who was distinguished for learning and holiness.
\switchcolumn*
\selectlanguage{latin}
Menáti, in território 
 Arvernénsi, sancti Meneléi Abbátis.
\switchcolumn
\selectlanguage{english}
At Menat, in the territory of 
 Auvergne, St. Meneleus, abbot.
\switchcolumn*
\selectlanguage{latin}
In monastério 
 Fontanéllæ, in Gállia, sancti Wandregísili Abbátis, miráculis clari; cujus 
 corpus ad Blandínum monastérium, in Flándria, póstea delátum fuit.
\switchcolumn
\selectlanguage{english}
In the monastery of Fontanelle in 
 France, Abbot St. Wandrille, famous for his miracles. His body was 
 afterwards translated to the monastery of Blandin, in Flanders.
\switchcolumn*
\selectlanguage{latin}
Ulyssipóne, in 
 Lusitánia, sancti Lauréntii a Brundísio, Sacerdótis et Confessóris; qui 
 Ordinis Minórum sancti Francísci Capuccinórum Miníster éxstitit Generális, 
 atque, divíni verbi prædicatióne et árduis pro Dei glória gestis præclárus, 
 a Leóne Décimo tértio, Summo Pontífice, Sanctórum fastis adscríptus est.
\switchcolumn
\selectlanguage{english}
At Lisbon in Portugal, St. Lawrence 
 of Brindisi, priest and confessor, superior general of the Order of Friars 
 Minor Capuchin of St. Francis. Illustrious for his preaching and his 
 arduous labour for the glory of God, he was canonized by Pope Leo XIII.
\switchcolumn*
\selectlanguage{latin}
Scythópoli, in 
 Palæstina, sancti Joséphi Cómitis.
\switchcolumn
\selectlanguage{english}
At Scythopolis in Palestine, St. 
 Joseph, a count.
\switchcolumn*
\selectlanguage{latin}
\end{paracol}


% ---- martyrology/mart07/mart0723.htm
\needspace{10\baselineskip}
\begin{paracol}{2}
\selectlanguage{latin}
\begin{center}{\color{gregoriocolor} Décimo Kaléndas Augústi. 
 Luna\dots\ }\end{center}
\switchcolumn
\selectlanguage{english}
\begin{center}{\color{gregoriocolor} The 
 23\textsuperscript{rd} Day of 
 July. The\dots\ Day of the Moon.}\end{center}
\end{paracol}

\noindent\begin{tabularx}{\linewidth}{*{19}{>{\centering\arraybackslash}X}}
 \textcolor{gregoriocolor}{a} & \textcolor{gregoriocolor}{b} & \textcolor{gregoriocolor}{c} & \textcolor{gregoriocolor}{d} & \textcolor{gregoriocolor}{e} & \textcolor{gregoriocolor}{f} & \textcolor{gregoriocolor}{g} & \textcolor{gregoriocolor}{h} & \textcolor{gregoriocolor}{i} & \textcolor{gregoriocolor}{k} & \textcolor{gregoriocolor}{l} & \textcolor{gregoriocolor}{m} & \textcolor{gregoriocolor}{n} & \textcolor{gregoriocolor}{p} & \textcolor{gregoriocolor}{q} & \textcolor{gregoriocolor}{r} & \textcolor{gregoriocolor}{s} & \textcolor{gregoriocolor}{t} & \textcolor{gregoriocolor}{u} \\
 28 & 29 & 30 & 1 & 2 & 3 & 4 & 5 & 6 & 7 & 8 & 9 & 10 & 11 & 12 & 13 & 14 & 15 & 16 \\
\end{tabularx}
\vspace{0.5\baselineskip}
\noindent\begin{tabularx}{\linewidth}{*{12}{>{\centering\arraybackslash}X}}
 \textcolor{gregoriocolor}{A} & \textcolor{gregoriocolor}{B} & \textcolor{gregoriocolor}{C} & \textcolor{gregoriocolor}{D} & \textcolor{gregoriocolor}{E} & F & \textcolor{gregoriocolor}{F} & \textcolor{gregoriocolor}{G} & \textcolor{gregoriocolor}{H} & \textcolor{gregoriocolor}{M} & \textcolor{gregoriocolor}{N} & \textcolor{gregoriocolor}{P} \\
 17 & 18 & 19 & 20 & 21 & 22 & 22 & 23 & 24 & 25 & 26 & 27 \\
\end{tabularx}

\begin{paracol}{2}
\selectlanguage{latin}
\lettrine[lines=2]{R}{avénnæ} natális sancti 
 Apollináris Epíscopi, qui, ab Apóstolo Petro Romæ ordinátus et Ravénnam 
 missus, pro fide Christi divérsas et multíplices pœnas perpéssus est; póstea, 
 Evangélium in Æmília prædicans, plúrimos ab idolórum cultu revocávit; 
 tandem, Ravénnam revérsus, gloriósum martyrium, sub Vespasiáno Cæsare, 
 complévit.
\switchcolumn
\selectlanguage{english}
\lettrine[lines=2]{A}{t} Ravenna, the birthday of the holy 
 bishop Apollinaris, who was consecrated at Rome by the Apostle Peter, and 
 sent to Ravenna, where he endured many different tribulations for the faith 
 of Christ. He afterwards preached the Gospel in Emilia, where he 
 converted many from the worship of idols. Finally, returning to 
 Ravenna, he completed his confession of Christ by a glorious martyrdom under 
 Vespasian Caesar.
\switchcolumn*
\selectlanguage{latin}
Cenómanis, in Gállia, 
 sancti Libórii, Epíscopi et Confessóris.
\switchcolumn
\selectlanguage{english}
At Le Mans in France, St. Liborius, 
 bishop and confessor.
\switchcolumn*
\selectlanguage{latin}
Romæ natális sanctæ 
 Birgíttæ Víduæ, quæ, post multas sanctórum locórum peregrinatiónes, divíno 
 affláta Spíritu, quiévit. Ipsíus autem festívitas octávo Idus Octóbris 
 celebrátur.
\switchcolumn
\selectlanguage{english}
At Rome, St. Bridget, widow, who, 
 after many pilgrimages to the holy places, fell asleep filled with the 
 Spirit of God. Her feast is observed on the 8th of October.
\switchcolumn*
\selectlanguage{latin}
Ibídem sancti Rásyphi 
 Mártyris.
\switchcolumn
\selectlanguage{english}
Also, St. Rasyphus, martyr.
\switchcolumn*
\selectlanguage{latin}
Item Romæ pássio sanctæ 
 Primitívæ, Vírginis et Mártyris.
\switchcolumn
\selectlanguage{english}
In the same city, the martyrdom of 
 St. Primitiva, virgin and martyr.
\switchcolumn*
\selectlanguage{latin}
Item sanctórum Mártyrum 
 Apollónii et Eugénii.
\switchcolumn
\selectlanguage{english}
Also the holy martyrs Apollonius and 
 Eugene.
\switchcolumn*
\selectlanguage{latin}
Eódem die natális 
 sanctórum Mártyrum Tróphimi et Theóphili, qui, sub Diocletiáno Imperatóre, 
 cæsi lapídibus et igne incénsi, demum, gládio percússi, martyrio coronántur.
\switchcolumn
\selectlanguage{english}
The same day, the birthday of the 
 holy martyrs Trophimus and Theophilus, who received their crown of martyrdom 
 by being beaten with stones, scorched with fire, and finally struck with the 
 sword, in the time of Emperor Diocletian.
\switchcolumn*
\selectlanguage{latin}
In Bulgária sanctórum 
 plurimórum Mártyrum, quos ímpius Imperátor Nicéphorus, Ecclésias Dei 
 devástans, divérso mortis génere, nimírum ense, láqueo, sagíttis, diútino cárcere et fame necári fecit.
\switchcolumn
\selectlanguage{english}
In Bulgaria, many holy martyrs, whom 
 the impious Emperor Nicephorus, while devastating the churches of God, put 
 to death in various ways: by the sword, by hanging, arrows, long 
 imprisonment, and by starvation.
\switchcolumn*
\selectlanguage{latin}
Romæ sanctárum Vírginum 
 Rómulæ, Redémptæ et Herúndinis, de quibus scribit sanctus Gregórius Papa.
\switchcolumn
\selectlanguage{english}
At Rome, the saintly virgins Romula, 
 Redempta, and Herundo, mentioned by Pope St. Gregory in his writings.
\switchcolumn*
\selectlanguage{latin}
\end{paracol}


% ---- martyrology/mart07/mart0724.htm
\needspace{10\baselineskip}
\begin{paracol}{2}
\selectlanguage{latin}
\begin{center}{\color{gregoriocolor} Nono Kaléndas Augústi. 
 Luna\dots\ }\end{center}
\switchcolumn
\selectlanguage{english}
\begin{center}{\color{gregoriocolor} The 
 24\textsuperscript{th} Day of 
 July. The\dots\ Day of the Moon.}\end{center}
\end{paracol}

\noindent\begin{tabularx}{\linewidth}{*{19}{>{\centering\arraybackslash}X}}
 \textcolor{gregoriocolor}{a} & \textcolor{gregoriocolor}{b} & \textcolor{gregoriocolor}{c} & \textcolor{gregoriocolor}{d} & \textcolor{gregoriocolor}{e} & \textcolor{gregoriocolor}{f} & \textcolor{gregoriocolor}{g} & \textcolor{gregoriocolor}{h} & \textcolor{gregoriocolor}{i} & \textcolor{gregoriocolor}{k} & \textcolor{gregoriocolor}{l} & \textcolor{gregoriocolor}{m} & \textcolor{gregoriocolor}{n} & \textcolor{gregoriocolor}{p} & \textcolor{gregoriocolor}{q} & \textcolor{gregoriocolor}{r} & \textcolor{gregoriocolor}{s} & \textcolor{gregoriocolor}{t} & \textcolor{gregoriocolor}{u} \\
 29 & 30 & 1 & 2 & 3 & 4 & 5 & 6 & 7 & 8 & 9 & 10 & 11 & 12 & 13 & 14 & 15 & 16 & 17 \\
\end{tabularx}
\vspace{0.5\baselineskip}
\noindent\begin{tabularx}{\linewidth}{*{12}{>{\centering\arraybackslash}X}}
 \textcolor{gregoriocolor}{A} & \textcolor{gregoriocolor}{B} & \textcolor{gregoriocolor}{C} & \textcolor{gregoriocolor}{D} & \textcolor{gregoriocolor}{E} & F & \textcolor{gregoriocolor}{F} & \textcolor{gregoriocolor}{G} & \textcolor{gregoriocolor}{H} & \textcolor{gregoriocolor}{M} & \textcolor{gregoriocolor}{N} & \textcolor{gregoriocolor}{P} \\
 18 & 19 & 20 & 21 & 22 & 23 & 23 & 24 & 25 & 26 & 27 & 28 \\
\end{tabularx}

\begin{paracol}{2}
\selectlanguage{latin}
\lettrine[lines=1]{V}{igília} sancti Jacóbi Apóstoli.
\switchcolumn
\selectlanguage{english}
\lettrine[lines=1]{T}{he} Vigil of St. James the Apostle.
\switchcolumn*
\selectlanguage{latin}
Tyri, apud lacum 
 Vulsínium, in Túscia, sanctæ Christínæ, Vírginis et Mártyris. Hæc 
 Virgo, cum patris idóla áurea et argéntea, in Christum credens, comminuísset 
 atque illórum frágmina paupéribus erogásset, ejúsdem patris jussu verbéribus 
 dilaniáta est, aliísque supplíciis diríssime cruciáta, et cum magno saxi 
 póndere in lacum projécta, sed ab Angelo liberáta; deínde, sub álio Júdice, 
 patris sui successóre, acerbióra torménta constánter pértulit; novíssime, 
 sub Juliáno Præside, post fornácem ardéntem, ubi quinque diébus illæsa 
 permánsit, post serpéntes virtúte Christi superátos, martyrii sui cursum 
 abscissióne linguæ et sagittárum infixióne complévit.
\switchcolumn
\selectlanguage{english}
At Tiro in Tuscany, on Lake Bolsena, 
 St. Christina, virgin and martyr. Because she believed in Christ, and 
 broke up her father's gold and silver idols to give them to the poor, she 
 was cruelly scourged at his command, subjected to other most severe 
 torments, and thrown with a heavy stone into the lake from which she was 
 drawn out by an angel. Then under another judge, who succeeded her 
 father, she bore courageously still more bitter tortures. Finally, 
 after she had been shut up by the governor Julian in a burning furnace for 
 five days without any injury, after being cured of the sting of serpents, 
 she ended her martyrdom by having her tongue cut out, and being pierced with 
 arrows.
\switchcolumn*
\selectlanguage{latin}
Romæ, via Tiburtína, 
 sancti Vincéntii Mártyris.
\switchcolumn
\selectlanguage{english}
At Rome, on the Tiburtine Way, St. 
 Vincent, martyr.
\switchcolumn*
\selectlanguage{latin}
Amitérni, in Vestínis, 
 pássio sanctórum mílitum octogínta trium.
\switchcolumn
\selectlanguage{english}
At Amiterno in Abruzzi, the 
 martyrdom of eighty-three holy soldiers.
\switchcolumn*
\selectlanguage{latin}
Eméritæ, in Hispánia, 
 sancti Victóris, viri militáris, qui, cum frátribus Stercátio et Antinógene, 
 in persecutióne Diocletiáni, per divérsa supplícia martyrium consummávit.
\switchcolumn
\selectlanguage{english}
At Merida in Spain, St. Victor, a 
 soldier who, with his two brothers, Stercatius and Antinogenes, by divers 
 torments fulfilled his martyrdom in the persecution of Diocletian.
\switchcolumn*
\selectlanguage{latin}
Item sanctórum Mártyrum 
 Menéi et Capitónis.
\switchcolumn
\selectlanguage{english}
Also, the holy martyrs Meneus and 
 Capito.
\switchcolumn*
\selectlanguage{latin}
In Lycia sanctárum 
 Mártyrum Nicétæ et Aquilínæ, quæ, beáti Christóphori Mártyris prædicatióne 
 ad Christum convérsæ, martyrii palmam obtruncatióne cápitis sumpsérunt.
\switchcolumn
\selectlanguage{english}
In Lycia, the holy martyrs Niceta 
 and Aquilina, who were converted to Christ by the preaching of the blessed 
 martyr Christopher, and gained the palm of martyrdom by being beheaded.
\switchcolumn*
\selectlanguage{latin}
Apud Sénonas sancti 
 Ursicíni, Epíscopi et Confessóris.
\switchcolumn
\selectlanguage{english}
At Sens, St. Ursicinus, bishop and 
 confessor.
\switchcolumn*
\selectlanguage{latin}
\end{paracol}


% ---- martyrology/mart07/mart0725.htm
\needspace{10\baselineskip}
\begin{paracol}{2}
\selectlanguage{latin}
\begin{center}{\color{gregoriocolor} Octávo Kaléndas Augústi. 
 Luna\dots\ }\end{center}
\switchcolumn
\selectlanguage{english}
\begin{center}{\color{gregoriocolor} The 
 25\textsuperscript{th} Day of 
 July. The\dots\ Day of the Moon.}\end{center}
\end{paracol}

\noindent\begin{tabularx}{\linewidth}{*{19}{>{\centering\arraybackslash}X}}
 \textcolor{gregoriocolor}{a} & \textcolor{gregoriocolor}{b} & \textcolor{gregoriocolor}{c} & \textcolor{gregoriocolor}{d} & \textcolor{gregoriocolor}{e} & \textcolor{gregoriocolor}{f} & \textcolor{gregoriocolor}{g} & \textcolor{gregoriocolor}{h} & \textcolor{gregoriocolor}{i} & \textcolor{gregoriocolor}{k} & \textcolor{gregoriocolor}{l} & \textcolor{gregoriocolor}{m} & \textcolor{gregoriocolor}{n} & \textcolor{gregoriocolor}{p} & \textcolor{gregoriocolor}{q} & \textcolor{gregoriocolor}{r} & \textcolor{gregoriocolor}{s} & \textcolor{gregoriocolor}{t} & \textcolor{gregoriocolor}{u} \\
 30 & 1 & 2 & 3 & 4 & 5 & 6 & 7 & 8 & 9 & 10 & 11 & 12 & 13 & 14 & 15 & 16 & 17 & 18 \\
\end{tabularx}
\vspace{0.5\baselineskip}
\noindent\begin{tabularx}{\linewidth}{*{12}{>{\centering\arraybackslash}X}}
 \textcolor{gregoriocolor}{A} & \textcolor{gregoriocolor}{B} & \textcolor{gregoriocolor}{C} & \textcolor{gregoriocolor}{D} & \textcolor{gregoriocolor}{E} & F & \textcolor{gregoriocolor}{F} & \textcolor{gregoriocolor}{G} & \textcolor{gregoriocolor}{H} & \textcolor{gregoriocolor}{M} & \textcolor{gregoriocolor}{N} & \textcolor{gregoriocolor}{P} \\
 19 & 20 & 21 & 22 & 23 & 24 & 24 & 25 & 26 & 27 & 28 & 29 \\
\end{tabularx}

\begin{paracol}{2}
\selectlanguage{latin}
\lettrine[lines=2]{S}{ancti} Jacóbi Apóstoli, 
 qui éxstitit beáti Joánnis Evangelístæ frater; et, prope festum Paschæ ab 
 Heróde Agríppa decollátus, primus ex Apóstolis corónam martyrii percépit. 
 Ejus sacra ossa, ab Hierosólymis ad Hispánias hoc die transláta, et in 
 últimis eárum fínibus apud Gallæciam recóndita, celebérrima illárum géntium 
 veneratióne, et frequénti Christianórum concúrsu, religiónis et voti causa 
 illuc adeúntium, pie colúntur.
\switchcolumn
\selectlanguage{english}
\lettrine[lines=2]{S}{t.} James the Apostle, brother of 
 the blessed evangelist John, who was beheaded by Herod Agrippa at about the 
 feast of Easter. He was the first of the apostles to receive the crown 
 of martyrdom. His sacred bones were on this day carried from Jerusalem 
 to Spain, and placed in the remote province of Galicia, where they are 
 devoutly honoured by the far-famed piety of the inhabitants, and the 
 frequent concourse of Christians, who visit them through piety and in fulfillment of vows.
\switchcolumn*
\selectlanguage{latin}
In Lycia sancti 
 Christóphori Mártyris, qui, sub Décio, virgis férreis attrítus, et a flammæ 
 æstuántis incéndio supérna Christi virtúte servátus, ad últimum, sagittárum 
 íctibus confóssus, cápitis obtruncatióne martyrium complévit.
\switchcolumn
\selectlanguage{english}
In Lycia, in the time of Decius, St. 
 Christopher, martyr. Being scourged with iron rods, cast into the 
 flames, from which he was saved by the power of Christ, and finally 
 transfixed with arrows and beheaded, he completed his martyrdom.
\switchcolumn*
\selectlanguage{latin}
Barcinóne, in Hispánia, 
 natális beáti Cucuphátis Mártyris, qui, in persecutióne Diocletiáni, sub 
 Daciáno Præside, plúrimis torméntis superátis, tandem, percússus gládio, 
 victor migrávit in cælum.
\switchcolumn
\selectlanguage{english}
At Barcelona in Spain, during the 
 persecution of Diocletian and under the governor Dacian, the birthday of the 
 holy martyr Cucuphas. After overcoming many torments, he was struck 
 with the sword, and thus went triumphantly to heaven.
\switchcolumn*
\selectlanguage{latin}
In Palæstína sancti 
 Pauli Mártyris, qui, in persecutióne Maximiáni Galérii, sub Firmiliáno 
 Præside, cápitis damnátus, et, exíguo orándi spátio impetráto, primum pro 
 contribúlibus suis, deínde pro Judæis qui ipsum damnáverat, et pro carnífice 
 a quo feriéndus erat, Deum toto corde precátus, martyrii corónam, abscíssis 
 cervícibus, accépit.
\switchcolumn
\selectlanguage{english}
In Palestine, St. Paul, a martyr in 
 the persecution of Maximian Galerius, under the governor Firmilian. He 
 was condemned to death, but having obtained a short period for prayer, he 
 besought God with all his heart, first for his own countrymen, then for the 
 Jews and the Gentiles, that they might embrace the true faith, next for the 
 multitude of spectators, and finally for the judge who had condemned him and 
 the executioner who was to strike him; after which he received the crown of 
 martyrdom by beheading.
\switchcolumn*
\selectlanguage{latin}
Furcónii, in Vestínis, 
 sanctórum Mártyrum Sipontinórum Floréntii et Felícis.
\switchcolumn
\selectlanguage{english}
At Forcono in Abruzzi, the holy 
 martyrs Florentius and Felix, natives of Siponte.
\switchcolumn*
\selectlanguage{latin}
Córdubæ, in Hispánia, 
 sancti Theodemíri, Mónachi et Mártyris.
\switchcolumn
\selectlanguage{english}
At Cordova, St. Theodemir, monk and 
 martyr.
\switchcolumn*
\selectlanguage{latin}
In Palæstína sanctæ 
 Valentínæ Vírginis, quæ, cum ad aram, ut immoláret, addúcta esset, eámque 
 cálcibus evertísset, prius diríssime cruciáta est; deínde, una cum sócia 
 Vírgine, in ignem conjécta cucúrrit ad Sponsum.
\switchcolumn
\selectlanguage{english}
In Palestine, St. Valentina, a 
 virgin, who was led to an altar to offer sacrifice, but overturning it with 
 her foot, she was cruelly tortured, and being cast into the fire with 
 another virgin, her companion, she went to her Spouse.
\switchcolumn*
\selectlanguage{latin}
Tréviris sancti 
 Magneríci, Epíscopi et Confessóris.
\switchcolumn
\selectlanguage{english}
At Treves, St. Magnericus, bishop 
 and confessor.
\switchcolumn*
\selectlanguage{latin}
\end{paracol}


% ---- martyrology/mart07/mart0726.htm
\needspace{10\baselineskip}
\begin{paracol}{2}
\selectlanguage{latin}
\begin{center}{\color{gregoriocolor} Séptimo Kaléndas Augústi. 
 Luna\dots\ }\end{center}
\switchcolumn
\selectlanguage{english}
\begin{center}{\color{gregoriocolor} The 
 26\textsuperscript{th} Day of 
 July. The\dots\ Day of the Moon.}\end{center}
\end{paracol}

\noindent\begin{tabularx}{\linewidth}{*{19}{>{\centering\arraybackslash}X}}
 \textcolor{gregoriocolor}{a} & \textcolor{gregoriocolor}{b} & \textcolor{gregoriocolor}{c} & \textcolor{gregoriocolor}{d} & \textcolor{gregoriocolor}{e} & \textcolor{gregoriocolor}{f} & \textcolor{gregoriocolor}{g} & \textcolor{gregoriocolor}{h} & \textcolor{gregoriocolor}{i} & \textcolor{gregoriocolor}{k} & \textcolor{gregoriocolor}{l} & \textcolor{gregoriocolor}{m} & \textcolor{gregoriocolor}{n} & \textcolor{gregoriocolor}{p} & \textcolor{gregoriocolor}{q} & \textcolor{gregoriocolor}{r} & \textcolor{gregoriocolor}{s} & \textcolor{gregoriocolor}{t} & \textcolor{gregoriocolor}{u} \\
 1 & 2 & 3 & 4 & 5 & 6 & 7 & 8 & 9 & 10 & 11 & 12 & 13 & 14 & 15 & 16 & 17 & 18 & 19 \\
\end{tabularx}
\vspace{0.5\baselineskip}
\noindent\begin{tabularx}{\linewidth}{*{12}{>{\centering\arraybackslash}X}}
 \textcolor{gregoriocolor}{A} & \textcolor{gregoriocolor}{B} & \textcolor{gregoriocolor}{C} & \textcolor{gregoriocolor}{D} & \textcolor{gregoriocolor}{E} & F & \textcolor{gregoriocolor}{F} & \textcolor{gregoriocolor}{G} & \textcolor{gregoriocolor}{H} & \textcolor{gregoriocolor}{M} & \textcolor{gregoriocolor}{N} & \textcolor{gregoriocolor}{P} \\
 20 & 21 & 22 & 23 & 24 & 25 & 25 & 26 & 27 & 28 & 29 & 30 \\
\end{tabularx}

\begin{paracol}{2}
\selectlanguage{latin}
\lettrine[lines=2]{D}{ormítio} sanctæ Annæ, 
 quæ mater éxstitit Immaculátæ Vírginis Genitrícis Dei Maríæ.
\switchcolumn
\selectlanguage{english}
\lettrine[lines=2]{T}{he} departure from this life of St. 
 Anne, mother of the Immaculate Virgin Mary, the Mother of God.
\switchcolumn*
\selectlanguage{latin}
Philíppis, in 
 Macedónia, natális sancti Erásti, qui, illic a beáto Paulo Apóstolo relíctus 
 Epíscopus, ibídem martyrio coronátus est.
\switchcolumn
\selectlanguage{english}
At Philippi in Macedonia, the 
 birthday of St. Erastus, who was appointed bishop of that place by the 
 blessed apostle Paul, and was there crowned with martyrdom.
\switchcolumn*
\selectlanguage{latin}
Romæ, via Latína, sanctórum Mártyrum Symphrónii, Olympii, Theodúli et Exsupériæ; qui (ut in 
 gestis sancti Stéphani Papæ légitur), ígnibus combústi, martyrii palmam 
 adépti sunt.
\switchcolumn
\selectlanguage{english}
At Rome, on the Latin Way, the holy 
 martyrs Symphronius, Olympius, Theodulus, and Exuperia, who (as we read in 
 the Acts of Pope St. Stephen) were burned alive, and thus obtained the palm 
 of martyrdom.
\switchcolumn*
\selectlanguage{latin}
In Portu Románo sancti 
 Hyacínthi Mártyris, qui, primo in ignem injéctus, deínde in profluéntem 
 præcipitátus, illæsus evásit; post hæc, sub Trajáno Imperatóre, a Leóntio 
 Consulári percússus gládio, vitam finívit. Ipsíus corpus Júlia Matróna 
 in prædio suo, juxta Urbem, sepelívit.
\switchcolumn
\selectlanguage{english}
At Porto, St. Hyacinth, martyr, who 
 was first thrown into the fire, and then cast into a stream without being 
 injured. Afterwards, under Emperor Trajan, being struck with the sword 
 by the exconsul Leontius, his martyrdom was fulfilled. His body was 
 buried by the matron Julia on her own estate near Rome.
\switchcolumn*
\selectlanguage{latin}
Verónæ sancti Valéntis, 
 Epíscopi et Confessóris.
\switchcolumn
\selectlanguage{english}
At Verona, St. Valens, bishop and 
 confessor.
\switchcolumn*
\selectlanguage{latin}
Romæ sancti Pastóris 
 Presbyteri, cujus nómine Títulus exstat in Vimináli, apud sanctam 
 Pudentiánam.
\switchcolumn
\selectlanguage{english}
At Rome, St. Pastor, a priest in 
 whose name a title exists in the church of St. Pudentiana, on the Viminal 
 Hill.
\switchcolumn*
\selectlanguage{latin}
In monastério sancti 
 Benedícti, in agro Mantuáno, sancti Simeónis, Mónachi et Eremítæ, qui, 
 multis miráculis clarus, in senectúte bona quiévit.
\switchcolumn
\selectlanguage{english}
In the monastery of St. Benedict, 
 near Mantua, St. Simeon, monk and hermit, who was renowned for many 
 miracles, and at an advanced age rested in the Lord.
\switchcolumn*
\selectlanguage{latin}
Lúere, in diœcési 
 Brixiénsi, sanctæ Bartholomǽæ Capitánio, Vírginis, Sorórum a Caritáte 
 Fundatrícis, puéllis instituéndis præcláræ, quam Pius Papa Duodécimus albo 
 sanctárum Vírginum adscrípsit.
\switchcolumn
\selectlanguage{english}
At Lovere, in the diocese of Brescia, 
 St. Bartholemea Capitanio, virgin, who founded the Sisters of Charity, 
 dedicated to teaching the young. Pope Pius XII added her name to the 
 catalogue of holy virgins.
\switchcolumn*
\selectlanguage{latin}
\end{paracol}


% ---- martyrology/mart07/mart0727.htm
\needspace{10\baselineskip}
\begin{paracol}{2}
\selectlanguage{latin}
\begin{center}{\color{gregoriocolor} Sexto Kaléndas Augústi. 
 Luna\dots\ }\end{center}
\switchcolumn
\selectlanguage{english}
\begin{center}{\color{gregoriocolor} The 
 27\textsuperscript{th} Day of 
 July. The\dots\ Day of the Moon.}\end{center}
\end{paracol}

\noindent\begin{tabularx}{\linewidth}{*{19}{>{\centering\arraybackslash}X}}
 \textcolor{gregoriocolor}{a} & \textcolor{gregoriocolor}{b} & \textcolor{gregoriocolor}{c} & \textcolor{gregoriocolor}{d} & \textcolor{gregoriocolor}{e} & \textcolor{gregoriocolor}{f} & \textcolor{gregoriocolor}{g} & \textcolor{gregoriocolor}{h} & \textcolor{gregoriocolor}{i} & \textcolor{gregoriocolor}{k} & \textcolor{gregoriocolor}{l} & \textcolor{gregoriocolor}{m} & \textcolor{gregoriocolor}{n} & \textcolor{gregoriocolor}{p} & \textcolor{gregoriocolor}{q} & \textcolor{gregoriocolor}{r} & \textcolor{gregoriocolor}{s} & \textcolor{gregoriocolor}{t} & \textcolor{gregoriocolor}{u} \\
 2 & 3 & 4 & 5 & 6 & 7 & 8 & 9 & 10 & 11 & 12 & 13 & 14 & 15 & 16 & 17 & 18 & 19 & 20 \\
\end{tabularx}
\vspace{0.5\baselineskip}
\noindent\begin{tabularx}{\linewidth}{*{12}{>{\centering\arraybackslash}X}}
 \textcolor{gregoriocolor}{A} & \textcolor{gregoriocolor}{B} & \textcolor{gregoriocolor}{C} & \textcolor{gregoriocolor}{D} & \textcolor{gregoriocolor}{E} & F & \textcolor{gregoriocolor}{F} & \textcolor{gregoriocolor}{G} & \textcolor{gregoriocolor}{H} & \textcolor{gregoriocolor}{M} & \textcolor{gregoriocolor}{N} & \textcolor{gregoriocolor}{P} \\
 21 & 22 & 23 & 24 & 25 & 26 & 26 & 27 & 28 & 29 & 30 & 1 \\
\end{tabularx}

\begin{paracol}{2}
\selectlanguage{latin}
\lettrine[lines=2]{N}{icomedíæ} pássio sancti 
 Pantaleónis médici, qui, pro fide Christi, a Maximiáno Imperatóre tentus, et 
 equúlei pœna ac lampadárum exustióne afflíctus, sed inter hæc, Dómino sibi 
 apparénte, refrigerátus, gládii tandem ictu martyrium consummávit.
\switchcolumn
\selectlanguage{english}
\lettrine[lines=2]{A}{t} Nicomedia, the martyrdom of St. 
 Pantaleon, a physician. For the faith of Christ he was apprehended by 
 Emperor Maximian, subjected to the torture and burned with torches, during 
 which torments he was comforted by an apparition of our Lord. He ended 
 his martyrdom by a stroke of the sword.
\switchcolumn*
\selectlanguage{latin}
Vigíliis, in Apúlia, sanctórum Mártyrum Mauri Epíscopi, Pantaleémonis et Sérgii; qui passi sunt 
 sub Trajáno.
\switchcolumn
\selectlanguage{english}
At Bisceglia in Apulia, the holy 
 martyrs Maur, a bishop, Pantaleon, and Sergius, who suffered under Trajan.
\switchcolumn*
\selectlanguage{latin}
Nicomedíæ sancti 
 Hermolái Presbyteri, cujus doctrína beátus Pantáleon ad fidem convérsus est; 
 itémque sanctórum Hermíppi et Hermócratis fratrum, qui, post multas pœnas 
 sibi illátas, capitáli senténtia, ob confessiónem Christi, a Maximiáno 
 Imperatóre puníti sunt.
\switchcolumn
\selectlanguage{english}
At Nicomedia, St. Hermolaus, priest, 
 by whose instructions blessed Pantaleon was converted to the faith. 
 Also, the Saints Hermippus and Hermocrates, brothers. After many 
 sufferings borne for the confession of Christ, they were condemned to death 
 by the same Maximian.
\switchcolumn*
\selectlanguage{latin}
Córdubæ, in Hispánia, 
 sanctórum Mártyrum Geórgii Diáconi, Aurélii et uxóris Natáliæ, Felícis et 
 uxóris Liliósæ, in persecutióne Arábica.
\switchcolumn
\selectlanguage{english}
At Cordova in Spain, during the Arab 
 persecution, the holy martyrs George, a deacon, Aurelius and his wife 
 Natalia, Felix and his wife Liliosa.
\switchcolumn*
\selectlanguage{latin}
Nolæ, in Campánia, sanctórum Mártyrum Felícis, Júliæ et Jucúndæ.
\switchcolumn
\selectlanguage{english}
At Nola in Campania, the holy 
 martyrs Felix, Julia, and Jucunda.
\switchcolumn*
\selectlanguage{latin}
Apud Homerítas, in 
 Arábia, commemorátio sanctórum Mártyrum, qui, sub Dúnaan tyránno, ob Christi 
 fidem, incéndio tráditi sunt.
\switchcolumn
\selectlanguage{english}
In the country of the Homerites in 
 Arabia, the commemoration of the holy martyrs, who were delivered to the 
 flames for the faith of Christ under the tyrant Dunaan.
\switchcolumn*
\selectlanguage{latin}
Ephesi natális 
 sanctórum septem Dormiéntium, scílicet Maximiáni, Malchi, Martiniáni, 
 Dionysii, Joánnis, Serapiónis et Constantíni.
\switchcolumn
\selectlanguage{english}
At Ephesus, the birthday of the 
 Seven Holy Sleepers, Maximian, Malchus, Martinian, Denis, John, Serapion, 
 and Constantine.
\switchcolumn*
\selectlanguage{latin}
Romæ sancti Cælestíni 
 Papæ Primi, qui damnávit Constantinopolitánum Epíscopum Nestórium, 
 Pelagiúmque fugávit; cujus étiam auctoritáte universális sancta Synodus 
 Ephesína advérsus eúndem Nestórium celebráta est.
\switchcolumn
\selectlanguage{english}
At Rome, Pope St. Celestine I, who 
 had condemned Nestorius, bishop of Constantinople, and put Pelagius to 
 flight. By his command the holy universal Council of Ephesus was also 
 held against the same Nestorius.
\switchcolumn*
\selectlanguage{latin}
Antisiodóri deposítio 
 beáti Æthérii, Epíscopi et Confessóris.
\switchcolumn
\selectlanguage{english}
At Auxerre, the death of blessed 
 Aetherius, bishop and confessor.
\switchcolumn*
\selectlanguage{latin}
Constantinópoli beátæ 
 Anthúsæ Vírginis, quæ, sub Constantíno Coprónymo, ob cultum sanctárum 
 Imáginum, verbéribus cæsa et exsílio relegáta, quiévit in Dómino.
\switchcolumn
\selectlanguage{english}
At Constantinople, blessed Anthusa, 
 virgin. After being scourged and banished by Constantine Copronymus 
 for the veneration of holy images, she rested in the Lord.
\switchcolumn*
\selectlanguage{latin}
\end{paracol}


% ---- martyrology/mart07/mart0728.htm
\needspace{10\baselineskip}
\begin{paracol}{2}
\selectlanguage{latin}
\begin{center}{\color{gregoriocolor} Quinto Kaléndas Augústi. 
 Luna\dots\ }\end{center}
\switchcolumn
\selectlanguage{english}
\begin{center}{\color{gregoriocolor} The 
 28\textsuperscript{th} Day of 
 July. The\dots\ Day of the Moon.}\end{center}
\end{paracol}

\noindent\begin{tabularx}{\linewidth}{*{19}{>{\centering\arraybackslash}X}}
 \textcolor{gregoriocolor}{a} & \textcolor{gregoriocolor}{b} & \textcolor{gregoriocolor}{c} & \textcolor{gregoriocolor}{d} & \textcolor{gregoriocolor}{e} & \textcolor{gregoriocolor}{f} & \textcolor{gregoriocolor}{g} & \textcolor{gregoriocolor}{h} & \textcolor{gregoriocolor}{i} & \textcolor{gregoriocolor}{k} & \textcolor{gregoriocolor}{l} & \textcolor{gregoriocolor}{m} & \textcolor{gregoriocolor}{n} & \textcolor{gregoriocolor}{p} & \textcolor{gregoriocolor}{q} & \textcolor{gregoriocolor}{r} & \textcolor{gregoriocolor}{s} & \textcolor{gregoriocolor}{t} & \textcolor{gregoriocolor}{u} \\
 3 & 4 & 5 & 6 & 7 & 8 & 9 & 10 & 11 & 12 & 13 & 14 & 15 & 16 & 17 & 18 & 19 & 20 & 21 \\
\end{tabularx}
\vspace{0.5\baselineskip}
\noindent\begin{tabularx}{\linewidth}{*{12}{>{\centering\arraybackslash}X}}
 \textcolor{gregoriocolor}{A} & \textcolor{gregoriocolor}{B} & \textcolor{gregoriocolor}{C} & \textcolor{gregoriocolor}{D} & \textcolor{gregoriocolor}{E} & F & \textcolor{gregoriocolor}{F} & \textcolor{gregoriocolor}{G} & \textcolor{gregoriocolor}{H} & \textcolor{gregoriocolor}{M} & \textcolor{gregoriocolor}{N} & \textcolor{gregoriocolor}{P} \\
 22 & 23 & 24 & 25 & 26 & 27 & 27 & 28 & 29 & 30 & 1 & 2 \\
\end{tabularx}

\begin{paracol}{2}
\selectlanguage{latin}
\lettrine[lines=2]{M}{edioláni} natális 
 sanctórum Mártyrum Nazárii, et Celsi púeri, quos Anolínus, sub rábie 
 persecutiónis quæ per Nerónem excitáta est, diu macerátos et afflíctos in cárcere, gládio feríri jussit.
\switchcolumn
\selectlanguage{english}
\lettrine[lines=2]{A}{t} Milan, the birthday of the holy 
 martyrs Nazarius and a boy named Celsus. While the persecution excited 
 by Nero was raging, they were beheaded by Anolinus, after long sufferings 
 and afflictions endured in prison.
\switchcolumn*
\selectlanguage{latin}
Romæ pássio sancti 
 Victóris Primi, Papæ et Mártyris.
\switchcolumn
\selectlanguage{english}
At Rome, the martyrdom of St. 
 Victor, pope and martyr.
\switchcolumn*
\selectlanguage{latin}
Item Romæ sancti 
 Innocéntii Primi, Papæ et Confessóris, qui ad Dóminum migrávit quarto Idus 
 Mártii.
\switchcolumn
\selectlanguage{english}
Also at Rome, St. Innocent, pope and 
 confessor, who passed to the Lord on the 12th of March.
\switchcolumn*
\selectlanguage{latin}
Thebáide, in Ægypto, 
 commemorátio plurimórum sanctórum Mártyrum, qui in Décii et Valeriáni 
 persecutióne passi sunt, quando Christiánis, optántibus pro Christi nómine 
 gládio pércuti, cállidus hostis, tarda ad montem supplícia conquírens, 
 ánimas cupiébat juguláre, non córpora. Ex eórum número unus, post 
 equúleos et láminas ac sartágines superátas, melle perúnctus, ligátis 
 mánibus post tergum, sub ardentíssimo sole, fucórum ac muscárum acúleis 
 expósitus fuit; álius, inter flores mólliter vinctus, cum ad eum, ut 
 concitáret in libídinem, impudicíssimum scortum venísset, præcísam morsu 
 linguam in blandiéntis fáciem éxspuit.
\switchcolumn
\selectlanguage{english}
In Thebais in Egypt, the 
 commemoration of many holy martyrs who suffered in the persecution of Decius 
 and Valerian. At this time, when Christians sought death by the sword 
 for the name of Christ, the crafty enemy devised certain slow torments to put 
 them to death, wishing to kill their souls much more than their bodies. 
 One of these Christians, after suffering the tortured of the rack, of hot 
 metal plates and of seething oil, was smeared with honey and exposed, in the 
 broiling heat of the sun, with his hands tied behind him, to the sting of 
 wasps and flies. Another, bound and placed among flowers, being 
 approached by a shameless woman for the purpose of exciting his passions, 
 bit through his tongue and spat it in her face.
\switchcolumn*
\selectlanguage{latin}
Ancyræ, in Galátia, sancti Eustáthii Mártyris, qui, váriis tormentórum genéribus excruciátibus, 
 atque demérsus in flumen, sed inde ab Angelo eréptus, demum, colúmba e cælis 
 adveniénte, ad præmia ætérna vocátus est.
\switchcolumn
\selectlanguage{english}
At Ancyra in Galatia, the holy 
 martyr Eustathius. After various torments he was plunged into a river, 
 but being delivered by an angel, was finally called to his eternal reward by 
 a dove coming from heaven.
\switchcolumn*
\selectlanguage{latin}
Miléti, in Cária, sancti Acátii Mártyris, qui, sub Licínio Imperatóre, post divérsas pœnas, 
 injéctus in fornácem ac Dei ope servátus illæsus, cápitis abscissióne 
 martyrium complévit.
\switchcolumn
\selectlanguage{english}
At Miletus, in the time of Emperor 
 Licinius, the holy martyr Acatius, who completed his martyrdom by having his 
 head struck off, after having undergone different torments and having been 
 thrown into a furnace, from which through the assistance of God he came 
 out uninjured.
\switchcolumn*
\selectlanguage{latin}
In Británnia minóre 
 sancti Sampsónis, Epíscopi et Confessóris.
\switchcolumn
\selectlanguage{english}
In Brittany, St. Sampson, bishop and 
 confessor.
\switchcolumn*
\selectlanguage{latin}
Lugdúni, in Gállia, 
 sancti Peregríni Presbyteri, cujus beatitúdinem glória miraculórum testátur.
\switchcolumn
\selectlanguage{english}
At Lyons in France, St. Peregrinus, 
 priest, whose happiness in heaven is testified by glorious miracles.
\switchcolumn*
\selectlanguage{latin}
\end{paracol}


% ---- martyrology/mart07/mart0729.htm
\needspace{10\baselineskip}
\begin{paracol}{2}
\selectlanguage{latin}
\begin{center}{\color{gregoriocolor} Quarto Kaléndas Augústi. 
 Luna\dots\ }\end{center}
\switchcolumn
\selectlanguage{english}
\begin{center}{\color{gregoriocolor} The 
 29\textsuperscript{th} Day of 
 July. The\dots\ Day of the Moon.}\end{center}
\end{paracol}

\noindent\begin{tabularx}{\linewidth}{*{19}{>{\centering\arraybackslash}X}}
 \textcolor{gregoriocolor}{a} & \textcolor{gregoriocolor}{b} & \textcolor{gregoriocolor}{c} & \textcolor{gregoriocolor}{d} & \textcolor{gregoriocolor}{e} & \textcolor{gregoriocolor}{f} & \textcolor{gregoriocolor}{g} & \textcolor{gregoriocolor}{h} & \textcolor{gregoriocolor}{i} & \textcolor{gregoriocolor}{k} & \textcolor{gregoriocolor}{l} & \textcolor{gregoriocolor}{m} & \textcolor{gregoriocolor}{n} & \textcolor{gregoriocolor}{p} & \textcolor{gregoriocolor}{q} & \textcolor{gregoriocolor}{r} & \textcolor{gregoriocolor}{s} & \textcolor{gregoriocolor}{t} & \textcolor{gregoriocolor}{u} \\
 4 & 5 & 6 & 7 & 8 & 9 & 10 & 11 & 12 & 13 & 14 & 15 & 16 & 17 & 18 & 19 & 20 & 21 & 22 \\
\end{tabularx}
\vspace{0.5\baselineskip}
\noindent\begin{tabularx}{\linewidth}{*{12}{>{\centering\arraybackslash}X}}
 \textcolor{gregoriocolor}{A} & \textcolor{gregoriocolor}{B} & \textcolor{gregoriocolor}{C} & \textcolor{gregoriocolor}{D} & \textcolor{gregoriocolor}{E} & F & \textcolor{gregoriocolor}{F} & \textcolor{gregoriocolor}{G} & \textcolor{gregoriocolor}{H} & \textcolor{gregoriocolor}{M} & \textcolor{gregoriocolor}{N} & \textcolor{gregoriocolor}{P} \\
 23 & 24 & 25 & 26 & 27 & 28 & 28 & 29 & 30 & 1 & 2 & 3 \\
\end{tabularx}

\begin{paracol}{2}
\selectlanguage{latin}
\lettrine[lines=2]{T}{arásci,} in Gállia 
 Narbonénsi, sanctæ Marthæ Vírginis, hóspitæ Salvatóris nostri ac soróris 
 beatórum Maríæ Magdalénæ et Lázari.
\switchcolumn
\selectlanguage{english}
\lettrine[lines=2]{A}{t} Tarascon, in the province of 
 Narbonne in France, St. Martha, virgin, the hostess of our Saviour and 
 sister of blessed Mary Magdalene and St. Lazarus.
\switchcolumn*
\selectlanguage{latin}
Romæ, via Aurélia, sancti Felícis Secúndi, Papæ et Mártyris; qui, ab Ariáno Imperatóre 
 Constántio, ob cathólicæ fídei defensiónem, e sede sua dejéctus, et 
 Ceræ, in Túscia, occúlte gládio necátus, glorióse occúbuit. Ejus 
 corpus, inde a Cléricis raptum, eádem via sepúltum fuit; póstea vero, ad 
 Ecclésiam sanctórum Cosmæ et Damiáni delátum, ibídem, Gregório Décimo tértio 
 Summo Pontífice, repértum est sub altári, una cum relíquiis sanctórum 
 Mártyrum Marci, Marcelliáni et Tranquillíni, atque in eódem loco, prídie 
 Kaléndas Augústi, simul cum eis recónditum. In quo étiam altári 
 invénta fuérunt córpora sanctórum Mártyrum Abúndii Presbyteri, et Abundántii 
 Diáconi; quæ, non multo post, ad Ecclésiam Societátis Jesu solémniter 
 transláta sunt prídie natális eórum.
\switchcolumn
\selectlanguage{english}
At Rome, on the Aurelian Way, St. 
 Felix II, pope and martyr. Being expelled from his See by the Arian 
 emperor Constantius for defending the Catholic faith, and being put to the 
 sword privately at Cera in Tuscany, he died gloriously. His body was 
 taken away from that place by clerics, and buried on the Aurelian Way. 
 It was afterwards brought to the Church of the Saints Cosmas and Damian, 
 where, under the Sovereign Pontiff Gregory XIII, it was found beneath the 
 altar with the relics of the holy martyrs Mark, Marcellian, and 
 Tranquillinus, and with the latter was put back in the same place on the 
 31st of July. In the same altar were also found the bodies of the holy 
 martyrs Abundius, a priest, and Abundantius, a deacon, which were shortly 
 after solemnly transferred to the church of the Society of Jesus, on the eve 
 of their feast.
\switchcolumn*
\selectlanguage{latin}
Item Romæ, via 
 Portuénsi, sanctórum Mártyrum Simplícii, Faustíni et Beatrícis, tempóribus 
 Diocletiáni Imperatóris. Horum duo primi, post multa et divérsa 
 supplícia, jussi sunt capitálem subíre senténtiam; Beátrix vero, eórum soror, 
 in confessióne Christi est in cárcere præfocáta.
\switchcolumn
\selectlanguage{english}
Also at Rome, on the Via Portuensis, 
 the holy martyrs Simplicius, Faustinus, and Beatrice, in the time of Emperor 
 Diocletian. The first two, after being subjected to many different 
 torments, were condemned to suffer death; Beatrice, their sister, was 
 smothered in prison for the confession of Christ.
\switchcolumn*
\selectlanguage{latin}
Romæ prætérea sanctórum 
 Mártyrum Lucíllæ et Floræ Vírginum, Eugénii, Antoníni, Theodóri, et Sociórum 
 decem et octo; qui sub Galliéno Imperatóre martyrium obiérunt.
\switchcolumn
\selectlanguage{english}
At Rome, likewise the holy martyrs 
 Lucilla and Flora, virgins, Eugenius, Antoninus, Theodore, and eighteen 
 companions, who underwent martyrdom in the reign of Emperor Gallienus.
\switchcolumn*
\selectlanguage{latin}
Item Romæ sanctæ 
 Serápiæ Vírginis, quæ, sub Hadriáno Príncipe, cum esset trádita duóbus 
 lascívis juvénibus et mínime potuísset illúdi, nec póstmodum ardéntibus 
 fácibus inflammári. Derílli Júdicis jussu fústibus cæsa est, dehinc 
 gládio decolláta. Ejus corpus a beáta Sabína in suo monuménto, juxta áream Vindiciáni, sepúltum est; sed memória ipsíus martyrii tértio Nonas 
 Septémbris celébrior habétur, quo die ambárum sarcóphagum ibi compósitum et 
 ornátum fuit, et locus oratiónis condígne dicátus.
\switchcolumn
\selectlanguage{english}
Again at Rome, St. Serapia, virgin. 
 Under Emperor Hadrian, she was delivered to two lustful young men, and as 
 she could not be corrupted, nor afterwards burned with lighted torches, she 
 was beaten with rods, and finally beheaded by order of the judge Derillus. 
 She was buried by blessed Sabina in her own tomb, near the field of 
 Vindician. But the commemoration of her martyrdom is celebrated more 
 solemnly on the 3rd of September, when their common tomb was finished and 
 adorned, and dedicated as a place of prayer.
\switchcolumn*
\selectlanguage{latin}
Gangris in Paphlagónia, 
 sancti Calliníci Mártyris, qui, virgis férreis verberátus aliísque 
 supplíciis afflíctus, tandem, in fornácem injéctus, spíritum Deo réddidit.
\switchcolumn
\selectlanguage{english}
At Gangra in Paphlagonia, St. 
 Callinicus, martyr, who was scourged with iron rods, and given over to other 
 torments. Being finally cast into a furnace, he gave up his soul to 
 God.
\switchcolumn*
\selectlanguage{latin}
In Norvégia sancti 
 Olávi, Regis et Mártyris.
\switchcolumn
\selectlanguage{english}
In Norway, St. Olaf, king and 
 martyr.
\switchcolumn*
\selectlanguage{latin}
Trecis, in Gállia, 
 sancti Lupi, Epíscopi et Confessóris; qui, cum beáto Germáno, ad expugnándam 
 Pelagianórum hæresim, in Británnia perréxit, urbémque Trecas a furóre Attilæ, 
 Gálliam omnem devastántis, assídua oratióne deféndit; demum, quinquagínta 
 duóbus annis sacerdótio venerabíliter functus, in pace quiévit.
\switchcolumn
\selectlanguage{english}
At Troyes in France, St. Lupus, 
 bishop and confessor, who went with blessed Germanus to England to 
 exterminate the Pelagian heresy, and by diligent prayer defended the city of 
 Troyes from the wrath of Attila, who was devastating all of France. At 
 length, having religiously discharged the functions of the priesthood for 
 fifty-two years, he rested in peace.
\switchcolumn*
\selectlanguage{latin}
In civitáte Briocénsi, 
 in Gállia, sancti Guliélmi, Epíscopi et Confessóris.
\switchcolumn
\selectlanguage{english}
At St. Brieuc in France, St. 
 William, bishop and confessor.
\switchcolumn*
\selectlanguage{latin}
Item deposítio beáti 
 Prósperi, Aurelianénsis Epíscopi.
\switchcolumn
\selectlanguage{english}
Also, the death of blessed Prosper, 
 bishop of Orleans.
\switchcolumn*
\selectlanguage{latin}
Apud Tudértum, in 
 Umbria, sancti Faustíni Confessóris.
\switchcolumn
\selectlanguage{english}
At Todi in Umbria, St. Faustinus, 
 confessor.
\switchcolumn*
\selectlanguage{latin}
In civitáte Mamiénsi 
 sanctæ Seraphínæ.
\switchcolumn
\selectlanguage{english}
At Mamia, St. Serafina.
\switchcolumn*
\selectlanguage{latin}
Romæ beáti Urbáni Papæ 
 Secúndi, qui, sancti Gregórii Séptimi vestígia secútus, doctrínæ et 
 religiónis stúdio enítuit, et fidéles Cruce signátos ad sacra Palæstínæ loca 
 ex infidélium domínio rediménda excitávit. Cultum autem, ab 
 immemorábili témpore eídem exhíbitum, Leo Décimus tértius, Póntifex Máximus, 
 ratum hábuit et confirmávit.
\switchcolumn
\selectlanguage{english}
At Rome, blessed Pope Urban II who 
 followed in the path of St. Gregory VII. He was resplendent for his 
 zeal for learning and religion, and aroused the faithful, signed with the 
 sign of the cross, to recover the holy places of Palestine from the power of 
 the infidels. Pope Leo XIII ratified and confirmed the veneration 
 shewn him from time immemorial.
\switchcolumn*
\selectlanguage{latin}
\end{paracol}


% ---- martyrology/mart07/mart0730.htm
\needspace{10\baselineskip}
\begin{paracol}{2}
\selectlanguage{latin}
\begin{center}{\color{gregoriocolor} Tértio Kaléndas Augústi. 
 Luna\dots\ }\end{center}
\switchcolumn
\selectlanguage{english}
\begin{center}{\color{gregoriocolor} The 
 30\textsuperscript{th} Day of 
 July. The\dots\ Day of the Moon.}\end{center}
\end{paracol}

\noindent\begin{tabularx}{\linewidth}{*{19}{>{\centering\arraybackslash}X}}
 \textcolor{gregoriocolor}{a} & \textcolor{gregoriocolor}{b} & \textcolor{gregoriocolor}{c} & \textcolor{gregoriocolor}{d} & \textcolor{gregoriocolor}{e} & \textcolor{gregoriocolor}{f} & \textcolor{gregoriocolor}{g} & \textcolor{gregoriocolor}{h} & \textcolor{gregoriocolor}{i} & \textcolor{gregoriocolor}{k} & \textcolor{gregoriocolor}{l} & \textcolor{gregoriocolor}{m} & \textcolor{gregoriocolor}{n} & \textcolor{gregoriocolor}{p} & \textcolor{gregoriocolor}{q} & \textcolor{gregoriocolor}{r} & \textcolor{gregoriocolor}{s} & \textcolor{gregoriocolor}{t} & \textcolor{gregoriocolor}{u} \\
 5 & 6 & 7 & 8 & 9 & 10 & 11 & 12 & 13 & 14 & 15 & 16 & 17 & 18 & 19 & 20 & 21 & 22 & 23 \\
\end{tabularx}
\vspace{0.5\baselineskip}
\noindent\begin{tabularx}{\linewidth}{*{12}{>{\centering\arraybackslash}X}}
 \textcolor{gregoriocolor}{A} & \textcolor{gregoriocolor}{B} & \textcolor{gregoriocolor}{C} & \textcolor{gregoriocolor}{D} & \textcolor{gregoriocolor}{E} & F & \textcolor{gregoriocolor}{F} & \textcolor{gregoriocolor}{G} & \textcolor{gregoriocolor}{H} & \textcolor{gregoriocolor}{M} & \textcolor{gregoriocolor}{N} & \textcolor{gregoriocolor}{P} \\
 24 & 25 & 26 & 27 & 28 & 29 & 29 & 30 & 1 & 2 & 3 & 4 \\
\end{tabularx}

\begin{paracol}{2}
\selectlanguage{latin}
\lettrine[lines=2]{R}{omæ} sanctórum Mártyrum 
 Abdon et Sennen Persárum, qui, sub Décio, caténis alligáti et Romam addúcti, 
 pro Christi fide prius plumbátis cæsi sunt, deínde gládio interfécti.
\switchcolumn
\selectlanguage{english}
\lettrine[lines=2]{A}{t} Rome, in the reign of Decius, the 
 holy Persian martyrs Abdon and Sennen, who were bound with chains, brought 
 to Rome, scourged with leaded whips for the faith of Christ, and then put to 
 the sword.
\switchcolumn*
\selectlanguage{latin}
Assísii, in Umbria, 
 sancti Rufíni Mártyris.
\switchcolumn
\selectlanguage{english}
At Assisi in Umbria, St. Rufinus, 
 martyr.
\switchcolumn*
\selectlanguage{latin}
Tubúrbi Lucernáriæ, in 
 Africa, pássio sanctárum Vírginum Máximæ, Donatíllæ et Secúndæ; quarum duæ 
 primæ, in persecutióne Valeriáni et Galliéni, acéto et felle potátæ, deínde 
 plagis acérrimis cæsæ et equúlei extensióne cruciátæ, cratículus étiam 
 exústæ et calce perfricátæ, póstmodum cum Secúnda, Vírgine duódecim annórum, 
 ad béstias simul projéctæ, sed ab his intáctæ, novíssime gládio jugulátæ 
 sunt.
\switchcolumn
\selectlanguage{english}
At Tuberbum Lucernarium in Africa, 
 the holy virgins and martyrs Maxima, Donatilla, and Secunda. The first 
 two, in the persecution of Valerian and Gallienus, were forced to drink 
 vinegar and gall, then scourged most severely, stretched on the rack, 
 burned on the gridiron, rubbed over with lime, and afterwards exposed to the 
 beasts with the virgin Secunda, twelve years old. But being untouched 
 by them, they were finally beheaded.
\switchcolumn*
\selectlanguage{latin}
Cæsaréæ, in Cappadócia, sanctæ Julíttæ Mártyris, quæ, cum bona sua, a quodam poténte sibi usurpáta, 
 in judícia repéteret, atque ille exceptiónem daret quod ut Christiáno non 
 debéret audíri, mox a Júdice jussa est thus idólis offérre, ut posset audíri. 
 Quod illa constánter recúsans, in ignem conjécta est, sicque spíritum Deo 
 réddidit; corpus autem a flamma remánsit illæsum. Ejus præclárus 
 laudes sanctus Basilíus Magnus egrégio encómio celebrávit.
\switchcolumn
\selectlanguage{english}
At Caesarea in Cappadocia, St. 
 Julitta, martyr. As she sought through the courts the restitution of 
 goods seized by a man of influence, the latter objected that, being a 
 Christian, her cause could not be pleaded. The judge commanded her to 
 offer sacrifice to the idols, that she might be heard. She refused 
 with great constancy, and being thrown into the fire, yielded her soul unto 
 God. Her body remained uninjured by the flames. St. Basil the 
 Great has proclaimed her praise in an excellent eulogy.
\switchcolumn*
\selectlanguage{latin}
Antisiodóri sancti Ursi, 
 Epíscopi et Confessóris.
\switchcolumn
\selectlanguage{english}
At Auxerre, St. Ursus, bishop and 
 confessor.
\switchcolumn*
\selectlanguage{latin}
\end{paracol}


% ---- martyrology/mart07/mart0731.htm
\needspace{10\baselineskip}
\begin{paracol}{2}
\selectlanguage{latin}
\begin{center}{\color{gregoriocolor} Prídie Kaléndas Augústi. 
 Luna\dots\ }\end{center}
\switchcolumn
\selectlanguage{english}
\begin{center}{\color{gregoriocolor} The 
 31\textsuperscript{st} Day of 
 July. The\dots\ Day of the Moon.}\end{center}
\end{paracol}

\noindent\begin{tabularx}{\linewidth}{*{19}{>{\centering\arraybackslash}X}}
 \textcolor{gregoriocolor}{a} & \textcolor{gregoriocolor}{b} & \textcolor{gregoriocolor}{c} & \textcolor{gregoriocolor}{d} & \textcolor{gregoriocolor}{e} & \textcolor{gregoriocolor}{f} & \textcolor{gregoriocolor}{g} & \textcolor{gregoriocolor}{h} & \textcolor{gregoriocolor}{i} & \textcolor{gregoriocolor}{k} & \textcolor{gregoriocolor}{l} & \textcolor{gregoriocolor}{m} & \textcolor{gregoriocolor}{n} & \textcolor{gregoriocolor}{p} & \textcolor{gregoriocolor}{q} & \textcolor{gregoriocolor}{r} & \textcolor{gregoriocolor}{s} & \textcolor{gregoriocolor}{t} & \textcolor{gregoriocolor}{u} \\
 6 & 7 & 8 & 9 & 10 & 11 & 12 & 13 & 14 & 15 & 16 & 17 & 18 & 19 & 20 & 21 & 22 & 23 & 24 \\
\end{tabularx}
\vspace{0.5\baselineskip}
\noindent\begin{tabularx}{\linewidth}{*{12}{>{\centering\arraybackslash}X}}
 \textcolor{gregoriocolor}{A} & \textcolor{gregoriocolor}{B} & \textcolor{gregoriocolor}{C} & \textcolor{gregoriocolor}{D} & \textcolor{gregoriocolor}{E} & F & \textcolor{gregoriocolor}{F} & \textcolor{gregoriocolor}{G} & \textcolor{gregoriocolor}{H} & \textcolor{gregoriocolor}{M} & \textcolor{gregoriocolor}{N} & \textcolor{gregoriocolor}{P} \\
 25 & 26 & 27 & 28 & 29 & 1 & 30 & 1 & 2 & 3 & 4 & 5 \\
\end{tabularx}

\begin{paracol}{2}
\selectlanguage{latin}
\lettrine[lines=2]{R}{omæ} natális sancti 
 Ignátii, Presbyteri et Confessóris, qui Fundátor éxstitit Societátis Jesu, 
 atque vir fuit sanctitáte et miráculis clarus, ac religiónis cathólicæ 
 ubíque dilatándæ studiosíssimus; quem Pius Undécimus, Póntifex Máximus, 
 cæléstem ómnium Exercitiórum spirituálium Patrónum constítuit.
\switchcolumn
\selectlanguage{english}
\lettrine[lines=2]{A}{t} Rome, the birthday of St. 
 Ignatius, priest and confessor, founder of the Society of Jesus, renowned 
 for sanctity and miracles, and most zealous for propagating the Catholic 
 religion in all parts of the world. Pope Pius XI declared him to be 
 the heavenly patron of all spiritual retreats.
\switchcolumn*
\selectlanguage{latin}
Medioláni sancti 
 Calimérii, Epíscopi et Mártyris; qui, in Antoníni persecutióne, comprehénsus, 
 vulnéribus confóssus, cervicibúsque gládio transverberátis, præceps in 
 púteum dejéctus, martyrii cursum confécit.
\switchcolumn
\selectlanguage{english}
At Milan, during the persecution of 
 Antoninus, St. Calimerius, bishop and martyr, who was arrested, covered with 
 wounds, and his throat transfixed with a sword. He completed his 
 martyrdom by being cast into a well.
\switchcolumn*
\selectlanguage{latin}
Cæsaréæ, in Mauritánia, pássio beáti Fábii Mártyris, qui, cum ferre vexílla præsidália recusáret, 
 primo in cárcerem trusus est, ibíque diébus áliquot deténtus; deínde, cum in 
 confessióne Christi, interrogátus semel et íterum, immóbilis perduráret, 
 senténtia capitáli a Júdice condemnátur.
\switchcolumn
\selectlanguage{english}
At Caesarea in Mauretania, the 
 martyrdom of the blessed martyr Fabius. Because he refused to carry 
 the banners of the governor of the province, he was thrown into prison for 
 some days, and as he persisted twice in confessing Christ when brought 
 before the judge, he was condemned to death.
\switchcolumn*
\selectlanguage{latin}
Synnadæ, in Phrygia 
 Pacatiána, sanctórum Mártyrum Demócriti, Secúndi et Dionysii.
\switchcolumn
\selectlanguage{english}
At Synnada in Phrygia Pacatiana, the 
 holy martyrs Democritus, Secundus, and Denis.
\switchcolumn*
\selectlanguage{latin}
In Syria sanctórum 
 trecentórum quinquagínta Monachórum Mártyrum, qui, ob defensiónem Synodi 
 Chalcedonénsis, ab hæréticis sunt occísi.
\switchcolumn
\selectlanguage{english}
In Syria, three hundred and fifty 
 monks, who became martyrs by being slain by the heretics for defending the 
 Council of Chalcedon.
\switchcolumn*
\selectlanguage{latin}
Ravénnæ tránsitus 
 sancti Germáni, Antisiodorénsis Epíscopi, génere, fide, doctrína et 
 miraculórum glória claríssimi; qui Británniam a Pelagianórum hærésibus 
 pénitus liberávit.
\switchcolumn
\selectlanguage{english}
At Ravenna, the death of St. 
 German, bishop of Auxerre, a man most renowned for his birth, faith, 
 learning, and glorious miracles, who freed England completely from the 
 heretical doctrines of the Pelagians.
\switchcolumn*
\selectlanguage{latin}
Tagáste, in Africa, 
 sancti Firmi Epíscopi, confessiónis glória conspícui.
\switchcolumn
\selectlanguage{english}
At Tagaste in Africa, St. Firmus, 
 bishop, illustrious by a glorious confession of the faith.
\switchcolumn*
\selectlanguage{latin}
Senis, in Túscia, 
 natális beáti Joánnis Columbíni, qui fuit Institútor Ordinis Jesuatórum, et 
 sanctitáte ac miráculis cláruit.
\switchcolumn
\selectlanguage{english}
At Siena in Tuscany, the birthday of 
 blessed John Colombini, founder of the Order of Gesuati, renowned for 
 sanctity and miracles.
\switchcolumn*
\selectlanguage{latin}
\end{paracol}

